\documentclass[mat2, tisk]{fmfdelo}
\usepackage{graphicx}
\usepackage{enumitem}
\usepackage{amsmath}
\usepackage{dsfont}
% chktex-file 3
% chktex-file 9
% \documentclass[fin2, tisk]{fmfdelo}
% \documentclass[isrm2, tisk]{fmfdelo}
% \documentclass[ped, tisk]{fmfdelo}
% Če pobrišete možnost tisk, bodo povezave obarvane,
% na začetku pa ne bo praznih strani po naslovu, …

%%%%%%%%%%%%%%%%%%%%%%%%%%%%%%%%%%%%%%%%%%%%%%%%%%%%%%%%%%%%%%%%%%%%%%%%%%%%%%%
% METAPODATKI
%%%%%%%%%%%%%%%%%%%%%%%%%%%%%%%%%%%%%%%%%%%%%%%%%%%%%%%%%%%%%%%%%%%%%%%%%%%%%%%

% - vaše ime
\avtor{Tadeja Možina}

% - naslov dela v slovenščini
\naslov{Lastni problem nelinearnih preslikav na stožcih}

% - naslov dela v angleščini
\title{The eigenproblem of nonlinear mappings on cones}

% - ime mentorja/mentorice s polnim nazivom:
%   - doc.~dr.~Ime Priimek
%   - izr.~prof.~dr.~Ime Priimek
%   - prof.~dr.~Ime Priimek
%   za druge variante uporabite ustrezne ukaze
\mentor{~izr.~prof.~dr.~Aljoša Peperko}
% \somentor{...}
% \mentorica{...}
% \somentorica{...}
% \mentorja{...}{...}
% \somentorja{...}{...}
% \mentorici{...}{...}
% \somentorici{...}{...}

% - leto magisterija
\letnica{2026}

% - povzetek v slovenščini
%   V povzetku na kratko opišite vsebinske rezultate dela. Sem ne sodi razlaga
%   organizacije dela, torej v katerem razdelku je kaj, pač pa le opis vsebine.
\povzetek{Tukaj napišemo povzetek vsebine. Sem sodi razlaga vsebine in ne opis tega, kako je delo organizirano.}

% - povzetek v angleščini
\abstract{An abstract of the work is written here. This includes a short description of
the content and not the structure of your work.}

% - klasifikacijske oznake, ločene z vejicami
%   Oznake, ki opisujejo področje dela, so dostopne na strani https://www.ams.org/msc/
\klasifikacija{15B48, 47L07}

% - ključne besede, ki nastopajo v delu, ločene s \sep
\kljucnebesede{stožci, lastni problem}

% - angleški prevod ključnih besed
\keywords{cones, eigenproblem}

% - neobvezna zahvala
\zahvala{
  Neobvezno.
  Zahvaljujem se \dots
}

% - ime datoteke z viri (vključno s končnico .bib), če uporabljate BibTeX
\literatura{literatura.bib}

%%%%%%%%%%%%%%%%%%%%%%%%%%%%%%%%%%%%%%%%%%%%%%%%%%%%%%%%%%%%%%%%%%%%%%%%%%%%%%%
% DODATNE DEFINICIJE
%%%%%%%%%%%%%%%%%%%%%%%%%%%%%%%%%%%%%%%%%%%%%%%%%%%%%%%%%%%%%%%%%%%%%%%%%%%%%%%

% naložite dodatne pakete, ki jih potrebujete
\usepackage{units}        % fizikalne enote kot \unit[12]{kg} s polovico nedeljivega presledka, glej primer v kodi
\usepackage{graphicx}     % za slike
% \usepackage{tikz}
% VEČ ZANIMIVIH PAKETOV
% \usepackage{array}      % več možnosti za tabele
% \usepackage[list=true,listformat=simple]{subcaption}  % več kot ena slika na figure, omogoči slika 1a, slika 1b
% \usepackage[all]{xy}    % diagrami
% \usepackage{doi}        % za clickable DOI entrye v bibliografiji
% \usepackage{enumerate}     % več možnosti za sezname

% Za barvanje source kode
% \usepackage{minted}
% \renewcommand\listingscaption{Program}

% Za pisanje psevdokode
% \usepackage{algpseudocode}  % za psevdokodo
% \usepackage{algorithm}
% \floatname{algorithm}{Algoritem}
% \renewcommand{\listalgorithmname}{Kazalo algoritmov}

% deklarirajte vse matematične operatorje, da jih bo LaTeX pravilno stavil
% \DeclareMathOperator{\...}{...}

% vstavite svoje definicije ...
\newcommand{\R}{\mathbb{R}}
\newcommand{\N}{\mathbb{N}}
\newcommand{\Z}{\mathbb{Z}}
% Lahko se zgodi, da je ukaz \C definiral že paket hyperref,
% zato dobite napako: Command \C already defined.
% V tem primeru namesto ukaza \newcommand uporabite \renewcommand
\newcommand{\C}{\mathbb{C}}
\newcommand{\Q}{\mathbb{Q}}

%%%%%%%%%%%%%%%%%%%%%%%%%%%%%%%%%%%%%%%%%%%%%%%%%%%%%%%%%%%%%%%%%%%%%%%%%%%%%%%
% ZAČETEK VSEBINE
%%%%%%%%%%%%%%%%%%%%%%%%%%%%%%%%%%%%%%%%%%%%%%%%%%%%%%%%%%%%%%%%%%%%%%%%%%%%%%%

\begin{document}

\section{Uvod}
Napišite kratek zgodovinski in matematični uvod.  Pojasnite motivacijo za problem, kje
nastopa, kje vse je bil obravnavan. Na koncu opišite tudi organizacijo dela -kaj je v
katerem razdelku.

Skozi celotno delo bomo gledali samo končno-dimenzionalne realne vektorske prostore. 

\section{Nelinearna Perron-Frobeniusova teorija}

Neki napisati.

\subsection{Klasična Perron-Frobeniusova teorija}

Na začetku se spomnimo nekaj definicij, ki jih že poznamo, in navedimo nekaj pomembnih 
izrekov, lem in posledic, ki jih bomo kasneje potrebovali.

\begin{definicija}
  Naj bo $A = (a_{ij})$ matrika velikosti $n \times n$.
  \begin{itemize}
    \item $A$ je \emph{nenegativna}, če velja $a_{ij} \geq 0$ za vsak $i$ in $j$,
    \item $A$ je \emph{pozitivna}, če velja $a_{ij} > 0$ za vsak $i$ in $j$.
  \end{itemize}
\end{definicija}

Podobno rečemo, da je vektor $x \in \R^n$ \emph{nenegativen} (oziroma \emph{pozitiven}), 
če so vse njegove koordinate nenegativne (oziroma pozitivne). 

\begin{definicija}
  Naj bo $A = (a_{ij})$ matrika velikosti $n \times n$. 
  
  \emph{Spekter} matrike $A$ je
  
  \begin{center}
    $\sigma(A) = \{ \lambda \in \C ~|~ Ax = \lambda x \; \text{za neki} \; x \in \C \setminus \{0\} \}$.
  \end{center}
  
  \emph{Spektralni radij} matrike $A$ je 

  \begin{center}
    $r(A) = \max\{ |\lambda | ~|~ \lambda \in \sigma(A) \}$.
  \end{center}
\end{definicija}

\begin{trditev}
  Za spektralni radij matrike $A \in \C ^ {n \times n}$ velja
  $r(A) = \lim\limits_{k \rightarrow \infty} || A^k ||^{\frac{1}{k}} $.
\end{trditev}

\begin{trditev}[Perron]\label{trd:perron}
  Naj bo $A$ pozitivna matrika. Potem veljajo naslednje trditve:
  \begin{itemize}
    \item $r(A) > 0$,
    \item $r(A)$ je enostavna ničla karakterističnega polinoma matrike $A$. Njen  
          pripadajoč normiran lastni vektor $v$ je pozitiven,
    \item Vsak nenegativen lastni vektor od $A$ je večkratnik vektorja $v$,
    \item Za vsako lastno vrednost $\lambda \in \sigma(A)$, za katero 
          drži $\lambda \neq r(A)$, velja $|\lambda| < r(A)$.
  \end{itemize}
\end{trditev}
\begin{proof}
  TODO
\end{proof}

Vemo že, da so vse pozitivne matrike nerazcepne. 

\begin{definicija}
  Naj bo $A = (a_{ij})$ nenegativna matrika velikosti $n \times n$. Matrika 
  $A$ je \emph{razcepna} natanko tedaj, ko obstaja permutacijska matrika $P$, da 
  velja 
  \begin{center}
    $P^T A P = 
    \begin{bmatrix}
      B & C \\
      0 & D
    \end{bmatrix}$,
  \end{center}

  kjer sta $B$ in $D$ kvadratni matriki, $0$ pa ničelna matrika.

  Matrika A je \emph{nerazcepna}, če taka permutacijska matrika $P$ ne 
  obstaja.
\end{definicija}

Hirto lahko vidimo, da je vsaka nenegativna $1 \times 1$ matrika nerazcepna.
Poglejmo sedaj, kako bi našo trditev~\ref{trd:perron} razširili tako, da bi veljala 
ne samo za pozitivne matrike, temveč tudi za nerazcepne matrike. To posplošitev 
Perronovega izreka je dokazal Frobenius.

\begin{trditev}[Perron-Frobenius]\label{trd:per-fro}
  Naj bo $A$ nenegativna nerazcepna $n \times n$ matrika. Potem veljajo naslednje trditve:
  \begin{itemize}
    \item Če je $A \neq [0]$, potem je $r(A) > 0$,
    \item $r(A)$ je enostavna ničla karakterističnega polinoma matrike $A$. Njen  
          pripadajoč normiran lastni vektor $v$ je pozitiven,
    \item Vsak nenegativen lastni vektor od $A$ je večkratnik vektorja $v$,
    \item Če ima $A$ natanko $q$ lastnih vrednosti $\lambda$, za katere
          velja $| \lambda | = r(A)$, potem so te lastne vrednosti 
          podane z $r(A) e^{2\pi i k / q}$.
  \end{itemize}
\end{trditev}

Številu $q$ v Perron-Frobeniusovi trditvi rečemo \emph{indeks cikličnosti matrike A}.

\subsubsection{Osnovne definicije o stožcih}

Poglejmo si sedaj, kaj sploh so stožci. Naj bo $V$ končno-dimenzionalen 
vektorski prostor.

\begin{definicija}
  Podmnožica $K$ končno-dimenzionalnega vektorskega prostora $V$ je \emph{stožec},
  če velja:
  \begin{enumerate}
    \item $K$ je konveksna množica,
    \item $\mu K \subseteq K$ za vsak $\mu \geq 0$,
    \item $K \cap (-K) = \{ 0 \}$.
  \end{enumerate}
\end{definicija}

\begin{definicija}
  Stožec je \emph{zaprt stožec}, če je zaprta množica v vektorskem prostoru $V$
  glede na standardno topologijo.
\end{definicija}

Ponovimo še nekaj oznak. Za podmnožico $S \subseteq V$ označimo \emph{notranjost} kot 
$\text{int}(S)$, \emph{zaprtje} kot $\text{cl}(S)$ in \emph{rob} kot $\partial S$.

\begin{definicija}
  Pravimo, da je stožec \emph{strnjen}, če ima neprazno notranjost.
\end{definicija}

Poglejmo si nekaj primerov polnih zaprtih stožcev.

\begin{zgled}
  Najbolj osnovna primera sta \emph{standardni pozitivni stožec}, ki ga označimo kot
  $\R_+ ^n = \{ ( x_1, \ldots, x_n ) \in \R^n ~|~ x_i \geq 0 ~\text{za}~ 1 \leq i \leq n \}$ 
  in \emph{Lorentzov stožec}, ki ga označimo kot 
  $\Lambda_{n+1} = \{ ( x_1, \ldots, x_{n+1} ) \in \R^{n+1} ~|~ 
  x_1^2 - x_2^2 - \ldots - x_{n+1}^2 \geq 0 ~\text{in}~ x_1 \geq 0\}$.% chktex 11
\end{zgled}

\begin{zgled}
  Poglejmo si še primer stožcev v prostoru realnih $n \times n$ simetričnih matrik. 
  Potem so pozitivne semidefinitne matrike, ki jih označimo z $\Pi_n (\R)$, strnjen
  zaprt stožec. Res, če je $M \in \Pi_n (\R)$, potem očitno za vsak $\mu \geq 0$ velja 
  $\mu M \in \Pi_n (\R)$ in $\Pi_n (\R) \cap (-\Pi_n (\R)) = \{ 0 \}$.
\end{zgled}

\begin{definicija}
  Dualni prostor prostora $V$ označimo z $V^*$, \emph{dualni stožec} $K^*$ pa je definiran kot
\[ K^* = \{ \varphi \in V^* : \varphi(x) \ge 0 \text{ za vse } x \in K \}.\]
\end{definicija}

\begin{definicija}
  \emph{Obraz zaprtega stožca} $K$ v $V$ je neprazen konveksna podmnožica $F \subset K$, 
  za katero velja: če sta $x, y \in K$ in je $\lambda x + (1-\lambda)y \in F$ 
  za nek $0 < \lambda < 1$, potem sta $x, y \in F$. 
\end{definicija}

Stožec $K$ sam in množica $\{0\}$ sta obraza in ju imenujemo \emph{neprava obraza}. 
Vsi ostali obrazi so pravi. Iz definicije obraza zaprtega stožca hitro opazimo, da 
vsi pravi obrazi stožca ne vsebujejo notranjih točk stožca, torej za $F$ obraz od $K$ 
velja $F \cap \operatorname{int}(K) \neq \varnothing$.

\begin{definicija}
  \emph{Relativna notranjost} konveksne množice $C \subseteq V$, označena z $\mathrm{relint}(C)$, 
  je notranjost $C$, gledane kot podmnožice afinega ovoja množice $C$ v $V$. 
\end{definicija}

Pripomnimo (glej~\cite[izrek 6.2]{rockafellar}), da sta $\mathrm{cl}(C)$ in $\mathrm{relint}(C)$ 
neprazni in konveksni, kadar je $C$ neprazen konveksen del prostora $V$.
To dejstvo lahko uporabimo, da dokažemo naslednjo trditev.
\begin{trditev}\label{trditev2.15}
  Vsak obraz zaprtega stožca je zaprt.
\end{trditev}
\begin{dokaz}
  Naj bo $F$ obraz zaprtega stožca $K$. Želimo $\mathrm{cl}(F) = F$. Vzemimo točko $z$  
  v relativni notranjosti $F$. Taka točka obstaja, ker je po definiciji obraza $F$ konveksen 
  in neprazen. Naj bo $x \in \mathrm{cl}(F)$, $x \neq z$. 
  Ker je $F$ obraz, po~\cite[izrek 6.2]{rockafellar} vemo, da sta $\mathrm{cl}(F)$ in 
  $\mathrm{relint}(F)$ neprazni in konveksni. Torej obstaja 
  $y \in \mathrm{cl}(F)$ tako, da je $z$ v relativni notranjosti daljice med $x$ in $y$, 
  zato je $x \in F$. (I don't understand this?)
\end{dokaz}


\begin{definicija}
  Obraz $F$ stožca $K$ imenujemo \emph{izpostavljeni obraz}, če obstaja $\varphi \in K^*$, 
za katerega velja
\[F = K \cap \{ x \in V : \varphi(x) = 0 \}.\]
\end{definicija}

\begin{definicija}
  Stožec $K$ v $V$ je \emph{poliedričen}, če je presečišče končno mnogo zaprtih polprostorov, tj. 
obstajajo $\varphi_1, \ldots, \varphi_m \in V^*$, tako da
\[K = \{ x \in V : \varphi_i(x) \ge 0 \text{ za vse } 1 \le i \le m \}.\]
\end{definicija}

Osnovni primer poliedrskega stožca je standardni pozitivni stožec $\R_+^n$, saj zanj 
lahko vzamemo funkcionale $\varphi_i(x) = x_i$ za $1 \le i \le n$.
Vsak obraz stožca ni nujno izpostavljen, vendar pa je vsak obraz poliedrskega stožca izpostavljen.

\begin{definicija}
  Obraz \( F \) poliedrskega stožca \( K \) imenujemo \emph{ličnica??} (facet), 
če velja \( \dim(F) = \dim(K) - 1 \). Tukaj \( \dim(F) \) označuje dimenzijo 
linearnega ovoja ploskve \( F \).
\end{definicija} 
Naslednji osnovni rezultat iz poliedrske 
geometrije~\cite[Poglavje~8.4]{201} bo uporaben v nadaljevanju.

\begin{lema}\label{lema2.19}
Če je \( K \subseteq V \) poliedrski stožec z \( N \) ličnicami, potem obstaja 
\( N \) linearnih funkcionalov \( \psi_1, \ldots, \psi_N \) takih, da

\[ K = \{ x \in V : \psi_i(x) \ge 0 \text{ za vse } 1 \le i \le N \}
    \cap \operatorname{span}(K)\]

in vsak linearni funkcional \( \psi_i \) ustreza natanko eni ličnici stožca \( K \).
\end{lema}

Funkcionale \( \psi_1, \ldots, \psi_N \) iz Leme~\ref{lema2.19}
imenujemo \emph{funkcionali, ki določajo ličnice}
poliedrskega stožca. 

\begin{zgled}\label{zgled2.20}
  Poglejmo si ličnice od \( \mathbb{R}^n_+ \). Obrazi tega stožca so očitno 
  \[
  F_I = \{ x \in \mathbb{R}^n_+ ~|~ x_i > 0 \text{ natanko tedaj, ko } i \in I \}
  \quad \text{za } I \subseteq \{1, \ldots, n\},
  \]
  Ker gledamo po podmnožicah množice z n elementi, je vseh obrazev natanko \( 2^n \).
  Ličnice stožca so vsi tisti obrazi, katerih dimenzija je 
  \( \dim(\mathbb{R}^n_+) - 1  = n - 1\), 
  torej vse podmnožice indeksne množice $I$ moči $n - 1$, teh pa je natanko $n$. 
  Ličnice so oblike
  \[
  F_i = \{ x \in \mathbb{R}^n_+ ~|~ x_i = 0 \text{ in } x_j > 0  \text{ za } j \neq i\}.
  \]
\end{zgled}



\subsubsection{Delna urejenost in invariantnost}

Stožec $K$ v  vektorskem prostoru $V$ inducira delno urejenost \(\leq_K\) na \(V\) z

\[
x \leq_K y \quad \text{natanko tedaj, ko je } y - x \in K.
\]

Če ima \(K\) neprazno notranjost, potem pišemo \(x \ll_K y\), kadar je
\(y - x \in \operatorname{int}(K)\). Prav tako pišemo \(x <_K y\), kadar
\(x \leq_K y\) in \(x \neq y\). Kadar je stožec \(K\) jasen iz konteksta,
pišemo samo \(x \leq y\), \(x < y\) in \(x \ll y\).

\begin{zgled}
  Poglejmo si primer \(K = \mathbb{R}^n_+\). Velja \(x \leq_K y\) natanko tedaj, 
  ko je $y - x \in \mathbb{R}^n_+$, kar pa pomeni, da je \(y_i - x_i \geq 0\) 
  za vse \(1 \leq i \leq n\) ali drugače povedano, da je \(x_i \leq y_i\) 
  za vse \(1 \leq i \leq n\). 
\end{zgled}

\begin{trditev}
  Linearni operator \(A : V \to V\) pusti stožec \(K\)
  invarianten natanko tedaj, ko \(0 \leq_K Ax\) za vse \(0 \leq_K x\) v \(V\).
\end{trditev}
\begin{dokaz}
  Recimo, da $Ax \in K$ za vsak $x \in K$. Če je \(0 \leq_K x\), sledi, 
  da $ x \in K$. Ker linearni operator $A$ ohranja stožec, vidimo, da 
  $Ax \in K$, torej velja \(0 \leq_K Ax\).

  V drugo smer: naj velja $0 \leq_K x \Rightarrow 0 \leq_K Ax$. Vzemimo 
  $x$ v stožcu, $x \in K$. Potem velja \(0 \leq_K x\), torej \(0 \leq_K Ax\), 
  kar pa pomeni $A(K) \subseteq K$.
\end{dokaz}

\begin{definicija}
Linearni operator \(A : V \to V\), ki pusti strnjen zaprt stožec \(K\) 
invarianten, imenujemo \emph{ireducibilen}, če noben pravi obraz stožca \(K\)
ni invarianten za \(A\).
\end{definicija}

\begin{zgled}
  Poglejmo si zgled ireducibilnosti na primeru stožca \(K = \mathbb{R}^n_+\) 
  in dokažimo, da je nenegativna matrika $A$ ireducibilna 
  glede na našo definicijo natanko tedaj, ko je nerazcepna v običajnem smislu. 

  Naj bo $A \geq 0$, saj taka matrika pusti stožec invarianten. Po~\ref{zgled2.20} 
  vemo, da so obrazi tega stožca oblike 
  $F_I = \{ x \in \mathbb{R}^n_+ ~|~ x_i > 0 \text{ natanko tedaj, ko } i \in I \}
  ~ \text{za } I \subseteq \{1, \ldots, n\}$.
  Obraz $F_I$ je invarianten za matriko $A$, če ima $x \in F_I$ ničle na indeksih $I$, 
  potem   jih mora imeti tudi $Ax$. Torej 
  ${(Ax)}_i = 0 ~\text{za vse } i \in I \text{ in vse } x \in F_I$. Če pogledamo 
  po komponentah, lahko zapišemo kot ${(Ax)}_i = \sum_{j=1}^n a_{ij} x_j = 0$, iz 
  česar pa sledi, da mora veljati $a_{ij} = 0 ~\text{za vse } i \in I,~ j \notin I$, 
  saj je $A$ nenegativna. Torej je $A$ po permutaciji indeksov obliko
  $A \sim
  \begin{bmatrix}
  * & * \\
  0 & *
  \end{bmatrix}$ in je razcepna.

  V drugo smer, če je $A$ razcepna, je po permutaciji indeksov oblike
  $A \sim
  \begin{bmatrix}
  B & C \\
  0 & D
  \end{bmatrix}$. Naj bo $I$ množica indeksov, ki ustreza spodnjemu bloku $D$.
  Vzemimo ploskev: $F_I = \{ x \in \mathbb{R}^n_+ \mid x_i = 0 \text{ za } i \in I \}$.
Za \( x \in F_I \) ima \( x \) obliko:
$x =
\begin{bmatrix}
x_1 \\
0
\end{bmatrix}$.
Potem:
\[
Ax =
\begin{bmatrix}
B & C \\
0 & D
\end{bmatrix}
\begin{bmatrix}
x_1 \\
0
\end{bmatrix}
=
\begin{bmatrix}
Bx_1 \\
0
\end{bmatrix}.
\]

Torej \( Ax \in F_I \), kar pomeni: $A(F_I) \subseteq F_I$, torej ni ireducibilna.
\end{zgled}

Naslednja trditev daje ekvivalenten način definiranja ireducibilnosti.

\begin{trditev}
  Linearni operator \(A : V \to V\), ki pusti strnjen zaprt stožec \(K\) invarianten, 
  je ireducibilen natanko tedaj, ko za nek \(\lambda > r(A)\) velja
\[
{(\lambda I - A)}^{-1}(K \setminus \{0\}) \subseteq \operatorname{int}(K).
\]
\end{trditev}
\begin{dokaz}
Recimo, da to ni res. Naj bo \(A\) ireducibilna in recimo, da obstaja
\(z \in K \setminus \{0\}\) tak, da za vsak \(\lambda > r(A)\) velja
\({(\lambda I - A)}^{-1}z \notin \operatorname{int}(K)\).
Zapišimo \(u = {(\lambda I - A)}^{-1}z\). 
Spomnimo se: ${(\lambda I - A)}^{-1}$ lahko z uporabo razvoja Neumannove vrste 
zapišemo kot ${(\lambda (I - \frac{1}{\lambda}A))}^{-1} = \sum_{k=0}^{\infty} \lambda^{-k-1} A^k $.
Po Hahn-Banachovem ločitvenem izreku~\cite[Theorem 11.6]{rockafellar} obstaja
\(\varphi \in K^* \setminus \{0\}\) tako, da je \(\varphi(u) = 0\) in
\(\varphi(x) > 0\) za vse \(x \in \operatorname{int}(K)\).
Definirajmo funkcional \(\psi \in V^*\) z
\[
\psi(x) = \varphi({(\lambda I - A)}^{-1}x)
        = \sum_{k=0}^{\infty} \lambda^{-k-1} \varphi(A^k x),
\]
in opazimo, da velja \(\psi(x) > 0\) za vse \(x \in \operatorname{int}(K)\), 
saj $A$ preslika $x \in K$ nazaj v notranjost stožca.
Očitno je \(\psi(z) = 0\), zato je množica


\[
F = \{x \in K : \psi(x) = 0\}
\]


pravi (izpostavljeni) obraz stožca \(K\).
Poleg tega za vsak \(x \in F\) velja


\[
0 \leq \psi(Ax)
   = \varphi(({(\lambda I - A)}^{-1})(Ax))
   = \sum_{k=0}^{\infty} \lambda^{-k-1} \varphi(A^{k+1}x) =
\]




\[
= \lambda \sum_{k=1}^{\infty} \lambda^{-k-1} \varphi(A^k x)
  \leq \lambda \psi(x) = 0.
\]


Od tod sledi, da je \(\psi(Ax) = 0\) za vse \(x \in F\).
Tako \(A\) pusti pravi obraz \(F\) invarianten, kar je prostislovje s predpostavko.

Za dokaz v drugo smer sedaj predpostavimo, da obstaja \(\lambda > r(A)\)
tak, da za vsak \(z \in K \setminus \{0\}\) velja
\({(\lambda I - A)}^{-1}z \in \operatorname{int}(K)\).
Naj obstaja pravi obraz \(F\) stožca \(K\) z lastnostjo \(A(F) \subseteq F\) 
in naj bo $z \in F$. Potem iz identitete
\[
{(\lambda I - A)}^{-1}
  = \sum_{k=0}^{\infty} \lambda^{-k-1} A^k
\]
sledi, da je 
\[
{(\lambda I - A)}^{-1} z = \sum_{k=0}^{\infty} \lambda^{-k-1} A^k z \in \operatorname{int}(K)
\]
torej je $z \in F \cap \operatorname{int}(K)$, iz česar pa sledi, da je 
$F \cap \operatorname{int}(K) \neq \varnothing$, kar pa je protislovje, saj je F 
pravi obraz.
\end{dokaz}

Opazimo, da nam ta dokaz pokaže tudi, da je linearni operator
\(A\) ireducibilen natanko tedaj, ko ne pusti nobenega pravega
izpostavljenega obraza stožca \(K\) invariantnega (?).

Kot smo omenili že prej, lahko izreka Perrona in Frobeniusa posplošimo
na linearne preslikave, ki pustijo stožec invarianten. To pomembno
ugotovitev sta naredila Krein in Rutman, ko sta preučevala linearne 
operatorje, ki pustijo stožec v (morda neskončnorazsežnem) 
normiranem prostoru invarianten.
Ker se v tej knjigi ukvarjamo le s končnorazsežnimi prostori, tukaj
podamo končnorazsežno različico njunega izreka.


\begin{izrek}[Krein-Rutman]\label{Krein-Rutman}
  Če je \(A : V \to V\) linearni operator, ki pusti strnjen
zaprt stožec \(K\) invarianten, potem je \(r(A)\) 
lastna vrednost operatorja \(A\) in
\(r(A)\) ima lastni vektor v \(K\).
\end{izrek}

Tako kot pri Perron-Frobeniusovem izreku 1.1.2 imamo tudi tukaj
enoličnost lastnega vektorja, če je \(A\) ireducibilen.


\begin{izrek}
  Če je \(A : V \to V\) ireducibilen linearni operator, ki pusti
  strnjen zaprt stožec \(K\) invarianten, potem velja:

  \begin{enumerate}
  \item \(r(A)\) je algebrajsko enostavna lastna vrednost operatorja
  \(A\), pripadajoči normirani lastni vektor \(v\) je enoličen in leži v
  \(\operatorname{int}(K)\). Poleg tega velja \(r(A) > 0\), če \(A \neq [0]\).

  \item Vsak nenegativen lastni vektor operatorja \(A\) je skalarni
  večkratnik vektorja \(v\).
  \end{enumerate}
\end{izrek}

Tretje točke Perron-Frobenius izreka~\ref{trd:per-fro} ni mogoče razširiti na poljubne
zaprte stožce. Najlažji protiprimer lahko vidimo v primeru Lorentzovega stožca $\Lambda_{3}$ in linearni
preslikavi $A : \mathbb{R}^{3} \to \mathbb{R}^{3}$, podani z

 

\[
A \begin{bmatrix}
x_{1} \\
x_{2} \\
x_{3}
\end{bmatrix}
=
\begin{bmatrix}
1 & 0 & 0 \\
0 & \cos\vartheta & -\sin\vartheta \\
0 & \sin\vartheta & \cos\vartheta
\end{bmatrix}
\begin{bmatrix}
x_{1} \\
x_{2} \\
x_{3}
\end{bmatrix}.
\]

 

Očitno velja $A(\Lambda_{3}) \subseteq \Lambda_{3}$ in $\sigma(A)=\{1, e^{\pm i\vartheta}\}$.
Preslikava $A$ ima tako tri lastne vrednosti in $r(A) = 1$.  Če je $\vartheta$ iracionalni 
večkratnik $2\pi$, potem $e^{i\vartheta}$ ni koren enote in točka (iii) ne drži.

Tretji del Perron-Frobeniusovega izreka~\ref{trd:per-fro} pa lahko posplošimo na poliedrske
stožce (glej~\cite{KreinRuthman}).

\begin{izrek}\label{izrek2.28}
Če je $A : V \to V$ linearna preslikava, ki ohranja strnjen poliedrski
stožec $K$ z $N$ ploskvami invarianten, potem za vsako $\lambda \in \sigma(A)$ za katero velja
$|\lambda| = r(A)$ obstaja $1 \le q \le N$ tako, da velja


\[
\lambda^{q} = {r(A)}^{q}.
\]


\end{izrek}

Izrek~\ref{izrek2.28} lahko uporabimo za dokaz naslednjega rezultata o iterativnem
obnašanju linearnih preslikav, ki ohranjajo poliedrski stožec.

\begin{izrek}
Če je $A : V \to V$ linearna preslikava, ki ohranja strnjen poliedrski
stožec $K$ z $N$ ploskvami invarianten, potem obstaja celo število $p \ge 1$ tako, da


\[
\lim_{k \to \infty} A^{kp} x
\]


obstaja za vsak $x \in K$, za katerega je zaporedje ${(\|A^{k}x\|)}_{k}$
omejeno. Poleg tega je $p$ red permutacije na $N$ elementih.
\end{izrek}

Za linearna preslikavo $A : V \to V$ je pogoj $A(K) \subseteq K$ ekvivalenten
trditvi, da iz $x \le_{K} y$ sledi $Ax \le_{K} Ay$ za $x, y \in V$. 

\begin{definicija}
  Naj bosta $K \subseteq V$ in $K' \subseteq V'$ stožca. 
  Preslikava $f : X \to X'$, kjer sta 
  $X \subseteq V$ in $X' \subseteq V'$, 
  \begin{itemize}
    \item \emph{ohranja red},
    če iz $x \le_{K} y$ sledi $f(x) \le_{K'} f(y)$ za $x, y \in X$,
    \item \emph{strogo ohranja red}, če iz $x \ll_{K} y$ sledi
    $f(x) \ll_{K'} f(y)$,
    \item \emph{obrača red}, če iz 
    $x \le_{K} y$ sledi $f(y) \le_{K'} f(x)$  za $x, y \in X$.
  \end{itemize}
\end{definicija}

Nelinearna Perron-Frobeniusova teorija se v prvi vrsti ukvarja 
preslikavami, ki ohranjajo red, in obravnava vprašanja, kot so:

\begin{itemize}
    \item Ali obstaja smiselna definicija spektralnega radija za
    preslikave $f : K \to K$, ki ohranjajo red, in ali obstaja ustrezen lastni
    vektor?
    \item Kdaj ima preslikava, ki ohranja red, lastni vektor v notranjosti
    stožca in kdaj je ta enoličen?
    \item Kako se dolgoročno obnašajo orbite
    $x, f(x), f^{2}(x), f^{3}(x), \ldots$ preslikave $f$, ki ohranja red?
    \item Kdaj ima preslikava, ki ohranja red, lastnost, da vsaka omejena
    orbita konvergira k periodični orbiti?
\end{itemize}

Ta vprašanja se pojavljajo v številnih problemih in so v središču
nelinearne Perron-Frobeniusove teorije. Kot bomo videli, obstajajo presenetljivo
podrobni odgovori za različne razrede preslikav, ki ohranjajo red.



\subsection{Stožci in delna urejenost}
Poglejmo si nekaj osnovnih lastnosti delne urejenosti, ki jo inducira stožec.
Naj bo $K \subseteq V$ stožec in $x, y \in V$ elementa vektorskega prostora $V$. 
Naslednja lema nam lahko pomaga določiti, kdaj velja $x \leq_K y$.

\begin{lema}\label{lema2.31}
Če je $K \subseteq V$ zaprt stožec, potem velja
\[
x \leq_K y \quad \Longleftrightarrow \quad \varphi(x) \leq \varphi(y)
\quad \text{za vse } \varphi \in K^*.
\]
\end{lema}

\begin{dokaz}
Če $x \not\leq_K y$, potem $y - x \notin K$. Ker je $K$ zaprt, 
konveksen stožec, po Hahn-Banachovem ločitvenem 
izreku~\cite[Theorem 11.6]{rockafellar} obstajata $\alpha \in \mathbb{R}$ in
$\psi \in V^*$ tako, da $\psi(y - x) < \alpha$ in
$\psi(v) > \alpha$ za vse $v \in K$. Ker je $K$ stožec, za vsak $\lambda \ge 0$ in
$v \in K$ velja $\lambda v \in K$. 
Če vzamemo $\lambda v$ namesto $v$, dobimo: $\psi(\lambda v) = \lambda \psi(v) > \alpha$.
Iz tega pa sledi $\psi(v) \ge 0$ za vse
$v \in K$, torej $\psi \in K^*$. Poleg tega je $\psi(0) = 0$, zato
$\alpha < 0$. Ker je $\psi$ linearna in $\psi(y-x) < 0$, tako dobimo $\psi(y) < \psi(x)$. 

Poglejmo še obratno implikacijo. Predpostavimo, da obstaja $\varphi \in K^*$, za katero velja
$\varphi(x) > \varphi(y)$. Tedaj je $\varphi(y - x) = \varphi(y) - \varphi(x) < 0$.
Ker pa je $\varphi \in K^*$, po definiciji velja $\varphi(v) \ge 0$ za vsak $v \in K$.
Če bi bilo $y - x \in K$, bi zato morali  $\varphi(y - x) \ge 0$,
kar je v protislovju z zgornjo strogo negativnostjo.
Torej $y - x \notin K$, kar pomeni $x \not\le_K y$.
\end{dokaz}

\begin{definicija}
  Pravimo, da element $y \in K$ \emph{dominira} element $x \in V$, če
  obstajata $\alpha, \beta \in \mathbb{R}$ tako, da velja
  $\alpha y \leq_K x \leq_K \beta y$.
\end{definicija}

Ta pojem določa ekvivalenčno
relacijo $\sim_K$ na $K$, kjer definiramo $x \sim_K y$ natanko tedaj,
ko $x$ dominira $y$ in $y$ dominira $x$. Zlahka preverimo, da
$x \sim_K y$ natanko tedaj, ko obstajata $0 < \alpha \le \beta$ tako,
da
\[
\alpha y \leq_K x \leq_K \beta y.
\]

\begin{definicija}
  Ekvivalenčne razrede v $K$ imenujemo \emph{deli} stožca in z
  $[x]$ označimo del, ki vsebuje $x$. Množico vseh delov stožca $K$
  označimo z $P(K)$.
\end{definicija}

\begin{zgled}
  Deli stožca $\mathbb{R}_+^2$ so
  $\{0\}$, $\operatorname{int}(\mathbb{R}_+^2)$,
  $\{(x_1, 0) \in \mathbb{R}_+^2~|~x_1 > 0\}$ in
  $\{(0, x_2) \in \mathbb{R}_+^2~|~x_2 > 0\}$.  
\end{zgled}

\begin{zgled}
  Poglejmo si še primer Lorentzovega stožca $\Lambda_3$. 
  Podobno kot v prejšnjem zgledu bosta
  $\{0\}$, $\operatorname{int}(\Lambda_3)$ ponovno dela stožca.
  Za rob stožca si pa poglejmo natančneje. Za vektor $x = (x_1, x_2, x_3)$ na robu velja 
  $x_1^2 = x_2^2 + x_3^2$. Če normaliziramo $x$ tako, da je $x_1 = 1$, potem dobimo,
  $1 = x_2^2 + x_3^2$. Ker pa je $x_2^2 + x_3^2 = 1$ krožnica, ki jo lahko parametriziramo 
  z enim kotom $\vartheta$, vidimo, da vsak robni vektor določa en kot $\vartheta$. 
  Hitro opazimo, da sta dva vektorja na robu ekvivalentna natanko tedaj, ko določata 
  isti kot $\vartheta$.
  Deli na robu so torej žarki, natančneje množice oblike
  \[  \{\lambda(1, \cos\vartheta, \sin\vartheta)~|~\lambda > 0\},
  \quad 0 \le \vartheta \le 2\pi.  \]
\end{zgled}



\begin{lema}\label{lema:deli_stozca}  
Če je $K \subseteq V$ zaprt stožec, potem so deli stožca $K$ natanko relativne notranjosti obrazov stožca $K$.
\end{lema}

\begin{dokaz}
Znano je (glej~\cite[Izrek 8.2]{rockafellar}), da relativne notranjosti obrazov stožca $K$ tvorijo razčlenitev $K$. 
Naj bo $[x]$ del elementa $x$ v $K$ in naj bo $F_x$ obraz stožca $K$, katerega relativna notranjost vsebuje $x$. 
Po trditvi~\ref{trditev2.15} je $F_x$ zaprt in hitro lahko preverimo, da je tudi stožec.

Predpostavimo, da je $y$ v relativni notranjosti $F_x$ in dokažimo, da je $x \sim_K y$. 
Ker je v relativni notranjosti, obstaja $\alpha > 0$ tako, da sta 
$y - \alpha x \in F_x$ in $x - \alpha y \in F_x$. 
Ker je \( F_x \subseteq K \), to pomeni:
$y - \alpha x \in K \;\Rightarrow\; \alpha x \le y$
in
$x - \alpha y \in K \;\Rightarrow\; \alpha y \le x$.
Od tod sledi $\alpha x \leq y \leq \frac{1}{\alpha} x$, in zato $x \sim_K y$.

V drugo smer: naj bo $y \in K$ tak, da $x \sim_K y$.
Naj bo $F_x$ obraz stožca $K$, katerega relativna notranjost vsebuje $x$. Ker je $x$ v relativni notranjosti $F_x$, lahko naredimo majhne 
premike od $x$ v smeri $y$, ne da bi zapustili $F_x$.
Za dovolj majhen $\delta > 0$ velja $x - \delta y \sim_K y$ in $y - \delta x \sim_K x$. 
Hitro lahko preverimo, da iz $x \sim_K y$ sledi tudi $\lambda x \sim_K \mu y$ za vse $\lambda, \mu > 0$.  
Za dovolj majhen $\varepsilon > 0$ definiramo $x_\varepsilon = (1+\varepsilon)x - \varepsilon y$ in $y_\varepsilon = (1+\varepsilon)y - \varepsilon x$.
Od tod sledi, da velja
\[
x_\varepsilon \sim_K y
\qquad\text{in}\qquad
y_\varepsilon \sim_K x.
\]
Opazimo, da velja
\[
x = a x_\varepsilon + (1-a) y_\varepsilon,
\qquad
a = \frac{1+\varepsilon}{1+2\varepsilon}, \text{ kjer je } 0 < a < 1.
\]
Ker je $x \in F_x$, iz definicije obraza sledi, da sta $x_\varepsilon$ in $y_\varepsilon$ v $F_x$.
Ker pa velja
\[
y = \lambda y_\varepsilon + (1-\lambda)x,
\qquad
\lambda = \frac{1}{1+\varepsilon}, \text{ kjer je } 0 < \lambda < 1,
\]
iz definicije obraza enako zaključimo, da je $y$ v relativni notranjosti $F_x$.
\end{dokaz}

Iz leme~\ref{lema:deli_stozca} sledi, da ima poliedrski stožec končno mnogo delov. 
\begin{definicija}
  Naj bo dan poliedrski stožec \(K\) z \(N\) ličnicami in funkcionali
  \(\psi_1,\ldots,\psi_N\), ki določajo ličnice. Za vsak \(x\in K\) definiramo množico
  \[
  I_x=\{\, i : \psi_i(x)>0 \,\}.
  \]
  Če je \(P\) del stožca \(K\), definiramo
  \[
  I(P)=\{\, i : \psi_i(x)>0 \text{ za nek } x\in P \,\}.
  \]
\end{definicija}


Množica delov \(\mathcal{P}(K)\) ima naravni delni red \(P \trianglelefteq Q\), podan z
\[
P \trianglelefteq Q
\quad \text{, če obstajata } x\in P \text{ in } y\in Q \text{ takšna, da }
y \text{ dominira } x.
\]

Ekvivalentno: \(P \trianglelefteq Q\) če za vsak \(x\in P\) in vsak \(y\in Q\)
velja, da \(y\) dominira \(x\).

\begin{lema}\label{lema2.38}
Če je \(K\subseteq V\) poliedrski stožec z \(N\) ličnicami, potem veljajo:
\begin{enumerate}[label= (\roman*)]
\item \(I_x=I_y\) natanko tedaj, ko \(x\sim_K y\).
\item Za vsak \(P\in\mathcal{P}(K)\) velja

\[
P=\{\, x\in K : \psi_i(x)>0 \text{ natanko tedaj, ko } i\in I(P) \,\}.
\]


\item \(P\trianglelefteq Q\) natanko tedaj, ko \(I(P)\subseteq I(Q)\).
\item \(|\mathcal{P}(K)|\le 2^N\).
\end{enumerate}
\end{lema}

\begin{dokaz}
Naj bodo \(\psi_1,\ldots,\psi_N\in K^*\) funkcionali, ki določajo ličnice stožca
\(K\). 

(i) Po lemi~\ref{lema2.19} vemo, da je \(x\le_K y\) ekvivalentno pogoju
      \(\psi_i(x)\le\psi_i(y)\) za vse \(i\). Iz tega sledi
      $I_x\subseteq I_y$ natanko tedaj, ko $y$ dominira $x$.
      Zato je \(x\sim_K y\) ekvivalentno \(I_x=I_y\). 

(ii) Če sta $x$ in $y$ v istem delu stožca $P$, potem velja $x \sim_K y$ in po (i) sledi $I_x = I_y$. 
      Če pa je $x \in P$ in $y \in K$ z lastnostjo, da velja $\psi_i(x) > 0$ natanko tedaj, ko 
      $i \in I(P)$, potem je $y \in P$, saj je $y \sim_K x$. Torej $x\in P$ natanko tedaj, ko $I_x=I(P)$.

(iii) 
      Naj bo sedaj \(P\trianglelefteq Q\). Tedaj obstajata \(x\in P\) in \(y\in Q\), takšna da
      \(y\) dominira \(x\). Tako dobimo $I(P)=I_x\subseteq I_y=I(Q)$.
      V obratno smer: če \(I(P)\subseteq I(Q)\), potem za vsak \(x\in P\) in \(y\in Q\) obstaja
      \(\beta>0\) tako, da $\psi_i(x)\le \beta\,\psi_i(y)$ za vse $i$.
      Po lemi~\ref{lema2.19} sledi, da \(y\) dominira \(x\), zato \(P\trianglelefteq Q\). 
      
(iv)      
      Iz točke (iii) sledi, da je vsak del stožca \(P\) določen z množico indeksov \(I(P)\subseteq\{1,\ldots,N\}\).
      Ker je število podmnožic množice z \(N\) elementi natanko \(2^N\), sledi, da je \(|\mathcal{P}(K)|\le 2^N\).
\end{dokaz}

Oblika stožca določa številne pomembne lastnosti, ki so ključne v nelinearni
Perron-Frobeniusovi teoriji. Pomemben razred stožcev so strogo konveksni stožci.

\begin{definicija}
Zaprt stožec $K\subseteq V$ je \emph{strogo konveksen}, če za vse
$x,y\in\partial K\setminus\{0\}$, kjer $x\neq \alpha y$ za noben $\alpha>0$,
velja
\[
\{\lambda x+(1-\lambda)y:0<\lambda<1\}\subseteq\operatorname{int}(K).
\]
\end{definicija}

\begin{zgled}
  Lorentzov stožec $\Lambda_{n+1}$ je primer strogo konveksnega stožca.
\end{zgled}

\begin{zgled}
  Stožec pozitivno semidefinitnih matrik $\Pi_n(\mathbb{R})$ za $n\ge 3$ je
  pomemben primer stožca, ki ni strogo konveksen. Poglejmo si na primeru, zakaj $\Pi_3(\mathbb{R})$ ni strogo konveksen.
  Spomnimo se, da so pozitivno semidefinitne matrike simetrične matrike, za katere velja, da $x^T A x \ge 0$ za vsak vektor $x$.
  Za tiste na robu stožca to pomeni, da obstaja vektor $x$ tako, da $x^T A x = 0$.
  Vzemimo matriki
  \[
  A =
  \begin{pmatrix}
  1 & 0 & 0\\
  0 & 0 & 0\\
  0 & 0 & 0
  \end{pmatrix},
  \qquad
  B =
  \begin{pmatrix}
  0 & 0 & 0\\
  0 & 1 & 0\\
  0 & 0 & 0
  \end{pmatrix}.
  \]
  Obe sta očitno pozitivno semidefinitni in ker velja
  $\det A = 0$ in $\det B = 0$,
  sta singularni, zato ležita na robu:
  $A,B \in \partial \Pi_3(\mathbb{R})$.
  Poleg tega $A$ in $B$ očitno nista skalarni večkratnici.
  Za $0<\lambda<1$ poglejmo konveksno kombinacijo
  \[
  C_\lambda = (1-\lambda)A + \lambda B
  =
  \begin{pmatrix}
  1-\lambda & 0 & 0\\
  0 & \lambda & 0\\
  0 & 0 & 0
  \end{pmatrix}.
  \]
  Za poljuben $x = (x_1,x_2,x_3)^\top$ dobimo
  \[
  x^\top C_\lambda x
  = (1-\lambda)x_1^2 + \lambda x_2^2 \ge 0,
  \]
  torej $C_\lambda$ je pozitivno semidefinitna matrika, $C_\lambda \in \Pi_3(\mathbb{R})$.  
  Po drugi strani pa za vektor $x = (0,0,x_3)^\top$ velja
  \[
  x^\top C_\lambda x = 0,
  \]
  torej je $C_\lambda$ singularna in zato leži na robu stožca: $C_\lambda \in \partial \Pi_3(\mathbb{R})$.
  Po definiciji to pomeni, da stožec $\Pi_3(\mathbb{R})$ ni strogo konveksen.
\end{zgled}

\begin{definicija}
  Naj bo $K$ stožec v vektorskem prostoru $V$ in $S\subseteq V$. Če je $u$ tak, da za vse $s\in S$ velja $s\le_K u$, 
  pravimo, da je $u$ \emph{zgornja meja} množice $S$. 
  Če ima vsaka zgornja meja $v \in V$ množice $S$ lastnost $u\le_K v$, $u$ imenujemo \emph{supremum} množice $S$ in
  označimo $u=\sup(S)$.

  Podobno definiramo \emph{spodnjo mejo} in \emph{infimum} množice $S$.
\end{definicija}

Označimo $x\wedge y=\inf(x,y)$ in $x\vee y=\sup(x,y)$.

\begin{zgled}
  Če je $K=\mathbb{R}^n_+$, potem je
  $\sup(x,y)_i=\max\{x_i,y_i\}$ in $\inf(x,y)_i=\min\{x_i,y_i\}$ za vsak $i = 1,\ldots,n$.
\end{zgled}

V splošnem pa supremum $\sup(x,y)$ in infimum $\inf(x,y)$ za poljubna $x,y \in V$ 
nista nujno definirana. 

\begin{definicija}
  Stožci, za katere $\sup(x,y)$ in $\inf(x,y)$ obstajata za vse 
  $x,y \in V$, se imenujejo \emph{miniedrski}.
\end{definicija}

Kot dejstvo bomo prevzeli, da je strnjen zaprt miniedrski stožec enak 
$\mathbb{R}_{+}^{n}$. V~\cite{KreinRuthman,deSzNagy1945} je bilo pokazano, 
da je strnjen zaprt stožec $K$ v $n$-dimenzionalnem realnem vektorskem prostoru $V$ 
miniedrski natanko tedaj, ko je $K$ \emph{simplicialen}, tj.\ obstajajo linearno 
neodvisni vektorji $x^{1},\dots,x^{n} \in V$ tako, da je
\[
K = \left\{ x \in V~\Big|~x = \sum_{i=1}^{n} \alpha_{i} x^{i},\ \alpha_{i} \ge 0 
\text{ za } 1 \le i \le n \right\}.
\]
Poglejmo si sedaj nekaj osnovnih dejstev o dualnih stožcih in normalnosti. 

\begin{definicija}
  \emph{Dualni stožec} stožca $K \subseteq V$ je definiran kot
  \[
  K^{*} = \{ \varphi \in V^{*}~|~\varphi(x) \ge 0 \text{ za vse } x \in K \}.
  \]
\end{definicija}

Na splošno $K^{*}$ ni nujno stožec, vendar pa velja naslednje.

\begin{lema}\label{lema2.46}
Če je $K \subseteq V$ strnjen zaprt stožec, potem je $K^{*}$ tudi strnjen zaprt stožec. Poleg tega, 
če $\varphi \in \operatorname{int}(K^{*})$, potem $\varphi(x) > 0$ za vse $x \in K$ 
kjer $x \ne 0$, in množica
\[
\Sigma_{\varphi} = \{ x \in K : \varphi(x) = 1 \}
\]
je kompaktno konveksna podmnožica prostora $V$.
\end{lema}
\begin{dokaz}
  TO DO!
\end{dokaz}

\begin{definicija}
Naj bo $K \subseteq V$ stožec in naj bo $V$ opremljen z normo $\|\cdot\|$.  
Pravimo, da je $K$ \emph{normalen}, če obstaja konstanta $\delta > 0$ taka, da iz $0 \leq x \leq y$ sledi
\[
\|x\| \leq \delta \, \|y\|.
\]
Infimum vseh takih $\delta > 0$ imenujemo \emph{konstanta normalnosti}.  
Normo $\|\cdot\|$ imenujemo \emph{monotona norma} za $K$, če je konstanta normalnosti enaka $1$.
\end{definicija}

\begin{lema}
Vsak zaprt stožec $K$ v normiranem prostoru $(V, \|\cdot\|)$ je normalen.
\end{lema}
\begin{dokaz}
  Predpostavimo, da $K$ ni normalen. Potem za vsak $ \delta > 0$ obstajata zaporedji 
  $(x_k)_k$ in $(y_k)_k$ v $K$ tako, da velja $0 \leq x_k \leq y_k$ in  
  $\|x_k\| > \delta \, \|y_k\|$. Ker to velja za vsak $\delta > 0$, lahko izberemo $\delta = k$ 
  in dobimo zaporedji $(x_k)_k$ in $(y_k)_k$ v $K$ tako, da velja $0 \leq x_k \leq y_k$ in  
  $\|x_k\| > k \, \|y_k\|$. Od tod sledi, da $\frac{\|y_k\| }{\|x_k\| } < \frac{1}{k}$, kar 
  pomeni, da $\lim\limits_{k \to \infty} \frac{\|y_k\| }{\|x_k\| }  = 0$. Brez škode za splošnost 
  lahko predpostavimo, da je $\|x_k\| = 1$ za vsak $k$. Iz limite torej sledi, da 
  $\lim\limits_{k \to \infty} \|y_k\| = 0$ in posledično $\lim\limits_{k \to \infty} y_k = 0$. 
  Ker so vsi $x_k$ del enotske sfere in je $V$ končnorazsežen, iz Bolzano-Weierstrassovega izreka 
  sledi, da obstaja podzaporedje $(x_{k_j})_j$ tako, 
  da $\lim\limits_{j \to \infty} x_{k_j} = x$ za nek $x \in K$. Ker je $K$ zaprt, sledi, da je 
  $x \in K$. Dalje velja $x_{k_j} \leq y_{k_j}$, 
  torej je $y_{k_j} - x_{k_j} \in K$. Iz $\lim\limits_{k \to \infty} y_k = 0$ in $\lim\limits_{j \to \infty} x_{k_j} = x$
  sledi, da je $\lim\limits_{j \to \infty} (y_{k_j} - x_{k_j}) = -x$. Ker je $K$ zaprt, sledi, da je $-x \in K$. 
  Po definiciji stožca je $ K \cap (-K) = \{0\}$, zato je $x = 0$. To pa je v protislovju z dejstvom, da je $\|x\| = 1$.
\end{dokaz}



\subsection{Preslikave, ki ohranjajo red}

Eden izmed najpomembnejših pojmov v nelinearni Perron-Frobeniusovi teoriji so preslikave, ki ohranjajo red,
$f : K \to K$, kjer je $K$ strnjen zaprt stožec v končnorazsežnem vektorskem prostoru $V$.
Poleg same definicije lahko za ugotavljanje, ali preslikava
$f : K \to K$
ohranja red glede na stožec $K$, uporabimo tudi Fréchetov odvod
$Df(x)$ v točki $x \in \operatorname{int}(K)$.
Spomnimo se sedaj definicije Fréchetovega odvoda.
\begin{definicija}
  Naj bo  $U \subseteq V$ odprta množica. Pravimo, da je preslikava
  $f : U \to V$ \emph{Fréchetovo odvedljiva} v točki $x \in U$, če obstaja linearna preslikava
  $L : V \to V$, za katero velja
  \[
  \lim_{h \to 0}
  \frac{\|f(x+h)-f(x)-Lh\|}{\|h\|}
  =0.
  \]
  Fréchetov odvod preslikave $f$ v točki $x$ označimo z
  \[
  Df(x)=L.
  \]
\end{definicija}

Spomnimo se sedaj nekaterih definicij.~\cite{Evans}
\begin{definicija}
Naj bo $A \subseteq \mathbb{R}^n$. Preslikava $f \colon A \to \mathbb{R}^m$
je \emph{Lipschitzevo zvezna}, če obstaja konstanta $C \geq 0$, da za vsak par točk $x, y \in A$ velja
\[
\|f(x)-f(y)\| \leq C\,\|x-y\|.
\]
\end{definicija}

Za dokaz naslednjega izreka bomo uporabili nekoliko naprednejše rezultate iz realne analize.
Natančneje, uporabili bomo Rademacherjev izrek~\cite{Rademacher}.
\begin{izrek}[Rademacherjev izrek]
  Vsaka lokalno Lipschitzeva preslikava
  $f : U \to V$ je 
  Fréchetovo odvedljiva za skoraj vse točke $x \in U$ glede na Lebesgueovo mero, norma odvoda
  $\|Df(x)\|$
  pa je enakomerno omejena na kompaktnih podmnožicah množice $U$.
\end{izrek}

Iz Radamacherjevega izreka sledi, da je preslikava $x \mapsto Df(x)$ Lebesgueovo merljiva. 
Zaradi poenostavitve naslednjega dokaza predpostavimo, da je preslikava
$f$ Fréchetovo odvedljiva na celotnem $K$.

\begin{izrek}\label{izrLipFreLeb}
Naj bo $K \subseteq V$ poln zaprt stožec in naj bo $U \subseteq \operatorname{int}(K)$ odprta konveksna množica.
Če je preslikava
$f : U \to K$
lokalno Lipschitzeva, potem Fréchetov odvod $Df(x)$ obstaja za skoraj vse
$x \in U$
glede na Lebesgueovo mero.
Poleg tega preslikava $f$ ohranja red glede na stožec $K$
natanko tedaj, ko za vsak $x \in U$, v katerem odvod obstaja, velja
\[
Df(x)(K)\subseteq K.
\]
\end{izrek}
\begin{dokaz}
  TO DO.
\end{dokaz}

Sedaj si bomo pogledali subhomogene preslikave, ki so pomemben razred preslikav, 
ki ohranjajo red. 
\begin{definicija}
  Naj bosta $K \subseteq V$ in $K' \subseteq V'$ stožca. Preslikava
  $f \colon K \to K'$ je 
  \begin{itemize}
    \item \emph{subhomogena}, če za vsak $x \in K$ in vsak $0 < \lambda < 1$ velja
      $\lambda f(x) \leq f(\lambda x)$,
    \item \emph{strogo subhomogena}, če za vsak
      $x \in K \setminus \{0\}$ in vsak $0 < \lambda < 1$ velja
      $\lambda f(x) \ll f(\lambda x)$,
    \item \emph{homogena stopnje $p$}, če za vsak
      $x \in K$ in vsak $\lambda > 0$ velja
      $f(\lambda x)=\lambda^p f(x)$.
    \item \emph{homogena}, če je homogena stopnje $p=1$.
  \end{itemize}
\end{definicija}





\section{Neekspanzivnost}

V tem poglavju bomo preučevali neekspanzivnost preslikav, ki ohranjajo red, glede na metriki, 
ki jih bomo definirali v nadaljevanju. Pokazali bomo, da sta Hilbertova in Thompsonova metrika
naravni metriki za proučevanje homogenih oziroma subhomogenih
preslikav, ki ohranjajo red na stožcih. Kot bomo videli, so homogene
preslikave, ki ohranjajo red, neekspanzivne glede na Hilbertovo in
Thompsonovo metriko, medtem ko so strogo subhomogene preslikave celo
skrčitev.

\subsection{Hilbertova in Thompsonova metrika}

V tem razdelku bomo pokazali, kako so homogene in subhomogene preslikave, ki ohranjajo red, povezane s preslikavami, 
ki niso ekspanzivne glede na Hilbertovo (projektivno) metriko in Thompsonovo (delno) metriko.

\begin{definicija}
  Naj bo $K$ stožec v vektorskem prostoru $V$. Za elementa
  $x,y\in K$, za katera velja $x \sim_K y$, $ y\neq0$, definiramo
  \[
  M(x/y)=\inf\{\beta>0 \mid x\le_K \beta y\},
  \]
  in
  \[
  m(x/y)=\sup\{\alpha>0 \mid \alpha y\le_K x\}.
  \]
\end{definicija}

Ko je stožec $K$ jasen iz konteksta, bomo enostavno zapisali $M(x/y)$ in $m(x/y)$.
Na podlagi teh količin definiramo Hilbertovo metriko.

\begin{definicija}
Naj bo $K$ stožec v vektorskem prostoru $V$.
\emph{Hilbertova projektivna metrika} je preslikava
$d_H:K\times K\longrightarrow[0,\infty]$,
podana z
\[
d_H(x,y)=
\log\!\left(
\frac{M(x/y)}{m(x/y)}
\right),
\]
za vse $x\sim_K y$ z $y\neq0$.

Poleg tega definiramo še $d_H(0,0)=0$,
v vseh ostalih primerih pa postavimo $d_H(x,y)=\infty$.
\end{definicija}

Naslednja trditev povzema osnovne lastnosti Hilbertove metrike.

\begin{trditev}\label{trd3.3}
Če je $K$ stožec v vektorskem prostoru $V$, potem velja:
\begin{enumerate}
\item $d_H(x,y) \ge 0$ in $d_H(x,y)=d_H(y,x)$ za vse $x,y\in K$,
\item $d_H(x,z) \le d_H(x,y)+d_H(y,z)$ za vse $x\sim_K y\sim_K z$,
\item Za vse $\alpha,\beta>0$ velja $d_H(\alpha x,\beta y)=d_H(x,y)$.
\end{enumerate}
\end{trditev}

\begin{dokaz}
Za dokaz prve trditve zadostuje obravnavati primer $x\sim_K y$ in $y\neq0$, saj ostali primeri sledijo 
neposredno iz definicije. V tem primeru je tudi $x\neq0$. Za vsaka $0<\alpha<m(x/y)$ in $0<M(x/y)<\beta$ 
velja $\alpha y\le_K x\le_K\beta y$, zato sledi $y\le_K(\beta/\alpha)y$. Od tod dobimo $\beta/\alpha\ge1$ 
za vse $0<\alpha<m(x/y)$ in $0<M(x/y)<\beta$. Torej je $M(x/y)/m(x/y)\ge1$, od koder sledi $d_H(x,y)\ge0$.

Ker je $\alpha$ pozitiven skalar, lahko hitro opazimo še, da z deljenjem z $\alpha$ dobimo 
$M(x/y)=\inf\{\alpha>0~|~x\le_K \alpha y\}=\inf\{\alpha>0~|~\alpha^{-1}x\le_K y\}=m(y/x)^{-1}$.
Na podoben način dobimo $m(x/y)=M(y/x)^{-1}$. Od tod sledi $d_H(x,y)= 
\log\!\left(\frac{M(x/y)}{m(x/y)}\right)=\log\!\left(\frac{m(y/x)^{-1}}{M(y/x)^{-1}}\right)=
\log\!\left(\frac{M(y/x)}{m(y/x)}\right)=d_H(y,x)$.

Za dokaz trikotniške neenakosti naj velja $x\sim_K y\sim_K z$ v $K$ in $z\neq0$. Za vsaka $0<\alpha<m(x/y)$ 
in $0<\gamma<m(y/z)$ velja $\alpha y\le_K x$ in $\gamma z\le_K y$, zato $\alpha\gamma z\le_K x$. Od tod 
sledi $m(x/z)\ge m(x/y)m(y/z)$. Na podoben način dobimo $M(x/z)\le M(x/y)M(y/z)$. Iz teh dveh neenakosti sledi
\[
\frac{M(x/z)}{m(x/z)}
\le
\frac{M(x/y)M(y/z)}
{m(x/y)m(y/z)}.
\]
Uporabimo dejstvo, da je logaritemska funkcija strogo naraščajoča ter da je logaritem produkta vsota logaritmov 
in dobimo $d_H(x,z)\le d_H(x,y)+d_H(y,z)$.

Tretja trditev sledi iz dejstva, da
\[
M(\alpha x/\beta y)=\frac{\alpha}{\beta}M(x/y)
\qquad\text{in}\qquad
m(\alpha x/\beta y)=\frac{\alpha}{\beta}m(x/y)
\]
za vse $\alpha,\beta>0$, saj $d_H(\alpha x, \beta y)= 
\log\!\left(\frac{M(\alpha x/\beta y)}{m(\alpha x/\beta y)}\right) = 
\log\!\left(\frac{\frac{\alpha}{\beta}M(x/y)}{\frac{\alpha}{\beta}m(x/y)}\right) = d_H(x,y)$.
\end{dokaz}

Hilbertova projektivna metrika pa ni prava metrika, saj na stožcu $K$ ne velja, da je 
$d_H(x,y)=0$ natanko tedaj, ko je $x=y$. Lahko pa pokažemo, da je Hilbertova metrika prava metrika 
na $\Sigma=P\cap X$, kjer je $P$ del stožca $K$, množica $X$ pa vsebuje
natanko enega predstavnika iz vsakega pozitivnega žarka.

\begin{izrek}\label{izr3.4}
  Naj bo $K$ zaprt stožec v vektorskem prostoru $V$. Potem velja $d_H(x,y)=0$
  natanko tedaj, ko obstaja $\lambda>0$, da je $x=\lambda y$.
  V tem primeru, če je $X\subseteq K$ taka množica, da za vsak
  $x\in K\setminus\{0\}$ obstaja natanko en $\lambda>0$, za katerega je
  $\lambda x\in X$ in je $P$ del stožca $K$, potem je $\Sigma=P\cap X$ 
  pravi metrični prostor glede na Hilbertovo metriko.
\end{izrek}
\begin{dokaz}
  Naj bo $K$ zaprt in naj velja $x\sim_K y$ v $K$, kjer je $y\neq0$. Ker je $K$ zaprt, sta supremum in infimum 
  v definicijah $M$ in $m$ dosežena, zato velja 
  \[
  m(x/y)y\le_K x\le_K M(x/y)y. \tag{2.1}
  \]
  Torej iz $d_H(x,y)=0$ sledi $\frac{M(x/y)}{m(x/y)}=1$, zato je
  \[
  y\le_K m(x/y)^{-1}x\le_K \frac{M(x/y)}{m(x/y)}y=y. 
  \]
  Iz tega sledi, da je $y = m(x/y)^{-1}x$ torej tak $\lambda>0$ obstaja. 
  
  Za obratno implikacijo uporabimo točko 3 prejšnje trditve.
  Ker je $x=\lambda y$ za nek $\lambda>0$, velja $d_H(x,y)=d_H(\lambda y,y)= d_H(y,y)=0$.

  Dokažimo še zadnjo trditev. Naj bosta $x,y\in\Sigma=P\cap X$. Če je $d_H(x,y)=0$, potem po prvem
  delu izreka obstaja $\lambda>0$, za katerega je $x=\lambda y$.
  Ker sta $x,y\in X$, velja $1\cdot y=y\in X$ in
  $\lambda y=x\in X$. Po predpostavki pa za vsak
  $y\in K\setminus\{0\}$ obstaja natanko en pozitiven skalar $\mu$,
  za katerega je $\mu y\in X$. Zato je $\lambda=1$ in posledično
  $x=y$.
  Obratno, če je $x=y$, potem je $d_H(x,y)=d_H(x,x)=0$.
  Torej za vse $x,y\in\Sigma$ velja $d_H(x,y)=0 \Longleftrightarrow x=y$.
\end{dokaz}

Poglejmo si, zakaj je bilo v prejšnjem izreku nujno predpostaviti, da je stožec $K$ zaprt. 
\begin{zgled}
  Vzemimo stožec, ki ni zaprt,
  \[
  K=\{x\in\mathbb{R}^2~|~x_1>0\}\cup\{(0,0)\}.
  \]
  Potem velja $x\sim_K y$ za vse $x,y\in K\setminus\{0\}$. 
  Naj bosta $x=(x_1,x_2)$ in $y=(y_1,y_2)\in K\setminus\{0\}$.
  Tedaj sta $x_1,y_1>0$. Označimo $r=\frac{x_1}{y_1}$.
  Ker je $\beta y-x=(y_1(\beta-r),\,\beta y_2-x_2)$, velja
  $x\le_K\beta y$ za vsak $\beta>r$. Zato je $M(x/y)=r=x_1/y_1$.
  Podobno je $x-\alpha y=(y_1(r-\alpha),\,x_2-\alpha y_2)$, zato velja
  $\alpha y\le_K x$ za vsak $\alpha<r$. Torej je $m(x/y)=r=x_1/y_1$.
  Sledi $M(x/y)=\frac{x_1}{y_1}=m(x/y)$, zato je
  $d_H(x,y)=\log\left(\frac{M(x_1/y_1)}{m(x_1/y_1)}\right)=0$ 
  ne glede na to, ali je $x$ skalarni večkratnik $y$ ali ne.
\end{zgled}

Pomembno različico Hilbertove metrike na stožcih je uvedel Thompson.

\begin{definicija}
Naj bo $K$ stožec v vektorskem prostoru $V$.
\emph{Thompsonova metrika} je preslikava
$d_T:K\times K\longrightarrow[0,\infty]$, podana z
\[
d_T(x,y)=\log\bigl(\max\{M(x/y),M(y/x)\}\bigr)
\]
za vse $x\sim_K y$ z $y\neq0$.
Poleg tega definiramo $d_T(0,0)=0$, v vseh ostalih primerih pa
postavimo $d_T(x,y)=\infty$.
\end{definicija}

S podobnimi argumenti kot v dokazih osnovnih lastnosti Hilbertove
metrike~\ref{trd3.3} in~\ref{izr3.4} lahko pokažemo, da je $d_T$ prava metrika na vsakem delu
stožca. (TO DO, napisati trditev in dokaz)

Hilbertova in Thompsonova metrika sta posebej uporabni pri preučevanju
homogenih in subhomogenih preslikav na stožcih, ki ohranjajo red.
Preden to podrobneje obravnavamo, si poglejmo nekaj osnovnih pojmov.

\begin{definicija}
Naj bosta $(X,d_X)$ in $(Y,d_Y)$ metrična prostora. Preslikava
$f:X\longrightarrow Y$ je \emph{neekspanzivna}, če za vse $x,y\in X$
velja
\[
d_Y\bigl(f(x),f(y)\bigr)\leq d_X(x,y).
\tag{2.4}
\]
Preslikava $f$ je \emph{skrčitev}, če je neenakost v~\textup{(2.4)}
stroga za vse različne $x,y\in X$.

Preslikava $f$ je \emph{Lipschitzeva skrčitev}, če obstaja 
$c\in[0,1)$, tako da za vse $x,y\in X$ velja
$d_Y(f(x),f(y))\leq c\,d_X(x,y)$.
Preslikava $f$ je \emph{izometrija}, če v~\textup{(2.4)} velja enakost
za vse $x,y\in X$.
\end{definicija}

\begin{trditev}
Naj bosta $K\subseteq V$ in $K'\subseteq V'$ zaprta stožca. Če je
$f:K\longrightarrow K'$ preslikava, ki ohranja red in je homogena
stopnje $r>0$, potem za vse $x,y\in K$, za katere velja $x\sim_K y$,
velja
\[
m(x/y)^r\leq m\bigl(f(x)/f(y)\bigr)
\qquad\text{in}\qquad
M\bigl(f(x)/f(y)\bigr)\leq M(x/y)^r.
\]
\end{trditev}

\begin{dokaz}
Naj bo $x\sim_K y$. Ker je $K$ zaprt stožec, velja
$m(x/y)y\leq_K x\leq_K M(x/y)y$. Ker $f$ ohranja red, dobimo 
$f(m(x/y) y)\leq_{K'}f(x)\leq_{K'} f(M(x/y) y)$ in ker je $f$ homogena stopnje $r>0$, velja
$m(x/y)^r f(y)\leq_{K'} f(x)\leq_{K'} M(x/y)^r f(y)$.
Po definiciji je $m (f(x)/f(y)) = \sup\{\alpha>0: \alpha f(y)\leq_{K'} f(x)\}$, torej je $m(x/y)^r\leq m(f(x)/f(y))$. Na podoben način je
dobimo, da je $M(f(x)/f(y))\leq M(x/y)^r$.
\end{dokaz}

Prejšnja trditev ima naslednjo posledico.

\begin{posledica}\label{posledica3.9}
Naj bosta $K\subseteq V$ in $K'\subseteq V'$ zaprta stožca. Če je
$f:K\longrightarrow K'$ preslikava, ki ohranja red in je homogena
stopnje $r>0$, potem za vse $x,y\in K$, za katere velja $x\sim_K y$,
velja
\[
d_H\bigl(f(x),f(y)\bigr)\leq r\,d_H(x,y)
\qquad\text{in}\qquad
d_T\bigl(f(x),f(y)\bigr)\leq r\,d_T(x,y).
\]
\end{posledica}
\begin{dokaz}
  Po prejšnji trditvi in lastnostih logaritma velja
  \[
  d_H(f(x),f(y))=\log\left(\frac{M(f(x/f(y)))}{m(f(x)/f(y))}\right)
  \leq
  \log\left(\frac{M(x/y)^r}{m(x/y)^r}\right)
  =r\,d_H(x,y)
  \]
  in
  \[
  \begin{aligned}
    d_T(f(x),f(y))
    &=
    \log\left(
    \max\left\{
    M\bigl(f(x)/f(y)\bigr),
    M\bigl(f(y)/f(x)\bigr)
    \right\}
    \right)\\
    &\leq
    \log\left(
    \max\left\{
    M(x/y)^r,
    M(y/x)^r
    \right\}
    \right)\\
    &=
    r\,d_T(x,y).
  \end{aligned}
  \]
\end{dokaz}

Posebej je vsaka homogena preslikava $f:K\longrightarrow K'$, ki
ohranja red, neekspanzivna glede na Hilbertovo in Thompsonovo metriko.

Iz prejšnje posledice sledi, da je bijektivna linearna preslikava
$f:K\longrightarrow K'$ izometrija glede na Hilbertovo in Thompsonovo
metriko. Res, če je $f$ linearna, potem velja $f(\lambda x)=\lambda f(x)$, torej je homogena 
stopnje $r=1$, zato velja $d_H(f(x),f(y))\leq d_H(x,y)$ in $d_T(f(x),f(y))\leq d_T(x,y)$. Ker je $f$ bijektivna, 
je $f^{-1}$ tudi homogena stopnje $r=1$ in po prejšnji posledicivelja
$d_H(x,y) = d_H(f^{-1}(f(x)),f^{-1}(f(y))) \leq d_H(f(x),f(y))$ in $d_T(x,y)\leq d_T(f(x),f(y))$. Sledi, da je
\[d_H(f(x),f(y))=d_H(x,y)
\qquad\text{in}\qquad
d_T(f(x),f(y))=d_T(x,y),\]
torej je $f$ izometrija glede na Hilbertovo in Thompsonovo metriko.

S podobnim argumentom kot v prejšnji trditvi dobimo naslednji
rezultat.

\begin{posledica}
Naj bosta $K\subseteq V$ in $K'\subseteq V'$ zaprta stožca. Če je
$f:K\longrightarrow K'$ preslikava, ki obrača red in je homogena
stopnje $r<0$, potem za vse $x,y\in K$, za katere velja $x\sim_K y$,
velja
\[
d_H\bigl(f(x),f(y)\bigr)\leq |r|\,d_H(x,y)
\qquad\text{in}\qquad
d_T\bigl(f(x),f(y)\bigr)\leq |r|\,d_T(x,y).
\]
\end{posledica}

\begin{dokaz}
Naj bo $x\sim_K y$. Ker je $K$ zaprt, velja
$m(x/y)y\leq_K x\leq_K M(x/y)y$. Ker $f$ obrača red in je homogena
stopnje $r<0$, dobimo
$M(x/y)^r f(y)\leq_{K'}f(x)\leq_{K'}m(x/y)^r f(y)$.
Zato je $m(f(x)/f(y))\geq M(x/y)^r$ in
$M(f(x)/f(y))\leq m(x/y)^r$. Sledi
\[
d_H\bigl(f(x),f(y)\bigr)
\leq
\log\left(\frac{m(x/y)^r}{M(x/y)^r}\right)
=
|r|\,d_H(x,y)
\]
in
\[
d_T\bigl(f(x),f(y)\bigr)
\leq
\log\bigl(\max\{M(x/y)^{-r},M(y/x)^{-r}\}\bigr)
=
|r|\,d_T(x,y).
\]
\end{dokaz}

Posebej je vsaka preslikava, ki obrača red in je homogena stopnje
$-1$, neekspanzivna glede na $d_H$ in $d_T$.

Hilbertova metrika je posebej uporabna pri preučevanju skaliranih
homogenih preslikav, ki ohranjajo red, kot kaže naslednja lema.

\begin{lema}
Naj bo $K\subseteq V$ strnjen zaprt stožec in naj bo
$f:K\longrightarrow K$ homogena preslikava, ki ohranja red, pri čemer
je $f(x)\neq0$ za vse $x\in K\setminus\{0\}$. Naj bo
$\varphi:K\longrightarrow[0,\infty)$ homogena preslikava, za katero
velja $\varphi(x)\neq0$ za vse $x\in K\setminus\{0\}$. Če je
$\Sigma=\{x\in K:\varphi(x)=1\}$ in je preslikava
$g:\Sigma\longrightarrow\Sigma$ podana z
\[
g(x)=\frac{f(x)}{\varphi(f(x))},
\]
potem je $g$ neekspanzivna glede na Hilbertovo metriko na
$\Sigma\cap P$ za vsak del $P$ stožca $K$. Če $f$ poleg tega strogo
ohranja red, potem je $g$ skrčitev.
\end{lema}

\begin{dokaz}
Naj bosta $x,y\in\Sigma$, za katera velja $x\sim_K y$, in označimo
$\alpha=m(x/y)$ ter $\beta=M(x/y)$. Ker je $K$ zaprt, velja
$\alpha y\leq_K x\leq_K\beta y$. Poleg  tega $f$ ohranja red, zato je
$f(\alpha y)\leq_K f(x)\leq_K f(\beta y)$, iz pogoja homogenosti pa sledi
$\alpha f(y)\leq_K f(x)\leq_K\beta f(y)$. Sedaj uporabimo definicijo 
preslikave $g$ 
\[
g(x)=\frac{f(x)}{\varphi(f(x))},
\qquad
g(y)=\frac{f(y)}{\varphi(f(y))}
\]
in s preureditvijo v $f(x)= \varphi(f(x)) g(x)$ in $f(y)= \varphi(f(y)) g(y)$ dobimo
\[
\alpha\frac{\varphi(f(y))}{\varphi(f(x))}g(y)
\leq_K g(x)
\leq_K
\beta\frac{\varphi(f(y))}{\varphi(f(x))}g(y).
\]
Iz definicij $m(g(x)/g(y))$ in $M(g(x)/g(y))$ sledi
\[
m(g(x)/g(y))\geq \alpha\frac{\varphi(f(y))}{\varphi(f(x))},
\qquad
M(g(x)/g(y))\leq \beta\frac{\varphi(f(y))}{\varphi(f(x))}.
\]
Zato je $d_H(g(x),g(y))\leq\log(\beta/\alpha)=d_H(x,y)$.

Če $f$ strogo ohranja red in $x\neq\lambda y$ za vsak $\lambda>0$,
potem iz $\alpha y\leq_K x\leq_K\beta y$ sledi
$\alpha f(y)\ll_K f(x)\ll_K\beta f(y)$.
Zato obstajata $\mu>\alpha$ in $\tau<\beta$, za katera velja
$\mu f(y)\leq_K f(x)\leq_K\tau f(y)$. Posledično
\[
d_H(g(x),g(y))
\leq\log(\tau/\mu)
<
\log(\beta/\alpha)
=
d_H(x,y).
\]
\end{dokaz}

Tipičen primer homogene funkcije
$\varphi:K\longrightarrow[0,\infty)$ iz prejšnje leme je
$\varphi(x)=\|x\|$ za neko normo $\|\cdot\|$ na $V$ in
$\varphi\in\operatorname{int}(K^*)$.

Thompsonova metrika je najbolj uporabna pri preučevanju subhomogenih
preslikav, ki ohranjajo red, na celotnem stožcu.

\begin{lema}
Naj bosta $K\subseteq V$ in $K'\subseteq V'$ strnjena zaprta stožca. Če
je $f:K\longrightarrow K'$ preslikava, ki ohranja red, potem je $f$
subhomogena natanko tedaj, ko je neekspanzivna glede na Thompsonovo
metriko na vsakem delu stožca $K$.

Če je
$f:\operatorname{int}(K)\longrightarrow\operatorname{int}(K')$
strogo subhomogena preslikava, ki ohranja red, potem je skrčitev.
\end{lema}

\begin{dokaz}
Predpostavimo najprej, da je $f$ subhomogena, in naj bosta
$x,y\in K$, za katera velja $x\sim_K y$. Tedaj je
$d_T(x,y)=\log\lambda$ za  $\lambda = \max\{M(x/y),M(y/x)\} \geq 1$. Ker je $K$ zaprt,
velja $y\leq_K M(y/x) \leq_K \lambda x$ in $x\leq_K M(x/y) \leq_K\lambda y$. Ker $f$ ohranja red
in je subhomogena ter $ 0 < \lambda^{-1} < 1 $, dobimo
\[
\lambda^{-1}f(y)
\leq_{K'}f(\lambda^{-1}y)
\leq_{K'}f(x),
\qquad
\lambda^{-1}f(x)
\leq_{K'}f(\lambda^{-1}x)
\leq_{K'}f(y).
\tag{2.5}
\]
Od tod sledi
$\max\{M(f(y)/f(x)),M(f(x)/f(y))\}\leq\lambda$ in zato
$d_T(f(x),f(y))\leq\log\lambda=d_T(x,y)$.

Za obratno implikacijo predpostavimo, da je $f$ neekspanzivna glede
na $d_T$ na vsakem delu stožca $K$. Naj bo $x\in K$ in
$y=\lambda^{-1}x$, kjer je $\lambda\geq1$. Če je $x\neq0$, velja
$d_T(x,y)=\log\lambda$, če pa je $x=0$, velja
$d_T(x,y)=0\leq\log\lambda$. Ker je $f$ neekspanzivna, dobimo
\[
\log M(f(x)/f(y))
\leq d_T(f(x),f(y)) \leq d_T(x,y)
\leq\log\lambda.
\]
Ker je logaritem naraščajoča funkcija, dobimo $M(f(x)/f(y)) \leq \lambda$. 
Ker je $K'$ zaprt, sledi $f(x) \leq \lambda f(y)$. Upoštevamo, da je 
$ y = \lambda^{-1} x$ ter obe strani delimo z $\lambda$ in dobimo
$\lambda^{-1}f(x)\leq_{K'}f(\lambda^{-1}x)$, kar pomeni, da je $f$
subhomogena.

Naj bo sedaj
$f:\operatorname{int}(K)\longrightarrow\operatorname{int}(K')$
strogo subhomogena preslikava, ki ohranja red, in naj bo
$\lambda = \max{M(x/y),M(y/x)} >1$. Podobno kot v~\textup{(2.5)} dobimo
\[
\lambda^{-1}f(y)
\ll_{K'}f(\lambda^{-1}y)
\leq_{K'}f(x),
\qquad
\lambda^{-1}f(x)
\ll_{K'}f(\lambda^{-1}x)
\leq_{K'}f(y)
\]
za vse različne $x,y\in\operatorname{int}(K)$. 
Iz $f(y)\ll_{K'}\lambda f(x)$ po definiciji strogega reda sledi, da je $z:=\lambda f(x)-f(y)\in\operatorname{int}(K')$. Definirajmo preslikavo $h\colon\mathbb{R}\to V'$ s predpisom $h(t)=z-tf(x)$. Preslikava $h$ je zvezna in velja $h(0)=z\in\operatorname{int}(K')$. Ker je $\operatorname{int}(K')$ odprta, je množica $h^{-1}(\operatorname{int}(K'))$ odprta v $\mathbb{R}$ in vsebuje $0$. Zato obstaja $\delta>0$, da za vsak $t\in[0,\delta)$ velja $h(t)=z-tf(x)\in\operatorname{int}(K')$.
Izberimo $0<\varepsilon<\min\{\delta,\lambda\}$ in postavimo $\mu_1=\lambda-\varepsilon$. Tedaj je $0<\mu_1<\lambda$ in velja $\mu_1f(x)-f(y)=(\lambda-\varepsilon)f(x)-f(y)=z-\varepsilon f(x)=h(\varepsilon)\in\operatorname{int}(K')\subseteq K'$. Zato je $f(y)\leq_{K'}\mu_1f(x)$, od koder po definiciji $M$ sledi $M(f(y)/f(x))\leq\mu_1<\lambda$.
 Zato velja $\max\{M(f(x)/f(y)),M(f(y)/f(x))\}<\lambda$ in posledično $d_T(f(x),f(y))<\log\lambda=d_T(x,y)$.
\end{dokaz}

Pri računanju $M(x/y)$ in $m(x/y)$ je včasih uporabna naslednja
ugotovitev. Naj bo $K\subseteq V$ strnjen zaprt stožec in
$u\in\operatorname{int}(K)$. Definirajmo
\[
\Sigma_u^*=\{\varphi\in K^*:\varphi(u)=1\}.
\]
Po lemi o dualnem stožcu~\ref{lema2.46} je $\Sigma_u^*$ kompaktna konveksna množica.
Poleg tega je $\varphi(u)>0$ za vse
$\varphi\in K^*\setminus\{0\}$, ker je $u\in\operatorname{int}(K)$.

Če je na primer $K=\mathbb{R}_+^n$ in $u=(1,1, \ldots, 1)$, potem je
\[
\Sigma_u^*
=
\left\{
\varphi\in\mathbb{R}_+^n:
\sum_i\varphi_i=1
\right\}.
\]

\begin{definicija}
Točka $z$ konveksne množice $S\subseteq V$ je \emph{ekstremna točka},
če ne obstajata različni točki $x,y\in S\setminus\{z\}$ in število
$0<\lambda<1$, za katere bi veljalo
$z=\lambda x+(1-\lambda)y$.
\end{definicija}

Znano je, da je vsaka kompaktna konveksna množica konveksna ogrinjača
svojih ekstremnih točk; glej~\cite{rockafellar}. (TO DO)

\begin{trditev}\label{konv-ogr}
Vsaka kompaktna konveksna množica v končnorazsežnem realnem vektorskem
prostoru je konveksna ogrinjača svojih ekstremnih točk.
\end{trditev}

\begin{posledica}\label{posledica:ekstremne_tocke_nepoliedrski}
Naj bo $K\subseteq V$ strnjen zaprt stožec,
$\varphi\in\operatorname{int}(K^*)$ in
\begin{align*}
\Sigma_\varphi=\{x\in K:\varphi(x)=1\}.
\end{align*}
Če ima $\Sigma_\varphi$ neskončno mnogo ekstremnih točk, potem
$K$ ni poliedrski stožec.
\end{posledica}

\begin{dokaz}
Pokažimo, da dve različni ekstremni točki množice $\Sigma_\varphi$
ne moreta pripadati istemu delu stožca $K$. Naj bosta
$x,y\in\Sigma_\varphi$ ekstremni točki in predpostavimo, da
$x\sim_K y$. Tedaj obstaja $\beta>0$, tako da $x\leq_K\beta y$.
Če postavimo $\alpha=1/\beta$, dobimo $\alpha x\leq_K y$, zato je
$y-\alpha x\in K$.
Ker je $\varphi$ pozitiven funkcional in
$\varphi(x)=\varphi(y)=1$, velja
$0\leq\varphi(y-\alpha x)=1-\alpha$,
torej $\alpha\leq1$.

Če je $\alpha<1$, postavimo
\begin{align*}
z=\frac{y-\alpha x}{1-\alpha}.
\end{align*}
Tedaj je $z\in\Sigma_\varphi$ in
$y=\alpha x+(1-\alpha)z$.
Ker je $y$ ekstremna točka, sledi $x=y$.
Če je $\alpha=1$, velja $y-x\in K$ in
$\varphi(y-x)=0$. Ker je $\varphi\in\operatorname{int}(K^*)$,
sledi $y-x=0$, torej ponovno $x=y$.

Različne ekstremne točke množice $\Sigma_\varphi$ zato pripadajo
različnim delom stožca $K$. Če ima $\Sigma_\varphi$ neskončno mnogo
ekstremnih točk, ima $K$ neskončno mnogo delov. Po točki~(iv)
leme~\ref{lema2.38} pa ima poliedrski stožec le končno mnogo delov.
Zato $K$ ni poliedrski.
\end{dokaz}

Iz leme~\ref{lema2.31} in linearnosti $\varphi$ sledi, da je $x\leq_K\beta y$ natanko tedaj,
ko je $\varphi(x)\leq\beta\varphi(y)$ za vse
$\varphi\in\Sigma_u^*$. Ker je $\Sigma_u^*$ konveksna ogrinjača svojih
ekstremnih točk, je ta pogoj ekvivalenten pogoju
$\varphi(x)\leq\beta\varphi(y)$ za vse $\varphi\in \varepsilon_u^*$, 
kjer je $\varepsilon_u^*$ množica ekstremnih točk množice $\Sigma_u^*$.
Zato
\begin{align}
  M(x/y)
  =
  \inf\{\beta>0:x\leq_K\beta y\}
  =
  \sup_{\varphi\in \varepsilon_u^*}
  \frac{\varphi(x)}{\varphi(y)}
  \label{oznaka2.6}
\end{align}
in
\begin{align}
  m(x/y)
  =
  \sup\{\alpha>0:\alpha y\leq_K x\}
  =
  \inf_{\varphi\in \varepsilon_u^*}
  \frac{\varphi(x)}{\varphi(y)}.
  \label{oznaka2.7}
\end{align}

V formulah~\eqref{oznaka2.6} in~\eqref{oznaka2.7} lahko množico $\varepsilon_u^*$
nadomestimo s poljubno podmnožico $S^*$ množice $\Sigma_u^*$, ki
vsebuje $\varepsilon_u^*$.






\subsection{Lorentzovi stožci}

Za $n\geq 2$ je eden izmed pomembnih primerov nepoliedrskih stožcev
Lorentzov stožec. Spomnimo se, da smo ga definirali kot
\begin{align*}
\Lambda_{n+1}
=
\{(x_1,\ldots,x_{n+1})\in\R^{n+1} ~|~
x_1^2-x_2^2-\cdots-x_{n+1}^2\geq0,\ x_1\geq0\}.
\end{align*}
Če vektor $x\in\R^{n+1}$ zapišemo kot $x=(s,u)$, kjer je
$s\in\R$ in $u\in\R^n$, lahko Lorentzov stožec pregledneje zapišemo kot
\begin{align*}
\Lambda_{n+1}
=
\{(s,u)\in\R\times\R^n ~|~ s\geq\|u\|_2\}.
\end{align*}

Definirajmo kvadratno formo
$Q:\R^{n+1}\longrightarrow\R$ s predpisom
\begin{align*}
Q((s,u))
=
s^2-\|u\|_2^2.
\end{align*}
Lorentzov stožec lahko tako zapišemo kot
\begin{align*}
\Lambda_{n+1}
=
\{(s,u)\in\R\times\R^n ~|~
Q((s,u))\geq0 \text{ in } s\geq0\}.
\end{align*}

Naj bo
\begin{align*}
H^n
=
\{(s,u)\in\R\times\R^n ~|~
Q((s,u))=1 \text{ in } s>0\}.
\end{align*}
Ker za $(s,u)\in H^n$ velja
$s=\sqrt{1+\|u\|_2^2}$,
je $H^n$ zgornji del dvodelnega hiperboloida
$\{x\in\R^{n+1} ~|~ Q(x)=1\}$
in leži v $\operatorname{int}(\Lambda_{n+1})$.

Definirajmo še simetrično bilinearno formo
$B:\R^{n+1}\times\R^{n+1}\longrightarrow\R$ s predpisom
\begin{align*}
B(x,y)
=
st-\langle u,v\rangle,
\end{align*}
kjer sta $x=(s,u)$ in $y=(t,v)$. Med $Q$ in $B$ velja
\begin{align*}
Q(x)=B(x,x).
\end{align*}

Za nadaljnje izračune Hilbertove metrike na
$\operatorname{int}(\Lambda_{n+1})$ definirajmo za
$x,y\in\operatorname{int}(\Lambda_{n+1})$ kvadratni polinom
\begin{align*}
p(\lambda)
=
\lambda^2-2B(x,y)\lambda+Q(x)Q(y).
\end{align*}
Njegovi ničli sta
\begin{align*}
\alpha(x,y)
&=
B(x,y)
-
\sqrt{B(x,y)^2-Q(x)Q(y)},\\
\beta(x,y)
&=
B(x,y)
+
\sqrt{B(x,y)^2-Q(x)Q(y)}.
\end{align*}
Funkciji $\alpha$ in $\beta$ sta pozitivno homogeni v obeh
spremenljivkah, saj velja 
$B(\lambda x,\lambda y)=  \lambda s \lambda t - \langle \lambda u, \lambda v \rangle = \lambda^2 B(x,y)$ 
in $Q(\lambda x)=\lambda^2 Q(x)$ za vse $\lambda>0$ in $x,y\in\R^{n+1}$. 
V dokazu naslednjega izreka bomo videli tudi, da za
$x,y\in H^n$ velja
$B(x,y)^2-Q(x)Q(y)\geq0$.

\begin{izrek}\label{izr:lorentz_M_m_H}
Za $x,y\in H^n$ velja
\begin{align*}
M(x/y;\Lambda_{n+1})
=
\beta(x,y)
\qquad\text{in}\qquad
m(x/y;\Lambda_{n+1})
=
\alpha(x,y).
\end{align*}
\end{izrek}

\begin{dokaz}
Če je $x=y$, sta enakosti očitni, saj sta obe strani enaki $1$. Predpostavimo zato, da sta
$x=(s,u)$ in $y=(t,v)$ različni točki množice $H^n$.

Velja $x\leq_{\Lambda_{n+1}}\lambda y$ natanko tedaj, ko je
$\lambda y-x\in\Lambda_{n+1}$, kar je ekvivalentno pogojema
\begin{align*}
Q(\lambda y-x)\geq0
\qquad\text{in}\qquad
\lambda t-s\geq0.
\end{align*}
Ker je $Q(x)=Q(y)=1$, dobimo
\begin{align*}
Q(\lambda y-x)
&=
Q((\lambda t-s,\lambda v-u))=\\
&=
(\lambda t-s)^2
-
\langle\lambda v-u,\lambda v-u\rangle=\\
&=
\lambda^2(t^2-\|v\|_2^2)
-2\lambda(st-\langle u,v\rangle)
+(s^2-\|u\|_2^2)=\\
&=
\lambda^2-2\lambda B(x,y)+1=\\
&=
p(\lambda).
\end{align*}
Zato je
\begin{align}
M(x/y;\Lambda_{n+1})
=
\inf\limits
\left\{
\lambda>0 ~|~
p(\lambda)\geq0
\text{ in }
\lambda\geq\frac{s}{t}
\right\}.
\label{eq:lorentz_M}
\end{align}

Po Cauchy-Schwarzovi neenakosti in neenakosti
$2ab\leq a^2+b^2$ dobimo
\begin{align*}
2stB(x,y)
&=
2s^2t^2-2st\langle u,v\rangle \geq
2s^2t^2-2st\|u\|_2\|v\|_2 \geq \\
&\geq
2s^2t^2-t^2\|u\|_2^2-s^2\|v\|_2^2=\\
&=
2s^2t^2-t^2(s^2-1)-s^2(t^2-1)=\\
&=
s^2+t^2.
\end{align*}
V drugi neenakosti smo uporabili
$2ab\leq a^2+b^2$ za $a=t\|u\|_2$ in $b=s\|v\|_2$, v predzadnji enakosti pa dejstvo, da je $Q(x)=Q(y)=1$, 
torej je $1 = s^2 - \|u\|_2^2 = t^2 - \|v\|_2^2$.

Poglejmo še, kdaj v zgornji neenakosti nastopi enakost.
Če sta $u=v=0$, iz $Q(x)=Q(y)=1$ in $s,t>0$ sledi
$s=t=1$, torej je $x=y$. V nasprotnem primeru enakost po
Cauchy--Schwarzovi neenakosti zahteva
$u=\mu v$ za neki $\mu>0$, enakost v
$2ab\leq a^2+b^2$ pa zahteva $t\|u\|_2=s\|v\|_2$.
Od tod sledi $\mu=s/t$ in zato
\begin{align*}
x
=
(s,u)
=
\left(s,\frac{s}{t}v\right).
\end{align*}
Ker je $Q(x)=Q(y)=1$, dobimo
\begin{align*}
1
=
Q(x)
&=
s^2-\left(\frac{s}{t}\right)^2\|v\|_2^2 =
\frac{s^2}{t^2}
\left(t^2-\|v\|_2^2\right) =\\
&=
\frac{s^2}{t^2}Q(y)\ =
\frac{s^2}{t^2}.
\end{align*}
Ker sta $s,t>0$, sledi $s=t$, zato je tudi $\mu=1$ in $x=y$.
Enakost torej nastopi natanko tedaj, ko je $x=y$.

Poleg tega iz zgornje neenakosti sledi
\begin{align*}
B(x,y)
\geq
\frac{s^2+t^2}{2st}
\geq
1
=
\sqrt{Q(x)Q(y)}.
\end{align*}
Zato je
\begin{align*}
B(x,y)^2-Q(x)Q(y)\geq0,
\end{align*}
torej sta $\alpha(x,y)$ in $\beta(x,y)$ realni ničli polinoma $p$.

Ker smo predpostavili $x\neq y$, je
\begin{align*}
2stB(x,y)>s^2+t^2,
\end{align*}
zato je
\begin{align*}
s^2-2stB(x,y)+t^2<0.
\end{align*}
Posledično je
\begin{align*}
t^2p(s/t)
=
s^2-2stB(x,y)+t^2
<0,
\end{align*}
torej je $p(s/t)<0$.

Ker sta $\alpha(x,y)$ in $\beta(x,y)$ ničli polinoma $p$ in je
njegov vodilni koeficient pozitiven, je $p(\lambda)<0$ natanko med
ničlama. Zato
\begin{align*}
\frac{s}{t}
\in
(\alpha(x,y),\beta(x,y)).
\end{align*}
Iz~\eqref{eq:lorentz_M} tako sledi
$M(x/y;\Lambda_{n+1})=\beta(x,y)$.

Podobno je
\begin{align*}
m(x/y;\Lambda_{n+1})
&=
\sup\limits
\{\lambda>0 ~|~
\lambda y\leq_{\Lambda_{n+1}}x\} =\\
&=
\sup\limits
\left\{
\lambda>0 ~|~
p(\lambda)\geq0
\text{ in }
\lambda\leq\frac{s}{t}
\right\}=\\
&=
\alpha(x,y).
\end{align*}
\end{dokaz}

Prejšnji izrek lahko s pomočjo homogenosti razširimo na celotno
notranjost Lorentzovega stožca.

\begin{posledica}
\label{pos:lorentz_M_m}
Za vse $x,y\in\operatorname{int}(\Lambda_{n+1})$ velja
\begin{align*}
M(x/y;\Lambda_{n+1})
&=
\frac{\beta(x,y)}{Q(y)},\\
m(x/y;\Lambda_{n+1})
&=
\frac{\alpha(x,y)}{Q(y)}.
\end{align*}
Poleg tega velja
\begin{align*}
d_H(x,y)
&=
\log\left(
\frac{\beta(x,y)}{\alpha(x,y)}
\right),\\
d_T(x,y)
&=
\log\left(
\max\limits
\left\{
\frac{\beta(x,y)}{Q(x)},
\frac{\beta(x,y)}{Q(y)}
\right\}
\right).
\end{align*}
\end{posledica}

\begin{dokaz}
Postavimo
\begin{align*}
\sigma=\sqrt{Q(x)}
\qquad\text{in}\qquad
\tau=\sqrt{Q(y)}.
\end{align*}
Ker sta $x,y\in\operatorname{int}(\Lambda_{n+1})$, velja
$Q(x)>0$ in $Q(y)>0$, zato sta $\sigma,\tau>0$. Poleg tega velja 
$Q(\sigma^{-1}x)
=\sigma^{-2}Q(x)=1$
in podobno $Q(\tau^{-1}y)=1$.
Ker množimo s pozitivnima skalarjema, ostaneta
$\sigma^{-1}x$ in $\tau^{-1}y$ v notranjosti Lorentzovega stožca.
Zato velja $\sigma^{-1}x,\tau^{-1}y\in H^n$.

Po izreku~\ref{izr:lorentz_M_m_H} in pozitivni homogenosti funkcije
$\beta$ dobimo
\begin{align*}
M(x/y;\Lambda_{n+1})
&=
\frac{\sigma}{\tau}
M(\sigma^{-1}x/\tau^{-1}y;\Lambda_{n+1})=\\
&=
\frac{\sigma}{\tau}
\beta(\sigma^{-1}x,\tau^{-1}y)=
\frac{\sigma}{\tau}\frac{1}{\tau \sigma}
\beta(x,y)=
\frac{1}{\tau^2}\beta(x,y)=\\
&=
\frac{\beta(x,y)}{Q(y)}.
\end{align*}
Podobno dobimo
\begin{align*}
m(x/y;\Lambda_{n+1})
&=
\frac{\sigma}{\tau}
m(\sigma^{-1}x/\tau^{-1}y;\Lambda_{n+1})=\\
&=
\frac{\sigma}{\tau}
\alpha(\sigma^{-1}x,\tau^{-1}y)=
\frac{\sigma}{\tau}\frac{1}{\tau \sigma}
\alpha(x,\tau y)=
\frac{1}{\tau^2}\alpha(x,y)=\\
&=
\frac{\alpha(x,y)}{Q(y)}.
\end{align*}

Formula za Hilbertovo metriko sledi neposredno iz njene definicije.
Ker je $\beta(x,y)=\beta(y,x)$, za Thompsonovo metriko dobimo
\begin{align*}
d_T(x,y)
&=
\log\left(
\max\limits\{M(x/y),M(y/x)\}
\right)=\\
&=
\log\left(
\max\limits
\left\{
\frac{\beta(x,y)}{Q(y)},
\frac{\beta(x,y)}{Q(x)}
\right\}
\right).
\end{align*}
\end{dokaz}

Na koncu omenimo še povezavo Lorentzovega stožca s hiperbolično
geometrijo. 
Poglejmo, kako na množici $H^n$ definiramo hiperbolično metriko.
Naj bo $x\in H^n$. Naj bo $x=(s,u)\in H^n$ in naj bo
$z=(t,v)\neq0$ tangentni vektor na $H^n$ v točki $x$.

Ker je $H^n$ del nivojske množice $Q(x)=1$, je tangentni vektor $z$
pravokoten na gradient funkcije $Q$. Ker je $\nabla Q(x)=(2s,-2u)$,
velja
$0 = \langle\nabla Q(x),z\rangle =
2st-2\langle u,v\rangle$,
torej
\begin{align}
st-\langle u,v\rangle=0.
\label{eq:tangentni_lorentz}
\end{align}

Pokažimo, da za vsak tak neničeln tangentni vektor velja
$B(z,z)<0$. Če je $u=0$, iz $Q(x)=1$ in $s>0$ sledi $s=1$.
Iz~\eqref{eq:tangentni_lorentz} nato dobimo $t=0$. Ker je $z\neq0$,
mora biti $v\neq0$, zato
\begin{align*}
B(z,z)
=
t^2-\|v\|_2^2
=
-\|v\|_2^2
<0.
\end{align*}

Predpostavimo sedaj, da je $u\neq0$. Iz~\eqref{eq:tangentni_lorentz}
in Cauchy-Schwarzove neenakosti sledi
\begin{align*}
s|t|
=
|\langle u,v\rangle|
\leq
\|u\|_2\|v\|_2.
\end{align*}
Ker je $x\in H^n$, velja
\begin{align*}
s^2-\|u\|_2^2=1,
\end{align*}
zato je $s>\|u\|_2$. Posledično je 
\begin{align*}
|t|
\leq
\frac{\|u\|_2}{s}\|v\|_2
<
\|v\|_2.
\end{align*}
Od tod dobimo
\begin{align*}
B(z,z)
=
t^2-\|v\|_2^2
<0.
\end{align*}

Tako za vsak $x\in H^n$ in vsak neničeln
$z\in T_x H^n$ velja $-B(z,z)>0$. Forma $-B$ je torej na
tangentnem prostoru $T_x H^n$ pozitivno določena, zato jo lahko
uporabimo za definicijo Riemannove metrike na $H^n$.

\begin{definicija}
Naj bo $\gamma:[a,b]\longrightarrow H^n$ odsekoma $C^1$ pot.
Njeno dolžino definiramo kot
\begin{align*}
L(\gamma)
=
\int_a^b
\sqrt{-B(\gamma'(t),\gamma'(t))}\,dt.
\end{align*}
Za $x,y\in H^n$ definiramo \emph{hiperbolično razdaljo}
\begin{align*}
d_{\mathrm{hyp}}(x,y)
=
\inf\limits_{\gamma}L(\gamma),
\end{align*}
kjer infimum vzamemo po vseh odsekoma $C^1$ poteh $\gamma$ v $H^n$,
ki povezujejo $x$ in $y$.

Metrični prostor $(H^n,d_{\mathrm{hyp}})$ imenujemo
\emph{hiperboloidni model} $n$-razsežnega hiperboličnega prostora.
\end{definicija}

Za hiperbolično razdaljo med točkama $x,y\in H^n$ velja
\begin{align}
\cosh d_{\mathrm{hyp}}(x,y)
=
B(x,y).
\label{eq:hiperbolicna_razdalja}
\end{align}

Pokažimo sedaj povezavo med hiperbolično, Hilbertovo in Thompsonovo
metriko. Za $x,y\in H^n$ velja $Q(x)=Q(y)=1$. Ker sta
$\alpha(x,y)$ in $\beta(x,y)$ ničli polinoma
$p(\lambda) =
\lambda^2-2B(x,y)\lambda+1$,
po Vietovih formulah velja
$\alpha(x,y)\beta(x,y)=1$.
Zato je
\begin{align*}
\alpha(x,y)=\beta(x,y)^{-1}.
\end{align*}
Po izreku~\ref{izr:lorentz_M_m_H} dobimo
\begin{align*}
\frac{1}{2}d_H(x,y)
&=
\frac{1}{2}
\log\left(
\frac{\beta(x,y)}{\alpha(x,y)}
\right)\\
&=
\frac{1}{2}
\log\bigl(\beta(x,y)^2\bigr)\\
&=
\log\beta(x,y).
\end{align*}

Ker je $B$ simetrična bilinearna forma in je $Q(x)=Q(y)=1$, velja
$\beta(x,y)=\beta(y,x)$.
Zato po izreku~\ref{izr:lorentz_M_m_H}
$M(x/y;\Lambda_{n+1})
= M(y/x;\Lambda_{n+1})
= \beta(x,y)$,
in posledično
\begin{align*}
d_T(x,y)
=
\log\beta(x,y).
\end{align*}
Torej je 
\begin{align}
d_T(x,y)
=
\frac{1}{2}d_H(x,y).
\label{eq:lorentz_H_T}
\end{align}

Po Vietovih formulah velja tudi
$\alpha(x,y)+\beta(x,y)
= 2B(x,y)$.
Ker je $\alpha(x,y)=\beta(x,y)^{-1}$, dobimo
\begin{align*}
B(x,y)
&=
\frac{\beta(x,y)+\beta(x,y)^{-1}}{2}=
\cosh\bigl(\log\beta(x,y)\bigr).
\end{align*}
Skupaj z~\eqref{eq:hiperbolicna_razdalja} sledi
$\cosh d_{\mathrm{hyp}}(x,y) =
\cosh\bigl(\log\beta(x,y)\bigr)$.
Ker je $\beta(x,y)\geq1$, sta
$d_{\mathrm{hyp}}(x,y)$ in $\log\beta(x,y)$ nenegativna. Funkcija
$\cosh$ je strogo naraščajoča na $[0,\infty)$, zato
$d_{\mathrm{hyp}}(x,y)
= \log\beta(x,y)$.
Za vse $x,y\in H^n$ tako velja
\begin{align*}
d_{\mathrm{hyp}}(x,y)
=
d_T(x,y)
=
\frac{1}{2}d_H(x,y).
\end{align*}





\subsection{Stožec pozitivno semidefinitnih simetričnih matrik}

V tem podpoglavju obravnavamo vektorski prostor
$\operatorname{Sym}_n$ simetričnih matrik velikosti $n\times n$ in
strnjen zaprt stožec $\Pi_n(\R)$, ki ga sestavljajo vse pozitivno
semidefinitne matrike iz $\operatorname{Sym}_n$.

Pri analizi Hilbertove in Thompsonove metrike na $\Pi_n(\R)$ je
koristno opaziti, da grupa obrnljivih matrik
$GL_n(\R)$ deluje tranzitivno na
$\operatorname{int}(\Pi_n(\R))$ s konjugacijo. To pomeni, da za poljubni
matriki $A,B\in\operatorname{int}(\Pi_n(\R))$ obstaja
$Q\in GL_n(\R)$, za katero velja
\begin{align*}
A=Q^T B Q.
\end{align*}

Na prostoru $\operatorname{Sym}_n$ definiramo skalarni produkt
\begin{align*}
\langle A,B\rangle
=
\operatorname{tr}(AB).
\end{align*}
Glede na ta skalarni produkt je stožec $\Pi_n(\R)$ sebi dualen, torej 
$\Pi_n(\R) = \Pi_n(\R)^* =
\left\{B\in\operatorname{Sym}_n ~|~ 
\langle A,B\rangle\geq0
\text{ za vsak } A\in\Pi_n(\R)
\right\}$.
Res, če je $B\in\Pi_n(\R)$, potem je tudi
$B^{1/2}\in\Pi_n(\R)$, saj so tudi njene lastne vrednosti nenegativne. Zato za vsak $A\in\Pi_n(\R)$ velja
$\langle A,B\rangle =
\operatorname{tr}(AB)=
\operatorname{tr}(B^{1/2}AB^{1/2})\geq0$.
Od tod sledi $B\in\Pi_n(\R)^*$.

Obratno, naj bo $B\in\operatorname{Sym}_n$ in predpostavimo, da
$\operatorname{tr}(AB)\geq0$
za vsak $A\in\Pi_n(\R)$. Po spektralnem izreku obstaja ortogonalna
matrika $U$, za katero je $B =UDU^T$,
kjer je $D=\operatorname{diag}(\lambda_1,\ldots,\lambda_n)$
in $\lambda_1\leq\lambda_2\leq\cdots\leq\lambda_n$
so lastne vrednosti matrike $B$. Tedaj za vsak
$A\in\Pi_n(\R)$ velja
\begin{align*}
\operatorname{tr}(AB)=
\operatorname{tr}(AU DU^T)=\operatorname{tr}(U^T A U D)\geq0.
\end{align*}
Če vzamemo $A=U E_i U^T$, kjer je $E_i$ matrika z $1$ na $i$-tem mestu in ničelami drugod, 
sledi, da mora biti $\lambda_i\geq0$
za vsak $i=1,\ldots,n$. Zato je $B\in\Pi_n(\R)$ in posledično $\Pi_n(\R)^* = \Pi_n(\R)$.

V nadaljevanju naj $I$ označuje identično matriko, ki pripada
$\operatorname{int}(\Pi_n(\R))$, in naj bo
\begin{align*}
\Sigma_I^*
=
\{A\in\Pi_n(\R) ~|~ \operatorname{tr}(A)=1\}.
\end{align*}

\begin{lema}\label{lema:ekstremne_tocke_psd}
Matrika $A\in\Sigma_I^*$ je ekstremna točka množice $\Sigma_I^*$
natanko tedaj, ko ima natanko eno neničelno lastno vrednost, ki je
enaka $1$.
\end{lema}

\begin{dokaz}
Naj $E_I^*$ označuje množico ekstremnih točk množice $\Sigma_I^*$,
$\mathcal{D}$ pa množico vseh matrik iz $\Sigma_I^*$, ki imajo natanko
eno neničelno lastno vrednost, ta pa je enaka $1$.

Naj bo najprej $A\in E_I^*$. Po spektralnem izreku obstaja ortogonalna
matrika $U$, za katero velja $A=U^T D U$,
kjer je $D=\operatorname{diag}(\lambda_1,\ldots,\lambda_n)$,
$0\leq\lambda_1\leq\ldots\leq\lambda_n$.
Ker je $A\in\Sigma_I^*$, velja $\sum\limits_{i=1}^n\lambda_i=1$.
Če z $e_i$ označimo $i$-ti standardni bazni vektor, lahko matriko $A$ zapišemo kot
\begin{align*}
A
=
U^T D U
=
\sum\limits_{i=1}^n
\lambda_i
U^T\operatorname{diag}(e_i)U.
\end{align*}
Vsaka matrika $U^T\operatorname{diag}(e_i)U$
pripada množici $\Sigma_I^*$. 

Predpostavimo, da sta vsaj dva izmed koeficientov $\lambda_i$
pozitivna. Izberimo indeks $j$, za katerega je $\lambda_j>0$.
Ker obstaja še vsaj en drug pozitiven koeficient in velja
$\sum\limits_{i=1}^n\lambda_i=1$, dobimo $0<\lambda_j<1$.
Označimo $A_i=U^T\operatorname{diag}(e_i)U$
in definirajmo matriko $C=\frac{1}{1-\lambda_j}
\sum\limits_{\substack{i=1\\ i\neq j}}^n
\lambda_i A_i$.
Ker velja
$\sum\limits_{\substack{i=1\\ i\neq j}}^n
\frac{\lambda_i}{1-\lambda_j}
=\frac{1-\lambda_j}{1-\lambda_j}=1$,
je $C$ konveksna kombinacija elementov množice $\Sigma_I^*$, zato
$C\in\Sigma_I^*$.
Sedaj lahko $A$ zapišemo kot
$A=\lambda_j A_j
+
\sum\limits_{\substack{i=1\\ i\neq j}}^n
\lambda_i A_i
=
\lambda_j A_j
+
(1-\lambda_j)
\left(
\frac{1}{1-\lambda_j}
\sum\limits_{\substack{i=1\\ i\neq j}}^n
\lambda_i A_i
\right)
=
\lambda_j A_j+(1-\lambda_j)C$.

Ker je $A$ ekstremna točka in je zgornji
zapis konveksna kombinacija elementov množice $\Sigma_I^*$ ter smo lastne vrednosti $\lambda_i$ 
uredili po velikosti, mora veljati
$\lambda_1=\ldots=\lambda_{n-1}=0$,
$\lambda_n=1$.
Zato je
$E_I^*\subseteq\mathcal{D}$.

Za obratno inkluzijo opazimo, da grupa ortogonalnih matrik $O(n)$
deluje tranzitivno na $\mathcal{D}$ s konjugacijo, saj imata poljubni matriki iz $\mathcal D$ enake lastne 
vrednosti $0,\ldots,0,1$, zato ju lahko po spektralnem izreku z ustrezno ortogonalno spremembo baze 
preslikamo eno v drugo. Poleg tega za vsak fiksen $U\in O(n)$ preslikava
$A
\longmapsto
U^T A U$
linearno preslika množico $\Sigma_I^*$ nase, zato ekstremne točke
preslika v ekstremne točke. Ker je $E_I^*\subseteq\mathcal{D}$
neprazna množica in $O(n)$ deluje tranzitivno na $\mathcal{D}$, sledi
$E_I^*
=
\mathcal{D}$.
\end{dokaz}

Iz leme~\ref{lema:ekstremne_tocke_psd} sledi, da ima
$\Sigma_I^*$ za $n\geq2$ neskončno mnogo ekstremnih točk, zato po
posledici~\ref{posledica:ekstremne_tocke_nepoliedrski}
stožec $\Pi_n(\R)$ ni poliedrski.

Po drugi strani $\Pi_n(\R)$ za $n\geq3$ ni strogo konveksen. Res, za
$0\leq\lambda\leq1$ velja
\begin{align*}
(1-\lambda)
\operatorname{diag}(1,0,0,\ldots,0)
+
\lambda
\operatorname{diag}(0,1,0,\ldots,0)
\in
\partial\Pi_n(\R).
\end{align*}
Tako daljica med dvema različnima točkama roba stožca ostane na
njegovem robu.

Stožec $\Pi_2(\R)$ pa je strogo konveksen in je linearno izomorfen
Lorentzovemu stožcu $\Lambda_3$. To vidimo s pomočjo bijektivne
linearne preslikave
\begin{align*}
\begin{pmatrix}
a & b\\
b & c
\end{pmatrix}
\longmapsto
\left(
\frac{a+c}{2},
\frac{a-c}{2},
b
\right),
\end{align*}
ki stožec $\Pi_2(\R)$ preslika na $\Lambda_3$.

Sedaj izračunajmo funkciji $M$ in $m$ za matriki
$X,Y\in\operatorname{int}(\Pi_n(\R))$. Pogoj
\begin{align*}
X\leq_{\Pi_n(\R)}\beta Y
\end{align*}
velja natanko tedaj, ko
\begin{align*}
Y^{-1/2}XY^{-1/2}
\leq_{\Pi_n(\R)}
\beta I.
\end{align*}
To je ekvivalentno pogoju
\begin{align*}
\sigma(Y^{-1}X)
=
\sigma(Y^{-1/2}XY^{-1/2})
\subseteq
(0,\beta],
\end{align*}
kjer $\sigma(A)$ označuje spekter matrike $A$. Zato je
\begin{align*}
M(X/Y;\Pi_n(\R))
=
\max\limits
\{\lambda:\lambda\in\sigma(Y^{-1}X)\}.
\end{align*}

Podobno pogoj
\begin{align*}
\alpha Y\leq_{\Pi_n(\R)}X
\end{align*}
velja natanko tedaj, ko
\begin{align*}
\alpha I
\leq_{\Pi_n(\R)}
Y^{-1/2}XY^{-1/2},
\end{align*}
kar je ekvivalentno pogoju
\begin{align*}
\sigma(Y^{-1}X)
=
\sigma(Y^{-1/2}XY^{-1/2})
\subseteq
[\alpha,\infty).
\end{align*}
Zato je
\begin{align*}
m(X/Y;\Pi_n(\R))
=
\min\limits
\{\lambda:\lambda\in\sigma(Y^{-1}X)\}.
\end{align*}

Dobljene rezultate lahko povzamemo v naslednji trditvi.

\begin{trditev}\label{trditev:metriki_psd}
Naj bosta $X,Y\in\operatorname{int}(\Pi_n(\R))$ in definirajmo
\begin{align*}
\lambda_+(X,Y)
&=
\max\limits
\{\lambda:\lambda\in\sigma(Y^{-1}X)\},\\
\lambda_-(X,Y)
&=
\min\limits
\{\lambda:\lambda\in\sigma(Y^{-1}X)\}.
\end{align*}
Tedaj velja
\begin{align*}
d_H(X,Y)
&=
\log\left(
\frac{\lambda_+(X,Y)}
{\lambda_-(X,Y)}
\right),\\
d_T(X,Y)
&=
\log\left(
\max\limits
\left\{
\lambda_+(X,Y),
\frac{1}{\lambda_-(X,Y)}
\right\}
\right).
\end{align*}
\end{trditev}

Enake formule lahko s podobnimi argumenti dobimo tudi za elemente
notranjosti stožca pozitivno semidefinitnih hermitskih matrik
$\Pi_n(\C)$.







\subsection{Konveksnost Hilbertovih krogel}

Pri dokazu obstoja lastnega vektorja bomo potrebovali naslednjo lastnost
zaprtih krogel glede na Hilbertovo metriko.

\begin{lema}\label{lema:konveksnost_hilbertovih_krogel}
Naj bo $K\subseteq V$ strnjen zaprt stožec,
$\varphi\in\operatorname{int}(K^*)$ in
$\Sigma=\{x\in K:\varphi(x)=1\}$. Potem je vsaka zaprta krogla
$B_r(x)=\{y\in\Sigma:d_H(x,y)\leq r\}$ konveksna podmnožica prostora
$V$.
\end{lema}

\begin{dokaz}
Naj bosta $y_1,y_2\in B_r(x)$. 
Za vsak $t \in [0,1]$ želimo dokazati, da je $y_t = (1-t)y_1 + t y_2$ element množice $B_r(x)$.
Označimo
$\alpha_i=m(y_i/x)$ in $\beta_i=M(y_i/x)$ za $i=1,2$. Ker je $K$
zaprt, sta infimum in supremum v definicijah količin $M$ in $m$
dosežena, zato velja
\begin{equation*}
\alpha_1x\leq_K y_1\leq_K\beta_1x
\qquad\text{in}\qquad
\alpha_2x\leq_K y_2\leq_K\beta_2x.
\end{equation*}
Ker sta $y_1,y_2\in B_r(x)$, 
sledi $d_H (x,y_i) \leq r$, torej velja
\begin{equation*}
\frac{\beta_i}{\alpha_i}
=
\frac{M(y_i/x)}{m(y_i/x)}
=
e^{d_H(y_i,x)}
\leq e^r
\end{equation*}
za $i=1,2$.

Naj bo $0\leq t\leq1$ in definirajmo
$y_t=(1-t)y_1+ty_2$,
$\alpha_t=(1-t)\alpha_1+t\alpha_2$ ter
$\beta_t=(1-t)\beta_1+t\beta_2$. Iz zgornjih neenakosti sledi
$\alpha_t x
\leq_K
y_t
\leq_K
\beta_t x$.
Poleg tega je
\begin{align*}
\frac{\beta_t}{\alpha_t}
&=
\frac{(1-t)\alpha_1}{\alpha_t}
\frac{\beta_1}{\alpha_1}
+
\frac{t\alpha_2}{\alpha_t}
\frac{\beta_2}{\alpha_2}\\
&\leq
\frac{(1-t)\alpha_1}{\alpha_t}e^r
+
\frac{t\alpha_2}{\alpha_t}e^r
=
e^r.
\end{align*}
Po definicijah količin $M$ in $m$ zato velja
$M(y_t/x)\leq\beta_t$ in $m(y_t/x)\geq\alpha_t$. 
Zato je $\frac{M(y_t/x)}{m(y_t/x)} \leq \frac{\beta_t}{\alpha_t}$.
Posledično je
\begin{equation*}
d_H(x,y_t)
=
\log\left(\frac{M(y_t/x)}{m(y_t/x)}\right)
\leq
\log\left(\frac{\beta_t}{\alpha_t}\right)
\leq r.
\end{equation*}

Ker je $\varphi$ linearna in sta $y_1,y_2\in\Sigma$, velja
$\varphi(y_t)=(1-t)\varphi(y_1)+t\varphi(y_2)=1$, zato je
$y_t\in\Sigma$. Torej je $y_t\in B_r(x)$ za vsak $0\leq t\leq1$,
kar pomeni, da je $B_r(x)$ konveksna.
\end{dokaz}







\section{Lastni vektorji in lastne vrednosti nelinearnih preslikav na stožcih}

V tem poglavju bomo postavili temelje spektralne teorije zveznih homogenih
preslikav $f\colon K\longrightarrow K$, ki ohranjajo red, kjer je $K$
strnjen zaprt stožec v končnorazsežnem vektorskem prostoru $V$.
Predstavili bomo dve smiselni definiciji stožčnega spektralnega radija 
$r_K(f)$ in pokazali, da sovpadata. Glavni rezultat, ki posplošuje
Krein-Rutmanov izrek~\ref{Krein-Rutman}, pove, da obstaja neničeln lastni vektor $x\in K$,
ki pripada stžčnemu spektralnemu radiju $r_K(f)$, torej
\begin{equation*}
f(x)=r_K(f)x.
\end{equation*}

Pokazali bomo tudi, da za stožčni spektralni radij $r_K(f)$ velja
minimaksna variacijska formula, podobna Collatz-Wielandtovi formuli za
spektralni radij nenegativnih matrik. Obravnavali bomo še zveznost
preslikave $r_K$ in analizirali 
\begin{definicija}
  \emph{Stožčni spekter} zvezne homogene
  preslikave $f\colon K\longrightarrow K$, ki ohranja red, kjer je $K$
  strnjen zaprt stožec v končnorazsežnem vektorskem prostoru $V$, je
  \begin{equation*}
  \sigma_K(f)
  =
  \{\lambda\geq 0:f(x)=\lambda x
  \text{ za neki }x\in K\setminus\{0\}\}.
  \end{equation*}
\end{definicija}


Preden začnemo razvijati teorijo, moramo razjasniti praktično vprašanje,
povezano z zveznimi razširitvami preslikav, ki ohranjajo red.

\subsection{Razširitve preslikav, ki ohranjajo red}

Pogosto obravnavamo subhomogeno preslikavo, ki ohranja red in je
definirana le na notranjosti zaprtega stožca. V takih primerih nas zanima, ali ima preslikava zvezno subhomogeno
razširitev na rob stožca, ki prav tako ohranja red. V tem podpoglavju
bomo obravnavali ta problem. Izkazalo se bo, da imajo pri njegovi
rešitvi pomembno vlogo poliedrski stožci.

\begin{definicija}
  Naj bo $K\subseteq V$ strnjen zaprt stožec in $x\in\partial K$.
  Pravimo, da zaporedje $(x_k)_k$ v $\operatorname{int}(K)$
  \emph{konvergira proti $x$ od zgoraj}, če $x_k\longrightarrow x$ in
  $x\ll_K x_k$ za vsak $k\geq1$. 
\end{definicija}

Taka zaporedja vedno obstajajo. Poglejmo si naslednji preprost primer.
\begin{zgled}\label{zgled4.3}
  Naj bo $u\in\operatorname{int}(K)$ in definirajmo zaporedje $x_k=x+\frac{u}{k}$.
  Pokazati želimo, da $(x_k)_k$ proti $x$ od zgoraj.
  Ker je $x_k-x=\frac{u}{k}$ in velja $u\in\operatorname{int}(K)$ ter je $\frac{1}{k}>0$, 
  velja tudi $\frac{u}{k}\in\operatorname{int}(K)$, saj je notranjost stožca zaprta za množenje s 
  pozitivnimi skalarji. Zato je $x_k-x\in\operatorname{int}(K)$, kar po definiciji pomeni $x\ll_K x_k$.
  Poleg tega velja $\|x_k-x\|=\left\|\frac{u}{k}\right\|=\frac{\|u\|}{k}\longrightarrow0$, zato je 
  $x_k\longrightarrow x$.
  Opazimo še, da je vsak $x_k$ v notranjosti stožca, torej velja $x_k=x+\frac{u}{k}\in\operatorname{int}(K)$.
\end{zgled}

\begin{lema}\label{lema:podzaporedje_od_zgoraj}
Naj bo $K\subseteq V$ strnjen zaprt stožec, $x\in\partial K$ in naj
zaporedje $(x_k)_k$ konvergira proti $x$ od zgoraj. Potem obstaja
podzaporedje $(x_{k_i})_i$, za katero velja
\begin{equation*}
x_m\ll_K x_{k_i}
\qquad\text{za vse }m\geq k_{i+1}.
\end{equation*}
\end{lema}

\begin{dokaz}
Podzaporedje $(x_{k_i})_i$ definiramo induktivno. Postavimo
$x_{k_1}=x_1$ in predpostavimo, da smo že izbrali indekse
$k_1<k_2<\cdots<k_i$, ki zadoščajo trditvi. Ker je
$x_{k_i}-x\in\operatorname{int}(K)$, obstaja okolica $U_i$ točke
$x_{k_i}-x$, za katero velja
$U_i\subseteq\operatorname{int}(K)$.

Naj bo $V_i$ odprta simetrična okolica točke $0$, izbrana tako, da
\begin{equation*}
x_{k_i}-x-v\in U_i
\qquad\text{za vsak }v\in V_i.
\end{equation*}
Ker $x_k\longrightarrow x$, obstaja $k_{i+1}>k_i$, tako da je
$x_m-x\in V_i$ za vsak $m\geq k_{i+1}$. Zato je
\begin{equation*}
x_{k_i}-x_m
=
x_{k_i}-x-(x_m-x)
\in U_i
\subseteq\operatorname{int}(K)
\end{equation*}
za vse $m\geq k_{i+1}$. Torej je $x_m\ll_K x_{k_i}$.
\end{dokaz}

Z uporabo prejšnje leme dobimo naslednji izrek o razširitvi.

\begin{izrek}\label{izr:razsiritev_preslikave_na_stozec}
Naj bosta $K\subseteq V$ in $K'\subseteq V'$ zaprta stožca, pri
čemer je $K$ strnjen, ter naj bo
$f\colon\operatorname{int}(K)\longrightarrow K'$ zvezna preslikava,
ki ohranja red. Potem veljajo naslednje trditve:
\begin{enumerate}[label= (\roman*)]
\item Če zaporedje $(x_k)_k$ proti $x\in\partial K$ od zgoraj,
potem obstaja $y_x\in K'$, za katerega velja
\begin{equation*}
\lim_{k\to\infty}f(x_k)=y_x.
\end{equation*}

\item Preslikava $F\colon K\longrightarrow K'$, ki je na
$\operatorname{int}(K)$ enaka $f$, za $x\in\partial K$ pa je
definirana s predpisom $F(x)=y_x$, je dobro definirana razširitev
preslikave $f$, ki ohranja red.

\item Če je $f$ poleg tega subhomogena oziroma homogena, je tudi
$F\colon K\longrightarrow K'$ subhomogena oziroma homogena.
\end{enumerate}
\end{izrek}

\begin{dokaz}
(i) Po lemi~\ref{lema:podzaporedje_od_zgoraj} obstaja podzaporedje
$(x_{k_i})_i$, za katero velja $x_m\ll_K x_{k_i}$ za vse
$m\geq k_{i+1}$. Posebej je
$0\leq_K x_{k_{i+1}}\ll_K x_{k_i}$
za vsak $i\geq1$. Ker $f$ ohranja red, dobimo
$0\leq_{K'} f(x_{k_{i+1}}) \leq_{K'} f(x_{k_{i}})$.
Ker je $K'$ zaprt stožec v končnorazsežnem prostoru, je normalen.
Zaporedje $f(x_{k_{i}})$ je padajoče zaporedje, zato obstaja $y_x\in K'$, za katerega velja
$\lim\limits_{i\to\infty} f(x_{k_i})=y_x$.

Pokažimo, da k $y_x$ konvergira celotno zaporedje $(f(x_k))_k$.
Naj bo $\varepsilon>0$ in izberimo $i\geq1$ tako, da je $\|f(x_{k_i})-y_x\|<\varepsilon$.
Naj bo $m\geq k_{i+1}$. Tedaj je $x_m\ll_K x_{k_i}$. Ker tudi
$(x_{k_i})_i$ konvergira proti $x$ od zgoraj, obstaja $j\geq i+1$,
za katerega je $x_{k_j}\leq_K x_m$.
V nasprotnem primeru bi za vse $j\geq i+1$ veljalo
$x_m\leq_K x_{k_j}$. Z limito in zaprtostjo stožca $K$
bi dobili $x_m\leq_K x$, kar je v protislovju z $x\ll_K x_m$.
Za vse dovolj velike $j$ torej velja $x_{k_j}\leq_K x_m\leq_K x_{k_i}$.
Ker $f$ ohranja red, sledi $f(x_{k_j})\leq_{K'}f(x_m)\leq_{K'}f(x_{k_i})$.
Iz normalnosti stožca $K'$ dobimo 
$\|f(x_m)-f(x_{k_j})\|
\leq
M\|f(x_{k_i})-f(x_{k_j})\|$.
Ko pošljemo $j\to\infty$, dobimo
$\|f(x_m)-y_x\|
\leq
M\|f(x_{k_i})-y_x\|
<
M\varepsilon$, kar pomeni $(f(x_k))_k$ konvergira proti $y_x$.

(ii) Preverimo, da je $F$ dobro definirana. Naj bosta $(x_k)_k$ in
$(\widetilde{x}_k)_k$ zaporedji, ki konvergirata k isti točki
$x\in\partial K$ od zgoraj. Definirajmo zaporedje $(z_k)_k$ s
predpisom
\begin{equation*}
z_{2k}=x_k,
\qquad
z_{2k-1}=\widetilde{x}_k.
\end{equation*}
Tudi $(z_k)_k$ konvergira proti $x$ od zgoraj. Po točki (i) zato
velja
\begin{equation*}
\lim_{k\to\infty}f(x_k)
=
\lim_{k\to\infty}f(z_k)
=
\lim_{k\to\infty}f(\widetilde{x}_k).
\end{equation*}
Vrednost $F(x)$ je torej neodvisna od izbire zaporedja, zato je $F$
dobro definirana razširitev preslikave $f$.

Pokažimo še, da $F$ ohranja red. Naj bosta $x,y\in K$ taka, da velja 
$x \leq_K y$, in naj bo $u\in\operatorname{int}(K)$. Po zgledu~\ref{zgled4.3} zaporedji
$(x+\frac{u}{k})_k$ in $(y+\frac{u}{k})_k$ konvergirata proti $x$ oziroma $y$ od
zgoraj. Za vsak $k\geq1$ velja
$x+\frac{u}{k}
\leq_K
y+\frac{u}{k}$,
zato je
$f\left(x+\frac{u}{k}\right)
\leq_{K'}
f\left(y+\frac{u}{k}\right)$.
Iz tega dobimo 
$\lim\limits_{k\to\infty} f\left(x+\frac{u}{k}\right) \leq_{K'} \lim\limits_{k\to\infty} f\left(y+\frac{u}{k}\right)$, 
iz česar pa sledi $F(x)\leq_{K'}F(y)$.

(iii) Predpostavimo nazadnje, da je $f$ subhomogena. Naj bo $x\in K$ in
$0<\lambda<1$. Če $(x_k)_k$ proti $x$ od zgoraj, potem
$(\lambda x_k)_k$ konvergira proti $\lambda x$ od zgoraj. Zato je
\begin{equation*}
\lambda F(x)
=
\lim_{k\to\infty}\lambda f(x_k)
\leq_{K'}
\lim_{k\to\infty}f(\lambda x_k)
=
F(\lambda x),
\end{equation*}
torej je $F$ subhomogena. Če je $f$ homogena, z podobnim argumentom
dobimo $F(\lambda x)=\lambda F(x)$ za vsak $\lambda>0$.
\end{dokaz}

Razširitev iz izreka~\ref{izr:razsiritev_preslikave_na_stozec} na
splošno ni nujno zvezna. To pokaže naslednji zgled.

\begin{zgled}\label{zgled:nezvezna_razsiritev}
Naj bo $\operatorname{Sym}_n (\R)$ realni vektorski prostor simetričnih
matrik velikosti $n\times n$, $\Pi_n(\R)$ pa stožec pozitivno
semidefinitnih matrik v $\operatorname{Sym}_n$. Postavimo
\begin{equation*}
K=\Pi_n(\R)\times\Pi_n(\R).
\end{equation*}
Vemo že, da je stožec
$\Pi_n(\R)$ strnjen in zaprt v $\operatorname{Sym}_n$. Ker je stožec $K$ produkt zaprtih množic, je zaprt v
$\operatorname{Sym}_n\times\operatorname{Sym}_n$.
Definirajmo
\begin{equation*}
f\colon\operatorname{int}(K)\longrightarrow\Pi_n(\R),
\qquad
f(A,B)=(A^{-1}+B^{-1})^{-1}.
\end{equation*}
Na pozitivno definitnih matrikah preslikava $A\mapsto A^{-1}$ očitno obrača
red. Zato iz $A_1\leq A_2$ in $B_1\leq B_2$ sledi
$A_2^{-1}+B_2^{-1}
\leq
A_1^{-1}+B_1^{-1}$,
po ponovni uporabi obračanja reda pa
$f(A_1,B_1)\leq f(A_2,B_2)$. Preslikava $f$ torej ohranja red.
Poleg tega je tudi očitno homogena stopnje $r = 1$. Zato je po posledici~\ref{posledica3.9} 
neekspanzivna glede na Thompsonovo metriko in posledično Lishitzeva s konstanto $1$ 
in zato zvezna na $\operatorname{int}(K)$.

Pokažimo sedaj, da razširitev $F$ iz
izreka~\ref{izr:razsiritev_preslikave_na_stozec} ni nujno zvezna. Poglejmo si primer 
za $n=2$. Vzemimo
\begin{equation*}
A=B=
\begin{bmatrix}
1&0\\
0&0
\end{bmatrix}
\in\partial\Pi_2(\R)
\end{equation*}
in za $k\in\N$ postavimo
\begin{equation*}
A_k=A+\frac{1}{k}I,
\qquad
B_k=B+\frac{1}{k}I.
\end{equation*}
Zaporedje $((A_k,B_k))_k$ konvergira proti $(A,B)$ od zgoraj. 
Ker za vsak $k\in\N$ velja $A_k=B_k =
\begin{bmatrix}
1+\frac{1}{k} & 0\\
0 & \frac{1}{k}
\end{bmatrix}$, dobimo
$F(A,B)
=
\lim\limits_{k\to\infty} f(A_k,B_k)
= \lim\limits_{k\to\infty} (2A_k^{-1})^{-1} 
= \lim\limits_{k\to\infty} 
\frac{1}{2}\begin{bmatrix}
1+\frac{1}{k}&0\\
0&\frac{1}{k}
\end{bmatrix}
=
\begin{bmatrix}
\frac12&0\\
0&0
\end{bmatrix}$.

Po drugi strani za $k,m\in\N$, kjer je $m>1$, definirajmo
\begin{equation*}
\widehat{A}_k=
\begin{bmatrix}
1&\frac1k\\
\frac1k&\frac{m}{k^2}
\end{bmatrix},
\qquad
\widehat{B}_k=
\begin{bmatrix}
1&\frac2k\\
\frac2k&\frac{4m}{k^2}
\end{bmatrix}.
\end{equation*}
Zaporedje $((\widehat{A}_k,\widehat{B}_k))_k$ prav tako konvergira proti
$(A,B)$ od zgoraj, vendar z izračunom lahko preverimo, da velja  
$\lim\limits_{k\to\infty}
f(\widehat{A}_k,\widehat{B}_k)
=
\begin{bmatrix}
a&0\\
0&0
\end{bmatrix}$,
$a=\frac{5m-5}{10m-9}<\frac12$.
Ker se limiti ne ujemata, preslikava $F$ ni zvezna v
točki $(A,B)$.
\end{zgled}

Stožec $\Pi_n(\R)$ ni poliedrski (TO DO: cite). Izkazalo se bo, da je razširitev
$F: K \to K'$ iz izreka~\ref{izr:razsiritev_preslikave_na_stozec} vedno
zvezna, kadar je domena poliedrski stožec. Razlog je naslednja
geometrijska lastnost, ki karakterizira poliedrske stožce.

\begin{definicija}\label{def:pogoj_G}
Naj bo $K\subseteq V$ zaprt stožec in $x\in K$. Pravimo, da $K$
\emph{zadošča pogoju G v točki $x$}, če za vsako zaporedje
$(x_k)_k\subseteq K$, za katero velja $\lim\limits_{k\to \infty} x_n = x$, in za
vsak $0<\lambda<1$ obstaja $m\in\N$ tako, da velja 
$\lambda x\leq_K x_k~\text{za vse }k\geq m$.

Pravimo, da $K$ \emph{zadošča pogoju G}, če pogoju G zadošča v vsaki
točki $x\in K$.
\end{definicija}

Če je $K$ strnjen zaprt stožec, je pogoj G izpolnjen v vsaki točki
$x\in\operatorname{int}(K)$. 
Res, naj bo $(x_k)_k \subseteq K$ zaporedje, za katero velja $\lim\limits_{k\to \infty} x_n = x$.
Pokazati moramo, da za vsak $0<\lambda<1$ obstaja $m\in\mathbb{N}$, tako da za vse $k\geq m$ velja 
$\lambda x\leq_K x_k$.
Po definiciji urejenosti, ki jo določa stožec $K$, je to enakovredno pogoju $x_k-\lambda x\in K$.
Razliko lahko zapišemo kot $x_k-\lambda x=(x_k-x)+(1-\lambda)x$.
Ker je $x\in\operatorname{int}(K)$ in je $1-\lambda>0$, je tudi $(1-\lambda)x\in\operatorname{int}(K)$, 
torej po definiciji notranjosti obstaja $U \subseteq K$, da je $x-x_n \in U$. 
Ker je $U\subseteq K$, za vse $k\geq m$ velja $x_k-\lambda x\in K$. Po definiciji urejenosti od tod sledi 
$\lambda x\leq_K x_k$. Stožec $K$ torej zadošča pogoju $G$ v točki $x$.

Da bi bil pogoj G izpolnjen tudi v vsaki točki
roba $\partial K$, pa mora biti stožec poliedrski.

\begin{lema}\label{lema:pogoj_G_poliedrski}
Strnjen zaprt stožec $K\subseteq V$ je poliedrski natanko tedaj, ko
zadošča pogoju G.
\end{lema}

\begin{dokaz}
Predpostavimo najprej, da je $K$ poliedrski. Potem 
obstajajo linearni funkcionali $\psi_1,\ldots,\psi_N$, za katere je
$K = \{x\in V:\psi_i(x)\geq0
\text{ za }i=1,\ldots,N\}$.
Naj bo $(x_k)_k$ zaporedje v $K$, ki konvergira proti $x\in K$, in naj
bo $0<\lambda<1$. Ker so $\psi_i$ linearni funkcionali, za vsak $i$ velja
$\lim\limits_{k\to\infty}\psi_i(x_k)=\psi_i(x)$.
Če je $\psi_i(x)>0$, potem
$\psi_i(x_k)-\psi_i(\lambda x)
\longrightarrow
(1-\lambda)\psi_i(x)>0$.
Zato obstaja $m_i\geq1$, tako da je
$\psi_i(x_k)\geq\psi_i(\lambda x)$ za vse $k\geq m_i$. Če je
$\psi_i(x)=0$, pa je
$0=\psi_i(\lambda x)\leq\psi_i(x_k)$
za vsak $k\geq1$, zato lahko postavimo $m_i=1$.

Sedaj vzemimo $m=\max\{m_i~|~i=1,\ldots,N\}$.
Tedaj je $\psi_i(\lambda x)\leq\psi_i(x_k)$ za vsak $i$ in vsak
$k\geq m$. Zato je $x_k-\lambda x\in K$, oziroma
$\lambda x\leq_K x_k$, za vse $k\geq m$. Stožec $K$ torej zadošča
pogoju G.

Dokažimo še obratno implikacijo. Naj $K$ zadošča pogoju G. Izberimo
$\varphi\in\operatorname{int}(K^*)$ in definirajmo
$\Sigma_\varphi
=
\{x\in K~|~\varphi(x)=1\}$.
Po Lemi~\ref{lema2.46} je $\Sigma_\varphi$ kompaktna konveksna
množica in $\varphi$ je pozitiven funkcional. Ker je $\Sigma_\varphi$ prerez stožca $K$ s 
hiperploskvijo $\{x\in V ~|~ \varphi(x)=1\}$, velja
$K = \{t x ~|~ t\geq0,\ x\in\Sigma_\varphi\}$.
Stožec $K$ bo torej poliedrski, če pokažemo, da je $\Sigma_\varphi$
poliedrska množica (\textit{Tukaj nekaj ni dobro razloženo - TO DO.}). 
Po trditvi~\ref{konv-ogr} zadostuje pokazati, da ima
$\Sigma_\varphi$ le končno mnogo ekstremnih točk.

Naj bo $x\in\Sigma_\varphi$ in $0<\lambda<1$. Iz pogoja G sledi, da
obstaja okolica $U_x$ točke $x$, za katero velja
$\lambda x\leq_K u$ 
za vsak $u\in U_x\cap K$.
Res, če take okolice ne bi bilo, bi lahko izbrali zaporedje
$(u_k)_k\subseteq K$, ki bi konvergiralo k $x$, vendar za noben $k$
ne bi veljalo $\lambda x\leq_K u_k$, kar bi bilo v protislovju s
pogojem G.

Naj bo $y\neq x$ element množice $U_x\cap\Sigma_\varphi$ in izberimo
$0<\alpha\leq\lambda$. Postavimo
\begin{equation*}
z
=
\frac{y-\alpha x}{1-\alpha}.
\end{equation*}
Ker je $\alpha x\leq_K\lambda x\leq_K y$, je $z\in K$. Poleg tega velja 
$\varphi(z)
=
\frac{\varphi(y)-\alpha\varphi(x)}{1-\alpha}
=
1$,
zato je $z\in\Sigma_\varphi$. Velja
$y=(1-\alpha)z+\alpha x$.
Ker je $y\neq x$, sta tudi $z$ in $x$ različna, zato $y$ ni
ekstremna točka množice $\Sigma_\varphi$. Torej vsaka okolica $U_x$
vsebuje največ eno ekstremno točko množice $\Sigma_\varphi$ in to je $x$.

Ker je $\Sigma_\varphi$ kompaktna, ima odprto pokritje
$\{U_x ~|~ x\in\Sigma_\varphi\}$ končno podpokritje. Zato ima
$\Sigma_\varphi$ le končno mnogo ekstremnih točk. Ker je kompaktna
in konveksna, je konveksna ogrinjača svojih ekstremnih točk, zato je
poliedrska množica. Posledično je tudi stožec $K$ poliedrski.
\end{dokaz}

Z uporabo pogoja G lahko sedaj dokažemo naslednji rezultat.

\begin{izrek}\label{izr:zvezna_razsiritev_poliedrski_stozec}
Naj bosta $K\subseteq V$ in $K'\subseteq V'$ zaprta stožca, pri
čemer je $K$ strnjen poliedrski stožec. Če je
$f\colon\operatorname{int}(K)\longrightarrow K'$ subhomogena oziroma
homogena preslikava, ki ohranja red, potem ima $f$ zvezno
subhomogeno oziroma homogeno razširitev
$F\colon K\longrightarrow K'$, ki ohranja red, definirano v izreku~\ref{izr:razsiritev_preslikave_na_stozec}.
\end{izrek}

\begin{dokaz}
Po Izreku~\ref{izr:razsiritev_preslikave_na_stozec} moramo pokazati
le še, da je $F$ zvezna. Ker je $f$ subhomogena in ohranja red, je po posledici~\ref{posledica3.9} 
neekspanzivna glede na Thompsonovo metriko na
$\operatorname{int}(K)$, zato je tam zvezna. Enako velja v
homogenem primeru. Ker se $F$ na $\operatorname{int}(K)$ ujema s
$f$, je $F$ zvezna na notranjosti stožca.

Naj bo sedaj $(x_m)_m\subseteq K$ zaporedje, ki konvergira proti
$x\in\partial K$. Izberimo naraščajoče zaporedje
$(\lambda_k)_k\subseteq(0,1)$, za katero je
$\lambda_k\longrightarrow1$. Ker poliedrski stožec $K$ zadošča
pogoju G, za vsak $k$ obstaja $m_k\geq1$, tako da
$\lambda_k x\leq_K x_m$ za vse $m\geq m_k$.
Po drugi strani, obstaja $u\in\operatorname{int}(K)$ in padajoče
zaporedje pozitivnih števil $(\delta_k)_k$, za katero je
$\delta_k\longrightarrow0$ in
\begin{equation*}
x-\delta_k u
\leq_K
x_m
\leq_K
x+\delta_k u
\qquad\text{za vse }m\geq m_k.
\end{equation*}
Ker $F$ ohranja red in je subhomogena, za vse $m\geq m_k$ dobimo
\begin{equation*}
\lambda_k F(x)
\leq_{K'}
F(\lambda_k x)
\leq_{K'}
F(x_m)
\leq_{K'}
F(x+\delta_k u)
=
f(x+\delta_k u).
\end{equation*}
Pri tem je v homogenem primeru prva neenakost enakost. Zato velja
\begin{equation*}
0
\leq_{K'}
F(x_m)-\lambda_k F(x)
\leq_{K'}
f(x+\delta_k u)-\lambda_k F(x)
\end{equation*}
za vse $m\geq m_k$.

Zaporedje $(x+\delta_k u)_k$ konvergira proti $x$ od zgoraj, zato po
definiciji razširitve velja
$f(x+\delta_k u)\longrightarrow F(x)$.
Ker tudi $\lambda_k F(x)\longrightarrow F(x)$, desna stran prejšnje
neenakosti konvergira k $0$. Stožec $K'$ je normalen, zato sledi
$F(x_m)\longrightarrow F(x)$. Preslikava $F$ je torej zvezna tudi
na robu stožca.
\end{dokaz}




\subsection{Stožčni spekter}

Analiza lastnih vrednosti in lastnih vektorjev homogenih preslikav, ki
na zaprtih stožcih v končnorazsežnih vektorskih prostorih ohranjajo red,
ima številne podobnosti s klasično Perron-Frobeniusovo teorijo, vendar
se od nje v nekaterih pomembnih pogledih tudi razlikuje. Obstajajo na
primer zvezne homogene preslikave $f\colon K\longrightarrow K$, ki
ohranjajo red in imajo neskončno mnogo različnih lastnih vrednosti. To
je lahko presenetljivo, zlasti ker je stožec $K$ končnorazsežen. V tem
podpoglavju bomo to vprašanje podrobneje obravnavali.

\begin{definicija}
  Naj bo $f\colon K\longrightarrow K$ homogena preslikava, ki ohranja red.
  Število $\lambda\geq0$ imenujemo \emph{lastna vrednost} preslikave $f$,
  če obstaja $x\in K\setminus\{0\}$, za katerega je $f(x)=\lambda x$.
  
  \emph{Stožčni spekter} preslikave $f$ označimo z
  \[
  \sigma_K(f)=\{\lambda\geq0:\lambda\text{ je lastna vrednost preslikave }f\}.
  \]
\end{definicija}

Naslednji rezultat je elementaren, vendar uporaben.

\begin{lema}\label{lema:primerjava_lastnih_vrednosti}
Naj bo $f\colon K\longrightarrow K$ homogena preslikava, ki ohranja red,
na zaprtem stožcu $K\subseteq V$. Naj bosta $x,y\in K\setminus\{0\}$ in
naj $y$ dominira $x$. Če velja $\lambda x\leq_K f(x)$ in
$f(y)\leq_K\mu y$, potem je $\lambda\leq\mu$.
\end{lema}

\begin{dokaz}
Primer $\lambda=0$ je trivialen, zato predpostavimo, da je $\lambda>0$.
Ker $y$ dominira $x$, obstaja $\beta>0$, tako da je $x\leq_K\beta y$.
Na predpostavki $\lambda x\leq_K f(x)$ uporabimo $f$, upoštevamo, da ohranja red in je homogena, in dobimo
$\lambda f(x)\leq_K f^2(x)$. Ponovno uporabimo predpostavko $\lambda x\leq_K f(x)$ in dobimo
$\lambda^2 x \leq_K \lambda f(x) \leq_K f^2(x)$.
Podobno lahko naredimo za predpostavko $f(y)\leq_K\mu y$.
Od tod za vsak $k\geq1$ sledi
\[
\lambda^k x
\leq_K f^k(x)
\leq_K \beta f^k(y)
\leq_K \beta\mu^k y,
\]
kjer upoštevamo, da $y$ dominira $x$. Delimo s $\beta\lambda^k$ in dobimo
\[
\left(\frac{\mu}{\lambda}\right)^k y-\frac{1}{\beta}x\in K
\qquad\text{za vsak }k\geq1.
\]
Poglejmo, kaj se zgodi, če pošljemo $k\to\infty$. Če bi veljalo $\mu<\lambda$, bi dobili $-x/\beta\in K$, 
kar ni mogoče, saj je $x\in K$, za stožec pa velja $K\cap(-K)=\{0\}$. Zato je $\lambda\leq\mu$.
\end{dokaz}

Kot posledico dobimo naslednji rezultat.

\begin{posledica}\label{posledica:isti_del_ista_lastna_vrednost}
Naj bo $f\colon K\longrightarrow K$ homogena preslikava, ki ohranja red,
na zaprtem stožcu $K\subseteq V$. Če sta $x$ in $y$ lastna vektorja
preslikave $f$ in velja $x\sim_K y$, potem pripadata isti lastni
vrednosti.
\end{posledica}

\begin{dokaz}
Naj bo $f(x)=\lambda x$ in $f(y)=\mu y$. Ker je $x\sim_K y$, $x$
dominira $y$ in $y$ dominira $x$. Iz
leme~\ref{lema:primerjava_lastnih_vrednosti} zato sledi $\lambda=\mu$.
\end{dokaz}

Ker je $\{0\}$ del vsakega stožca $K$, iz posledice~\ref{posledica:isti_del_ista_lastna_vrednost} sledi, 
da je število različnih lastnih
vrednosti homogene preslikave $f\colon K\longrightarrow K$, ki ohranja
red, največ $m-1$, kjer je $m$ število delov stožca $K$. Za $K=\R_+^n$ lahko preprosto skonstruiramo primer, 
ki pokaže, da je ta zgornja meja dosežena.

\begin{zgled}
  Naj bo $K=\R_+^n$. Za $I\subseteq\{1,\ldots,n\}$ definirajmo $\mathds{1}^I\in\R_+^n$ s
  predpisom
  $ (\mathds{1}^I)_i=1$ za $i\in I$ in $(\mathds{1}^I)_i=0$ sicer.
  Za vsako neprazno množico $I\subseteq\{1,\ldots,n\}$ izberimo še
  $\lambda_I>0$ tako, da iz $I\subseteq J$ in $I\neq J$ sledi
  $\lambda_I<\lambda_J$. Definirajmo $f\colon\R_+^n\longrightarrow\R_+^n$ z
  \[
  f(x)=
  \bigvee_{\emptyset\neq I\subseteq\{1,\ldots,n\}}
  \lambda_I\left(\min_{i\in I}x_i\right)\mathds{1}^I,
  \qquad x\in\R_+^n.
  \]
  Preslikava $f$ je očitno zvezna, homogena in ohranja red. Če je
  $\emptyset\neq J\subseteq\{1,\ldots,n\}$, potem velja
  \[
  f(\mathds{1}^J)
  =
  \bigvee_{\emptyset\neq I\subseteq J}\lambda_I\mathds{1}^I
  =
  \lambda_J\mathds{1}^J,
  \]
  saj je $\lambda_I<\lambda_J$ za vsak $I\subseteq J$ z $I\neq J$.
  Preslikavo $f$ lahko torej izberemo tako, da ima $2^n-1$ različnih
  lastnih vrednosti.

  Poglejmo si na primeru $n=2$, torej $K=\R_+^2$. 
  Stožec $\mathbb{R}_+^2$ ima štiri dele:
  $\{0\}, \operatorname{int}(\mathbb{R}_+^2), \{(x_1,0):x_1>0\}$ in $\{(0,x_2):x_2>0\}$.
  Zato je največje možno število različnih lastnih vrednosti enako
  $4-1=3$.
  Neprazne podmnožice množice $\{1,2\}$ so
  $\{1\}$, $\{2\}$ in $\{1,2\}$. Pripadajoči vektorji so:
  \[
  \mathds{1}^{\{1\}}=(1,0), \qquad
  \mathds{1}^{\{2\}}=(0,1), \qquad
  \mathds{1}^{\{1,2\}}=(1,1).
  \]
  Izberimo
  $\lambda_{\{1\}}=1$, $\lambda_{\{2\}}=2$ in
  $\lambda_{\{1,2\}}=3$ ter definirajmo
  $f\colon\mathbb{R}_+^2\longrightarrow\mathbb{R}_+^2$ s predpisom
  \begin{align*}
  f(x_1,x_2)
  &=\bigvee_{\emptyset\neq I\subseteq\{1,2\}}
  \lambda_I\left(\min_{i\in I}x_i\right)\mathds{1}^I =\\
  &=
  x_1(1,0)
  \vee
  2x_2(0,1)
  \vee
  3\min\{x_1,x_2\}(1,1) =\\
  &=
  \bigl(
  \max\{x_1,3\min\{x_1,x_2\}\},
  \max\{2x_2,3\min\{x_1,x_2\}\}
  \bigr).
  \end{align*}
  Tedaj velja $f(1,0)=(1,0), f(0,1)=2(0,1), f(1,1)=3(1,1)$.
  Preslikava $f$ ima torej tri različne lastne vrednosti $1$, $2$ in $3$,
  zato je zgornja meja dosežena.
\end{zgled}

Izkazalo se bo, da je zgornja meja $m-1$ dosežena za vsak poliedrski
stožec. Pred dokazom uvedimo nekaj oznak. 

\begin{definicija}
  Naj bo $K\subseteq V$ strnjen
  poliedrski stožec z $N$ ličnicami in naj bodo
  $\psi_1,\ldots,\psi_N$ funkcionali, ki določajo njegove ličnice. Za
  $r<0$ in neprazno množico $I\subseteq\{1,\ldots,N\}$ definirajmo
  $M_r(I)\colon\operatorname{int}(K)\longrightarrow[0,\infty)$ s predpisom
  \[
  M_r(I)(x)
  =
  \left(\sum_{i\in I}\psi_i(x)^r\right)^{1/r},
  \qquad x\in\operatorname{int}(K).
  \]
\end{definicija}

Preslikava $M_r(I)$ je na $\operatorname{int}(K)$ očitno zvezna in homogena, saj so 
$\psi_i$ linearni funkcionali, ki pa so na $\operatorname{int}(K)$ pozitivni.
Preverimo še, da ohranja red. Naj bosta $x,y\in\operatorname{int}(K)$ in naj velja $x\leq_K y$. Potem je $\psi_i(x)\leq\psi_i(y)$ za vsak $i\in I$, zato je
\[
M_r(I)(x)
=
\left(\sum_{i\in I}\psi_i(x)^r\right)^{ 1/r}
\leq
\left(\sum_{i\in I}\psi_i(y)^r\right)^{1/r}
=
M_r(I)(y).
\]
Po izreku~\ref{izr:razsiritev_preslikave_na_stozec} ima torej $M_r(I)$ zvezno
homogeno razširitev na $\partial K$, ki ohranja red. 
Poglejmo na robu: Naj bo $x\in\partial K$ in predpostavimo, da za neki $i\in I$ velja $\psi_i(x)=0$. 
Vzemimo zaporedje $(x_k)_k\subseteq\operatorname{int}(K)$, ki konvergira proti $x$.
Tedaj velja $\psi_i(x_k)\longrightarrow\psi_i(x)=0$, pri čemer je $\psi_i(x_k)>0$ za vsak $k\in\N$. Za 
$(x_k)_k$ je vsota definirana z $\left(\sum_{i\in I}\psi_j(x)^r\right)^{1/r}$. Ker so vsi členi nenegativni, je 
$\sum_{j\in I}\psi_j(x_k)^r \geq \psi_i(x_k)^r$. Pri potenciranju z $\frac{1}{r}$ pa se neenakost obrne, saj je 
$r<0$ in dobimo $0 \leq M_r(I)(x_k) \leq \psi_i(x_k)$, torej je $M_r(I)(x)=0$. 

Za $r=-\infty$ definirajmo
$M_{-\infty}(I)\colon K\longrightarrow[0,\infty)$ s predpisom
\[
M_{-\infty}(I)(x)=\min_{i\in I}\psi_i(x),
\qquad x\in K.
\]

\begin{izrek}\label{izrek:stevilo_lastnih_vrednosti_poliedrski}
Naj bo $K\subseteq V$ strnjen poliedrski stožec z $m$ obrazi. Če je
$f\colon K\longrightarrow K$ homogena preslikava, ki ohranja red, potem
velja $|\sigma_K(f)|\leq m-1$. Poleg tega na $K$ obstaja zvezna homogena
preslikava, ki ohranja red in ima $m-1$ različnih lastnih vrednosti.
\end{izrek}

\begin{dokaz}
Iz leme~\ref{lema:deli_stozca} sledi, da ima poliedrski stožec z $m$
obrazi natanko $m$ delov. Če izpustimo trivialni del $\{0\}$ in uporabimo
posledico~\ref{posledica:isti_del_ista_lastna_vrednost}, dobimo, da je
$m-1$ zgornja meja za $|\sigma_K(f)|$.

Skonstruirajmo sedaj primer z $m-1$ različnimi lastnimi vrednostmi. Naj bodo 
$\psi_1,\ldots,\psi_N$ funkcionali, ki določajo ličnice
stožca $K$, množico vseh delov stožca $K$ pa označimo z
$\mathcal{P}(K)$. Tedaj je $|\mathcal{P}(K)|=m$. Spomnimo se, da je vsak
del stožca $K$ relativna notranjost nekega njegovega obraza po lemi~\ref{lema:deli_stozca}. Poleg tega
po točki (ii) leme~\ref{lema2.38} za vsak del stožca $P\in\mathcal{P}(K)$ velja
\[
P
=
\{x\in K ~|~ \psi_i(x)>0\text{ natanko tedaj, ko }i\in I(P)\},
\]
kjer je
$I(P)=\{i ~|~ \psi_i(x)>0\text{ za neki }x\in P\}$.
Za vsak $P\in\mathcal{P}(K)$, $P\neq\{0\}$, izberimo $z^P\in P$.
Skonstruirali bomo zvezno homogeno preslikavo, ki ohranja red in za
katero so točke $z^P$ lastni vektorji z med seboj različnimi lastnimi
vrednostmi.

Naj bo $r\in[-\infty,0)$. Uporabili bomo zvezno homogeno preslikavo
$f_r\colon K\longrightarrow K$, ki ohranja red in je oblike
\begin{equation}\label{eq:preslikava_fr}
f_r(x)
=
\sum_{\substack{P\in\mathcal{P}(K)\\P\neq\{0\}}}
\lambda_P M_r(I(P))(x)u^P,
\qquad x\in K,
\end{equation}
kjer bomo $\lambda_P>0$ in $u^P\in P$ ustrezno izbrali.
Ker je $M_r(I(P))$ zvezna, homogena in ohranja red ter je $\lambda_P>0$, je tudi $f_r$ zvezna, 
homogena in ohranja red. Poleg tega je $u^P\in P \subseteq K$, zato je $f_r(x)\in K$ za vsak $x\in K$.

Spomnimo se, da je množica $\mathcal{P}(K)$ delno urejena z relacijo
$\trianglelefteq$, pri čemer $P\trianglelefteq Q$ pomeni, da za vsak
$x\in P$ in $y\in Q$ obstaja $\beta>0$, tako da je $x\leq_K\beta y$.
Po lemi~\ref{lema2.38} velja tudi
$P\trianglelefteq Q$ natanko tedaj, ko je $I(P)\subseteq I(Q)$.
Zato za vsak $x\in Q$ in vsak $P\in\mathcal{P}(K)$,
$P\neq\{0\}$, velja
$M_r(I(P))(x)>0 \Longleftrightarrow P\trianglelefteq Q$,
sicer pa je $M_r(I(P))(x)=0$.

Števila $\lambda_P>0$ in vektorje $u^P\in P$ bomo izbrali induktivno
glede na višino dela $P$ v delno urejeni množici
$(\mathcal{P}(K),\trianglelefteq)$. Če je $P\in\mathcal{P}(K)$ in
$\dim P=1$, ima $P$ višino $1$. V tem primeru postavimo $u^P=z^P$ in
izberemo $\lambda_P>0$ tako, da so pozitivna števila
\[
\mu_P=\lambda_P M_r(I(P))(u^P)
\]
med seboj različna za vse dele $P$ višine $1$. Za take dele zato velja
$f_r(z^P)=\mu_P z^P$.

Predpostavimo sedaj, da smo $u^P\in P$ in $\lambda_P>0$ že izbrali za
vse $P\trianglelefteq Q$, za katere je $P\neq\{0\}$ in $P\neq Q$.
Tedaj je
\[
\begin{aligned}
f_r(z^Q)
&=
\sum_{\substack{P\in\mathcal{P}(K)\\P\neq\{0\}}}
\lambda_P M_r(I(P))(z^Q)u^P\\
&=
\sum_{\substack{P\trianglelefteq Q\\P\neq\{0\}}}
\lambda_P M_r(I(P))(z^Q)u^P.
\end{aligned}
\]
Označimo
\[
w^Q
=
\sum_{\substack{P\trianglelefteq Q\\P\neq\{0\},\,P\neq Q}}
\lambda_P M_r(I(P))(z^Q)u^P.
\]
Za vsako dovolj veliko število $\mu_Q>0$ velja
$\mu_Q z^Q-w^Q\in Q$ (???). Izberimo $\lambda_Q>0$ in $u^Q\in Q$ tako, da
\[
\lambda_Q M_r(I(Q))(z^Q)u^Q=\mu_Q z^Q-w^Q.
\]
Ker je $M_r(I(Q))(z^Q)>0$, je vektor $u^Q\in Q$ enolično določen, ko
izberemo $\mu_Q>0$ in $\lambda_Q>0$. Iz konstrukcije sledi
\[
f_r(z^Q)=w^Q+\mu_Q z^Q-w^Q=\mu_Q z^Q.
\]
Torej števila $\mu_Q>0$ lahko izberemo tako, da ima $f_r$ natanko
$|\mathcal{P}(K)|-1=m-1$ različnih lastnih vrednosti.
\end{dokaz}

Če je $-\infty<r<0$, je preslikava $f_r$, skonstruirana v
izreku~\ref{izrek:stevilo_lastnih_vrednosti_poliedrski}, neskončnokrat
odvedljiva na $\operatorname{int}(K)$. Vprašamo se lahko, ali obstaja
zvezna homogena preslikava na $K$, ki ohranja red, ima lastnosti iz
izreka~\ref{izrek:stevilo_lastnih_vrednosti_poliedrski}, je razreda
$C^1$ na $\operatorname{int}(K)$ in za katero se $Df(x)$ zvezno razširi
v točko $0$. Če je $\dim K>1$, je odgovor očitno negativen: če bi namreč
$L$ označeval limito $Df(x)$ pri $x\to0$, bi moral imeti $L$ kar $m-1$
različnih lastnih vrednosti, kjer je $m$ število obrazov poliederskega stožca $K$.

Naslednji rezultat pokaže, da stožčni spekter ni nujno končen, kadar
stožec ni poliedrski. Za njegov dokaz potrebujemo Straszewiczev izrek
(glej tudi~\cite[str.~167]{rockafellar}), ki pravi, da so ekstremne
točke zaprte konveksne množice $S$ v končnorazsežnem vektorskem prostoru
$V$ goste v množici ekstremnih točk množice $S$.

\begin{definicija}
  Točka $x$ zaprte konveksne množice $S\subseteq V$ je 
  \emph{ekstremna}, če obstaja hiperravnina $H\subseteq V$ z
  $\dim H=\dim V-1$, za katero velja $H\cap S=\{x\}$.
\end{definicija}

\begin{izrek}[Straszewiczev izrek]\label{Straszewicz}
  Za vsako zaprto konveksno množico $S$ je
  množica ekstremnih točk $S$ gosta podmnožica množice ekstremnih točk
  množice $S$. Tako je vsaka ekstremna točka limita nekega zaporedja ekstremnih
  točk.
\end{izrek}
\begin{dokaz}
  TO DO.
\end{dokaz}

\begin{izrek}\label{izrek:neskoncno_lastnih_vrednosti}
Naj bo $K\subseteq V$ strnjen zaprt stožec. Tedaj obstaja zvezna
homogena preslikava $f\colon K\longrightarrow K$, ki ohranja red in ima
neskončno mnogo različnih lastnih vrednosti, natanko tedaj, ko je $K$
nepoliedrski.
\end{izrek}

\begin{dokaz}
Po izreku~\ref{izrek:stevilo_lastnih_vrednosti_poliedrski} dobimo 
implikacijo iz leve proti desni, torej zadostuje, da
za vsak strnjen zaprt nepoliedrski stožec $K\subseteq V$ skonstruiramo
zvezno homogeno preslikavo $f\colon K\longrightarrow K$, ki ohranja red
in ima neskončno mnogo različnih lastnih vrednosti.

Izberimo $u\in\operatorname{int}(K)$. Če je $K$ nepoliedrski, je tudi
$K^*$ nepoliedrski (glej:~\cite[Trditev~19.2.2, Izrek~14.1]{rockafellar}). Definirajmo
\[
\Sigma^*=\{\varphi\in K^*:\varphi(u)=1\}.
\]
Po lemi~\ref{lema2.46} je $\Sigma^*$ kompaktna konveksna množica v
$V^*$. Ker je $K^*$ nepoliedrski, tudi $\Sigma^*$ ni poliederska, zato ima
neskončno mnogo ekstremnih točk. Po Straszewiczevem izreku so
ekstremne točke množice $\Sigma^*$ goste v njenih ekstremnih
točkah. Zato obstaja zaporedje $(\varphi_k)_k$ med seboj različnih
ekstremnih točk množice $\Sigma^*$, tako da
$\varphi_k\longrightarrow\psi\in\Sigma^*$ in $\varphi_k\neq\psi$ za
vsak $k\geq1$.

Ker je $\varphi_k$ ekstremna točka množice $\Sigma^*$, je
$F=\{\lambda\varphi_k ~|~ \lambda\geq0\}$ izpostavljeni obraz stožca $K^*$.
Ker je $K^{**}=K$, obstaja (?????) $x^k\in K$ z $\|x^k\|=1$, tako da je
$\varphi_k(x^k)=0$ in
$\varphi(x^k)>0$ za vsak
$\varphi\in\Sigma^*\setminus\{\varphi_k\}$. Posebej velja
\begin{equation}\label{eq:pozitivnost_funkcionalov}
\psi(x^k)>0
\qquad\text{in}\qquad
\varphi_m(x^k)>0\quad\text{za vsak }m\neq k.
\end{equation}

Definirajmo $f\colon K\longrightarrow K$ s predpisom
\[
f(x)
=
\sum_{k=1}^{\infty}
\left(\frac12\right)^k
\lambda_k
\inf_{m\neq k}\varphi_m(x)\,x^k,
\qquad x\in K,
\]
kjer je $(\lambda_k)_k$ omejeno zaporedje strogo pozitivnih realnih
števil. Če je $k\neq q$, potem je
\[
\lambda_k\inf_{m\neq k}\varphi_m(x^q)=0,
\]
saj je $\varphi_q(x^q)=0$. Če pa je $k=q$, potem po~\eqref{eq:pozitivnost_funkcionalov} in zaradi
$\varphi_k\longrightarrow\psi$, $\varphi_k\neq\varphi_q$ za vsak
$k\neq q$ ter $\varphi_q\neq\psi$ velja
\[
\lambda_q\inf_{m\neq q}\varphi_m(x^q)>0.
\]
Zato za vsak $q\geq1$ dobimo
\[
f(x^q)
=
\left(\frac12\right)^q
\lambda_q
\inf_{m\neq q}\varphi_m(x^q)\,x^q.
\]
To pokaže, da je $x^q$ lastni vektor preslikave $f$. Števila
$0<\lambda_q\leq1$ lahko izberemo tako, da ima $f$ neskončno mnogo
različnih lastnih vrednosti.
\end{dokaz}

Stožčni spekter preslikave $f$, skonstruirane v dokazu
izreka~\ref{izrek:neskoncno_lastnih_vrednosti}, je števno neskončna
množica. Naslednji zgled pokaže, da lahko stožčni spekter zvezne homogene
preslikave, ki ohranja red, vsebuje neštevno neskončno množico tudi tedaj, ko je stožec
končnorazsežen.

\begin{zgled}\label{zgled:kontinuum_stozcnega_spektra}
Zgled definiramo na stožcu pozitivno semidefinitnih simetričnih matrik
velikosti $2\times2$, $\Pi_2(\R)$, podobne zglede pa lahko konstruiramo
tudi na $\Pi_n(\R)$ za $n>2$. Spomnimo se, da je $\Pi_2(\R)$ sebi
dualen stožec (TO DO: Section 2.4) glede na skalarni produkt
$\langle R,S\rangle=\operatorname{tr}(RS)$.

Naj bo
\[
A=
\begin{bmatrix}
1&0\\
0&0
\end{bmatrix}
\in\partial\Pi_2(\R)
\]
in definirajmo $f\colon\Pi_2(\R)\longrightarrow\Pi_2(\R)$ s predpisom
\[
f(X)=\bigl(\operatorname{tr}(XA)X\bigr)^{1/2},
\qquad X\in\Pi_2(\R).
\]
Če je $X\leq Y$, potem je
$\operatorname{tr}(XA)\leq\operatorname{tr}(YA)$, saj je
$A\in\Pi_2(\R)$. Zato velja
\[
\operatorname{tr}(XA)X
\leq
\operatorname{tr}(YA)X
\leq
\operatorname{tr}(YA)Y.
\]
Ker poleg tega tudi preslikava $R\longmapsto R^{1/2}$ na $\Pi_2(\R)$ ohranja red, sledi
$f(X)\leq f(Y)$, torej $f$ ohranja red. Preslikava $f$ je tudi homogena: 
$f(\alpha X)=\bigl(\operatorname{tr}(\alpha XA)\alpha X\bigr)^{1/2}
=\alpha\bigl(\operatorname{tr}(XA)X\bigr)^{1/2}=\alpha f(X)$ za vsak $\alpha\geq0$ in $X\in\Pi_2(\R)$.

Za $0\leq\alpha\leq1$ postavimo
\[
X_\alpha
=
\begin{bmatrix}
1-\alpha&0\\
0&\alpha
\end{bmatrix}
\in\Pi_2(\R).
\]
Tedaj je $\operatorname{tr}(X_\alpha A)=1-\alpha$ in
\[
\begin{aligned}
f(X_\alpha)
&=
\begin{bmatrix}
(1-\alpha)^2&0\\
0&(1-\alpha)\alpha
\end{bmatrix}^{1/2}\\
&=
\begin{bmatrix}
1-\alpha&0\\
0&\sqrt{(1-\alpha)\alpha}
\end{bmatrix}.
\end{aligned}
\]
Če je $f(X_\alpha)=\lambda X_\alpha$, potem je $\lambda=1$ in
$\sqrt{(1-\alpha)\alpha}=\alpha$, kar je enakovredno pogoju
$\alpha=0$ ali $\alpha=1/2$. Zato sta $X_0$ in $X_{1/2}$ lastna vektorja
preslikave $f$ z lastno vrednostjo $1$. 
Vemo že: če je $x \in \operatorname{int}(K)$, potem $x$ dominira vsak element stožca. 
Naj bo $\lambda$ še ena lastna vrednost preslikave $f$, torej $f(X) = \lambda X$. Ker je
$X_{1/2}\in\operatorname{int}(\Pi_2(\R))$ in $f(X_{1/2}) = 1 X_{1/2}$, iz
leme~\ref{lema:primerjava_lastnih_vrednosti} sledi, da je največja lastna
vrednost preslikave $f$ enaka $1$, saj je $X_{1/2}$ dominira vsak element stožca.

Naj bo sedaj $Z\in\partial\Pi_2(\R)$ in $\operatorname{tr}(Z)=1$. 
Ker je $Z \in \partial \Pi_2(\mathbb{R})$, je matrika $Z$ singularna, zato je vsaj ena
njena lastna vrednost enaka $0$.
Ker je sled matrike enaka vsoti njenih lastnih vrednosti, velja
$\operatorname{tr}(Z)=\lambda_1+\lambda_2=1$.
Ker je ena izmed lastnih vrednosti enaka $0$, mora biti druga enaka $1$.
Tedaj
lahko $Z$ za neki $0\leq\vartheta\leq2\pi$ zapišemo kot
\[
\begin{aligned}
Z
&=
\begin{bmatrix}
\cos\vartheta&-\sin\vartheta\\
\sin\vartheta&\cos\vartheta
\end{bmatrix}
X_0
\begin{bmatrix}
\cos\vartheta&\sin\vartheta\\
-\sin\vartheta&\cos\vartheta
\end{bmatrix}\\
&=
\begin{bmatrix}
\cos^2\vartheta&\cos\vartheta\sin\vartheta\\
\cos\vartheta\sin\vartheta&\sin^2\vartheta
\end{bmatrix}.
\end{aligned}
\]
Hitro lahko preverimo, da velja $Z \cdot Z=Z$, zato je $Z^{1/2}=Z$. Zato je
\[
f(Z)
=
\bigl(\operatorname{tr}(ZA)Z\bigr)^{1/2}
=
|\cos\vartheta|Z^{1/2}
=
|\cos\vartheta|Z.
\]
Torej je vsak $Z\in\partial\Pi_2(\R)$ z $\operatorname{tr}(Z)=1$
lastni vektor z lastno vrednostjo $|\cos\vartheta|$. Stožčni spekter
preslikave $f$ je zato enak intervalu $[0,1]$.
\end{zgled}

V zvezi z zgledom~\ref{zgled:kontinuum_stozcnega_spektra} omenimo še
naslednji odprti problem.

\medskip
\noindent\textbf{Odprti problem.}
Kateri končnorazsežni zaprti stožci $K$ dopuščajo zvezno homogeno
preslikavo $f\colon K\longrightarrow K$, ki ohranja red in za katero
stožčni spekter $\sigma_K(f)$ vsebuje kontinuum?








\subsection{Stožčni spektralni radij}

Glavni del Perron-Frobeniusovega izreka se nanaša na spektralni radij
in pripadajoči lastni vektor. V tem podpoglavju bomo obravnavali
posplošitve spektralnega radija na nelinearne preslikave stožcev.
Sprva bomo predpostavili le, da je $f\colon K\longrightarrow K$ zvezna
homogena preslikava na strnjenem zaprtem stožcu $K$ v
končnorazsežnem vektorskem prostoru $V$.

Obravnavali bomo dve naravni definiciji stožčnega spektralnega radija,
ki ga označimo z $r_K(f)$, in pokazali, da sovpadata. Za razliko od
linearnega primera ni nujno, da obstaja $x\in K\setminus\{0\}$, za
katerega velja $f(x)=r_K(f)x$. V naslednjem podpoglavju bomo pokazali,
da je $r_K(f)$ lastna vrednost preslikave $f$, če dodatno
predpostavimo, da $f$ ohranja red. V linearnem primeru te težave ni,
saj linearna preslikava $f\colon K\longrightarrow K$ nujno ohranja red.

Naj bo $K\subseteq V$ zaprt stožec in
$f\colon K\longrightarrow K$ zvezna homogena preslikava. Ker je $f$
zvezna v $0$ in homogena, lahko definiramo njeno normo na naslednji način.

\begin{definicija}
  Normo $\|f\|_K$ preslikave $f$ na stožcu $K$ definiramo kot
  \[\|f\|_K=\sup\{\|f(x)\| ~|~ x\in K,\ \|x\|\leq1\},\] kjer je
  $\|\cdot\|$ neka fiksna norma na $V$.
\end{definicija}

Ker je $f$ homogena, vemo že, da velja
$\|f(x)\|\leq\|f\|_K\|x\|$ za vsak $x\in K$. Zato za vse
$k,m\geq1$ in $x\in K$ velja
$\|f^{k+m}(x)\|\leq\|f^k\|_K\|f^m\|_K\|x\|$, od koder sledi
$\|f^{k+m}\|_K\leq\|f^k\|_K\|f^m\|_K$.

Po Feketejevi subaditivni lemi~\cite{Fekete1923} zato velja
$\lim\limits_{m\to\infty}\|f^m\|_K^{1/m}
=\inf\limits_{m\geq1}\|f^m\|_K^{1/m}$.
Za linearne preslikave je Bonsall to limito imenoval stožčni
spektralni radij. 

\begin{definicija}
  Za zvezno homogeno preslikavo
  $f\colon K\longrightarrow K$ definiramo njen
  \emph{Bonsallov stožčni spektralni radij} kot
  \[\widehat r_K(f)=\lim\limits_{m\to\infty}\|f^m\|_K^{1/m}=\inf_{m\geq1}\|f^m\|_K^{1/m}.\]
\end{definicija}

Vrednost $\widehat r_K(f)$ ni odvisna od izbire norme, saj že vemo, da so vse
norme na končnorazsežnem vektorskem prostoru ekvivalentne.

Drugo naravno definicijo dobimo z opazovanjem orbit posameznih točk.

\begin{definicija}
  Za vsak $x\in K\setminus\{0\}$ definiramo
  \[\mu_K(x,f)=\limsup_{m\to\infty}\|f^m(x)\|^{1/m}\]
  in nato \emph{stožčni spektralni radij} preslikave $f$ kot
  \begin{align*}
  r_K(f)
  =
  \sup\{\mu_K(x,f) ~|~ x\in K\setminus\{0\}\}.
  \end{align*}
\end{definicija}


Za vsak $x\in K\setminus\{0\}$ in $m\geq1$ velja
$\|f^m(x)\|^{1/m}\leq\|f^m\|_K^{1/m}\|x\|^{1/m}$.
Ker je $\|x\|^{1/m}\longrightarrow1$, iz definicij sledi
$\mu_K(x,f)\leq\widehat r_K(f)$ za vsak $x\in K\setminus\{0\}$ in zato
$r_K(f)\leq\widehat r_K(f)$.

Z uporabo predpostavke, da je $V$ končnorazsežen, bomo pokazali, da
pravzaprav velja $r_K(f)=\widehat r_K(f)$. V neskončnorazsežnih
Banachovih prostorih ta enakost ni nujno izpolnjena.

\begin{izrek}\label{izr:enakost_stozcnih_spektralnih_radijev}
Naj bo $K\subseteq V$ zaprt stožec in
$f\colon K\longrightarrow K$ zvezna homogena preslikava. Potem velja
$r_K(f)=\widehat r_K(f)$.
\end{izrek}

\begin{dokaz}
Ker že vemo, da je $r_K(f)\leq\widehat r_K(f)$, moramo pokazati samo še
obratno neenakost. Predpostavimo nasprotno, da je
$r_K(f)<\widehat r_K(f)$.

Za $\lambda>0$ definirajmo $f_\lambda(x)=f(\lambda x)=\lambda f(x)$.
Iz homogenosti sledi
$r_K(f_\lambda)=\lambda r_K(f)$ in
$\widehat r_K(f_\lambda)=\lambda\widehat r_K(f)$.
Ker je $r_K(f)<\widehat r_K(f)$, lahko $\lambda>0$ izberemo tako, da po
zamenjavi $f$ s $f_\lambda$ velja
$r_K(f)<1<\widehat r_K(f)$.

Za $\alpha>0$ označimo
$B_\alpha=\{x\in K ~|~ \|x\|\leq\alpha\}$.
Ker je $r_K(f)<1$, po definiciji velja, da je $\mu_K(x,f)<1$ za vsak $x\in K\setminus\{0\}$, 
torej obstaja $q$, da je $\limsup\limits_{m\to\infty}\|f^m(x)\|^{1/m} < q < 1$.
Sledi $\limsup\limits_{m\to\infty}\|f^m(x)\| < q^m < 1$. Ker 
je $0 < q < 1$, velja $q^m \to 0$, torej 
za vsak $x\in K$ obstaja $m_x\geq1$, tako da je
$f^{m_x}(x)\in B_{1/3}$. Zaradi zveznosti preslikave $f^{m_x}$ obstaja
relativno odprta okolica $V_x$ točke $x$ v $K$, tako da je
$\|f^{m_x}(y)\|<1/2$ za vsak $y\in V_x$.

Množice $V_x$, kjer je $x\in B_1$, tvorijo odprto pokritje kompaktne
množice $B_1$, zato obstaja končno podkritje. Sledi, da obstajajo $x_1,\ldots,x_p\in B_1$, tako da
$B_1\subseteq\bigcup\limits_{j=1}^p V_{x_j}$. Zaradi preglednosti postavimo
$V_j=B_1\cap V_{x_j}$ in $m_j=m_{x_j}$. Tedaj je
$B_1=\bigcup\limits_{j=1}^p V_j$ in $\|f^{m_j}(y)\|<1/2$ za vsak $y\in V_j$.

Naj bo $M=\max\limits_{1\leq j\leq p}m_j$ in definirajmo
$W=\bigcup\limits_{j=0}^{M}f^j(B_1)$. Pokazali bomo, da je $f(W)\subseteq W$.
Naj bo $w\in W$. Potem je $w=f^j(x)$ za neki $x\in B_1$ in
$0\leq j\leq M$. Če je $j<M$, je
$f(w)=f^{j+1}(x)\in W$.
Če pa je $j=M$, obstaja $1\leq k\leq p$, tako da je $x\in V_k$.
Tedaj je $\xi:=f^{m_k}(x)\in B_{1/2}\subseteq B_1$. Ker je
$1\leq m_k\leq M$, sledi $M-m_k+1\leq M$ in zato
$f(w)=f^{M-m_k+1}(\xi)\in W$. Torej je $f(W)\subseteq W$.

Množica $B_1$ je kompaktna in preslikava $f$ je zvezna. Ker je zvezna slika kompaktne množice kompaktna, 
je množica $W$ kompaktna in zato omejena. Torej obstaja $R>0$, tako da je
$\|w\|\leq R$ za vsak $w\in W$. Ker je $B_1\subseteq W$ in
$f(W)\subseteq W$, za vsak $x\in B_1$ in $m\geq1$ velja
$\|f^m(x)\|\leq R$. Posledično je $\|f^m\|_K\leq R$ za vsak $m\geq1$,
zato
\begin{align*}
\widehat r_K(f)
=
\lim_{m\to\infty}\|f^m\|_K^{1/m}
\leq
\lim_{m\to\infty}R^{1/m}
=
1.
\end{align*}

To je v protislovju s predpostavko $\widehat r_K(f)>1$. Zato mora
veljati $r_K(f)=\widehat r_K(f)$.
\end{dokaz}

Stožčni spektralni radij ima naslednje osnovne lastnosti.

\begin{trditev}\label{trd:lastnosti_stozcnega_spektralnega_radija}
Naj bo $K\subseteq V$ zaprt stožec in
$g\colon K\longrightarrow K$ zvezna homogena preslikava. Potem velja:
\begin{enumerate}[label=(\roman*)]
\item Za vsak $k\geq1$ velja $r_K(g^k)=r_K(g)^k$.
\item Če za neki $x\in K\setminus\{0\}$ in $k\geq1$ velja
$g^k(x)=\lambda^k x$, potem je $\lambda\leq r_K(g)$.
\end{enumerate}

Če sta poleg tega $K_1\subseteq V_1$ in $K_2\subseteq V_2$ zaprta
stožca ter 
$f_1\colon K_1\longrightarrow K_2$ in
$f_2\colon K_2\longrightarrow K_1$ zvezni homogeni preslikavi.
Potem velja
$r_{K_1}(f_2\circ f_1)=r_{K_2}(f_1\circ f_2)$.
\end{trditev}

\begin{dokaz}
Po izreku~\ref{izr:enakost_stozcnih_spektralnih_radijev} zadostuje,
da trditve dokažemo za Bonsallov stožčni spektralni radij. Za
$k\geq1$ velja

\begin{align*}
\widehat r_K(g^k)
&=
\lim_{m\to\infty}\|g^{km}\|_K^{1/m}
=
\lim_{m\to\infty}
\left(\|g^{km}\|_K^{1/(mk)}\right)^k
=
\widehat r_K(g)^k.
\end{align*}

Za dokaz točke (ii) naj bo $g^k(x)=\lambda^kx$ za neki
$x\in K\setminus\{0\}$. Tedaj je $g^{km}(x)=\lambda^{km}x$ za vsak
$m\geq1$ in zato dobimo 
\begin{align*}
\lambda
=
\lim_{m\to\infty}\|g^{km}(x)\|^{1/(mk)}
\leq
\limsup_{j\to\infty}\|g^j(x)\|^{1/j}
\leq
r_K(g).
\end{align*}

Dokažimo še zadnjo trditev. Naj bo $\|\cdot\|_1$ norma na $V_1$ in
$\|\cdot\|_2$ norma na $V_2$. Definirajmo
$\|f_1\|=\sup\{\|f_1(x)\|_2 ~|~ x\in K_1,\ \|x\|_1\leq1\}$ in podobno
$\|f_2\|=\sup\{\|f_2(y)\|_1 ~|~ y\in K_2,\ \|y\|_2\leq1\}$.
Za vsak $m\geq1$ in $x\in K_1$ velja
\begin{align*}
\|(f_2\circ f_1)^{m+1}(x)\|_1
&=
\left\|
f_2\left((f_1\circ f_2)^m(f_1(x))\right)
\right\|_1\\
&\leq
\|f_2\|\,
\|(f_1\circ f_2)^m\|_{K_2}\,
\|f_1\|\,\|x\|_1.
\end{align*}

Zato je
$\|(f_2\circ f_1)^{m+1}\|_{K_1}
\leq\|f_2\|\|(f_1\circ f_2)^m\|_{K_2}\|f_1\|$.
Če je $\|f_1\|=0$ ali $\|f_2\|=0$, je
$\widehat r_{K_1}(f_2\circ f_1)=0
\leq\widehat r_{K_2}(f_1\circ f_2)$.

Če sta $\|f_1\|$ in $\|f_2\|$ neničelni, dobimo
\begin{align*}
\widehat r_{K_1}(f_2\circ f_1)
&\leq
\lim_{m\to\infty}
\|f_2\|^{1/(m+1)}
\left(
\|(f_1\circ f_2)^m\|_{K_2}^{1/m}
\right)^{m/(m+1)}
\|f_1\|^{1/(m+1)}\\
&=
\widehat r_{K_2}(f_1\circ f_2).
\end{align*}

Po izreku~\ref{izr:enakost_stozcnih_spektralnih_radijev} zato velja
$r_{K_1}(f_2\circ f_1)\leq r_{K_2}(f_1\circ f_2)$.
Če zamenjamo vlogi $K_1$ in $K_2$, dobimo še obratno neenakost, zato
je $r_{K_1}(f_2\circ f_1)=r_{K_2}(f_1\circ f_2)$.
\end{dokaz}

Naravno se je vprašati, ali je stožčni spektralni radij zvezne
homogene preslikave $f\colon K\longrightarrow K$ vedno njena lastna
vrednost. Definirajmo
\begin{equation*}
  \widetilde r_K(f)=
  \sup\{\lambda\geq0 ~|~ f(x)=\lambda x
  \text{ za neki }x\in K\setminus\{0\}\}.
\end{equation*}
Po točki (ii) trditve~\ref{trd:lastnosti_stozcnega_spektralnega_radija}
velja $\widetilde r_K(f)\leq r_K(f)$. Morda bi pričakovali, da vedno
velja enakost, vendar naslednji zgled pokaže, da je lahko neenakost
stroga tudi v končnorazsežnem prostoru.

\begin{zgled}\label{zgled:spektralni_radij_ni_lastna_vrednost}
Poglejmo Lorentzov stožec
$\Lambda_3=\{(x_1,x_2,x_3)\in\R^3:
x_1^2-x_2^2-x_3^2\geq0,\ x_1\geq0\}$.
Naj $U_\vartheta\colon\R^2\longrightarrow\R^2$ označuje rotacijo za
kot $\vartheta$ okoli izhodišča in predpostavimo, da je
$\vartheta/\pi$ iracionalno število. Definirajmo
$f_\vartheta\colon\Lambda_3\longrightarrow\Lambda_3$ s predpisom

\begin{align*}
f_\vartheta(x)
=
\left(
\frac12x_1+\sqrt{x_2^2+x_3^2},
U_\vartheta(x_2,x_3)
\right).
\end{align*}

Preprosto lahko preverimo, da je $f_\vartheta(\Lambda_3)\subseteq\Lambda_3$
in da je $f_\vartheta$ zvezna homogena preslikava. 
Res, označimo $r = \sqrt{x_2^2+x_3^2}$. Očitno je $\frac12x_1+r \geq r \geq 0$ za vse $x\in\Lambda_3$. 
Ker rotacija ne spremeni dolžine vektorja, je $\|U_\vartheta(x_2,x_3)\|=\|(x_2,x_3)\|=r$ in 
velja $\left(\frac12x_1+\sqrt{x_2^2+x_3^2}\right)^2 - r^2 \geq \frac14x_1^2 \geq 0$, torej 
velja $f_\vartheta(x)\in\Lambda_3$ za vsak $x\in\Lambda_3$. Da je $f_\vartheta$ homogena, 
sledi iz linearnosti rotacije in linearne funkcije $x_1 \mapsto \frac12 x_1 + r$. Zveznost je očitna.

Poglejmo si $f_\vartheta^2$. Računamo $f_\vartheta^2(x_1, x_2, x_3) = f_\vartheta(\frac12x_1+\sqrt{x_2^2+x_3^2}, U_\vartheta(x_2,x_3))=
f_\vartheta(\frac{1}{2} x_1 + r, U_\vartheta(x_2,x_3)) = (\frac{1}{2} (\frac{1}{2} x_1 + r) + \| U_\vartheta(x_2,x_3) \|,
U_\vartheta^2(x_2,x_3))$. Ker velja, da je $\| U_\vartheta(x_2,x_3) \| = r$, dobimo
$f_\vartheta^2(x_1, x_2, x_3) = 
\left(\frac{1}{4} x_1 + \frac{3}{2} r, U_\vartheta^2(x_2,x_3)\right)$.
Splošno lahko pokažemo, da za vsak $k\geq1$ velja
\begin{align*}
f_\vartheta^k(x)
=
\left(
\left(\frac12\right)^k x_1
+
\sqrt{x_2^2+x_3^2}
\sum_{j=0}^{k-1}\left(\frac12\right)^j,
U_\vartheta^k(x_2,x_3)
\right).
\end{align*}

Ker je $\vartheta/\pi$ iracionalno, velja
$U_\vartheta(x_2,x_3)=\lambda(x_2,x_3)$ za neko $\lambda\geq0$
natanko tedaj, ko je $(x_2,x_3)=(0,0)$. Zato je
$x\in\Lambda_3\setminus\{0\}$ lastni vektor preslikave $f_\vartheta$
natanko tedaj, ko je $x_1>0$ in $x_2=x_3=0$. V tem primeru je
$f_\vartheta(x)=\frac12x$, zato je
$\widetilde r_{\Lambda_3}(f_\vartheta)=1/2$.

Izberimo sedaj $x=(x_1,x_2,x_3)$ z $x_1\geq0$ in
$x_1^2=\frac{1}{2}=x_2^2+x_3^2$. Tedaj je $\|x\|_2=1$ in iz zgornjega izraza
za $f_\vartheta^k(x)$ sledi
$\|f_\vartheta^k\|_{\Lambda_3}\geq1$ za vsak $k\geq1$. Zato je
$r_{\Lambda_3}(f_\vartheta)\geq1$ in posledično
$\widetilde r_{\Lambda_3}(f_\vartheta) = \frac12
< 1 \leq
r_{\Lambda_3}(f_\vartheta)$.

Torej je največja lastna vrednost $\frac{1}{2}$, stožčni spektralni radij pa je vsaj $1$. 
Zato spektralni radij v splošnem ni nujno lastna vrednost.
\end{zgled}

Da bi zagotovili, da je stožčni spektralni radij lastna vrednost,
potrebujemo dodatne predpostavke. V naslednjem podpoglavju bomo
pokazali, da zadostuje predpostavka, da $f$ ohranja red. Pravzaprav
zadostuje nekoliko splošnejši pogoj: obstaja zaprt stožec
$C\subseteq V$ s $K\subseteq C$, tako da je
$f\colon K\longrightarrow K$ zvezna homogena preslikava, ki ohranja
red glede na $C$.

Predpostavka, da $f$ ohranja red glede na $C$ namesto glede na $K$,
se morda zdi umetna. V številnih naravnih primerih $f$ ohranja red
glede na $K$, lahko pa se zgodi, da je lažje dokazati, da $f$ ohranja
red glede na večji stožec $C$. To ponazarja naslednji zgled.

\begin{zgled}\label{zgled:harmonicna_preslikava}
Naj bodo za $1\leq i,j\leq n$ števila $\alpha_{ij}\geq0$, za katera
velja $\sum\limits_{j=1}^n\alpha_{ij}>0$ za vsak $i$. Definirajmo
$f\colon\operatorname{int}(\R_+^n)\longrightarrow\R_+^n$ s predpisom
\[f_i(x)=\left(\sum\limits_{j=1}^n\alpha_{ij} \frac{1}{x_j}\right)^{-1}\] za
$1\leq i\leq n$. Preslikava $f$ je očitno zvezna, homogena in ohranja red
glede na $\R_+^n$. Po
izreku~\ref{izr:zvezna_razsiritev_poliedrski_stozec} ima zvezno
homogeno razširitev $\R_+^n \to \R_+^n$, ki prav tako ohranja red.

Naj bo
$K=\{x\in\R^n ~|~ 0\leq x_1\leq x_2\leq\cdots\leq x_n\} \subseteq \R_+^n$. Očitno je to zaprt stožec.
Trdimo, da je $f(K)\subseteq K$ natanko tedaj, ko za
$1\leq i<n$ in $1\leq k\leq n$ velja
\begin{align}
\sum\limits_{j=1}^k\alpha_{ij}
\geq
\sum\limits_{j=1}^k\alpha_{(i+1)j}.
\label{eq:4.3}
\end{align}

Dokažimo najprej iz desne proti levi. Zaradi zveznosti preslikave $f$ je $f(K)\subseteq K$ natanko tedaj,
ko je $f(x)\in K$ za vsak $x\in K$ z
$0<x_1\leq x_2\leq\cdots\leq x_n$. Za tak $x$ je pogoj
$f(x)\in K$ ekvivalenten pogoju
\begin{align}
\sum\limits_{j=1}^n\alpha_{ij}x_j^{-1}
\geq
\sum\limits_{j=1}^n\alpha_{(i+1)j}x_j^{-1}
\qquad
\text{za }1\leq i<n.
\label{eq:4.4}
\end{align}

Postavimo $y_j=x_j^{-1}$. Tedaj je
$y_1\geq y_2\geq\cdots\geq y_n>0$, pri čemer je $y_j\to0$, ko
$x_j\to\infty$. Zgornji pogoj je zato ekvivalenten neenakosti
$\sum\limits_{j=1}^n\alpha_{ij}y_j
\geq\sum\limits_{j=1}^n\alpha_{(i+1)j}y_j$
za vse $y_1\geq y_2\geq\cdots\geq y_n\geq0$ 
oziroma $\sum\limits_{j=1}^n(\alpha_{ij}-\alpha_{(i+1)j})y_j
\geq 0$.

Za fiksen $i$ postavimo
$\beta_j=\alpha_{ij}-\alpha_{(i+1)j}$,
$S_k=\sum\limits_{j=1}^k\beta_j$ in $S_0=0$.
Če izrazimo $\beta_j$ v terminih $S_k$, dobimo
$\beta_1=S_1-S_0 = S_1$, $\beta_2=S_2-\beta_1=S_2-S_1$, 
$\beta_3= S_3 - \beta_2 - \beta_1 = S_3 - (S_2-S_1) - S_1 = S_3 - S_2$ in 
splošno $\beta_j=S_j-S_{j-1}$ za $1\leq j\leq n$.
S sumacijo po delih dobimo
\begin{align*}
\sum\limits_{j=1}^n\beta_j y_j
= \sum\limits_{j=1}^n (S_j-S_{j-1}) y_j
=\sum\limits_{k=1}^{n-1}S_k(y_k-y_{k+1})+S_n y_n.
\end{align*}

Ker je $y_1\geq y_2\geq\cdots\geq y_n\geq0$ in po zgornjem pogoju
$S_k\geq0$ za vsak $1\leq k\leq n$, je desna stran nenegativna.
Zato velja neenakost~\eqref{eq:4.4}, torej $f(K)\subseteq K$.

Za dokaz trditve iz leve proti desni fiksirajmo $1\leq k\leq n$ in izberimo
$y_j=1$ za $1\leq j\leq k$ ter $y_j=0$ za $j>k$. Tedaj dobimo
$\sum\limits_{j=1}^k\alpha_{ij}
\geq\sum\limits_{j=1}^k\alpha_{(i+1)j}$.
S tem smo pokazali želeno ekvivalenco $f(K)\subseteq K$ natanko tedaj,
ko velja~\eqref{eq:4.3}.
\end{zgled}

Čeprav je $f(K)\subseteq K$, se lahko zgodi, da $f$ ne ohranja reda
glede na $K$. Na primer lahko uporabimo izrek~\ref{izrLipFreLeb}, da pokažemo, da $f$ ne ohranja reda glede na
$K$, če velja
\begin{align*}
\alpha_{nn}\bigl(\alpha_{(n-1)1}\bigr)^2
-
\alpha_{(n-1)n}\bigl(\alpha_{n1}\bigr)^2
<0.
\end{align*}

\begin{zgled}
  Pokažimo, da $f$ ne ohranja reda glede na $K$, če velja
  $\alpha_{nn}(\alpha_{(n-1)1})^2-\alpha_{(n-1)n}(\alpha_{n1})^2<0$. 
  Iz prejšnjega zgleda vemo, da je $f(K)\subseteq K$.

  Po izreku~\ref{izrLipFreLeb} bi moralo v primeru, ko bi $f$ ohranjala red glede na $K$, za vsak
  $x\in\operatorname{int}(K)$ veljati $Df(x)(K)\subseteq K$. Pokazali bomo, da ta
  pogoj ni izpolnjen.

  Vzemimo $e_n=(0,\ldots,0,1)\in K$. Za $1\leq i\leq n$ označimo
  $A_i(x)=\sum\limits_{j=1}^n\alpha_{ij}x_j^{-1}$, tako da je
  $f_i(x)=A_i(x)^{-1}$. Ker je
  $\frac{\partial A_i}{\partial x_n}(x)=-\alpha_{in}x_n^{-2}$, dobimo
  \begin{align*}
  \frac{\partial f_i}{\partial x_n}(x)
  &=-A_i(x)^{-2}\frac{\partial A_i}{\partial x_n}(x) \\
  &=\frac{\alpha_{in}}{x_n^2A_i(x)^2}.
  \end{align*}
  Ker je $Df(x)e_n$ ravno $n$-ti stolpec matrike $Df(x)$, je njegova $i$-ta
  koordinata enaka
  $\frac{\alpha_{in}}{x_n^2A_i(x)^2}$.

  Za $M>1$ izberimo
  $x=(1,M,M^2,\ldots,M^{n-1})\in\operatorname{int}(K)$. Tedaj pri
  $M\to\infty$ velja
  $\sum\limits_{j=1}^n\alpha_{ij}x_j^{-1}\to\alpha_{i1}$. Zato
  \[\alpha_{(n-1)n}
  \left(\sum\limits_{j=1}^n\alpha_{nj}x_j^{-1}\right)^2
  -
  \alpha_{nn}
  \left(\sum\limits_{j=1}^n\alpha_{(n-1)j}x_j^{-1}\right)^2
  \longrightarrow
  \alpha_{(n-1)n}(\alpha_{n1})^2
  -
  \alpha_{nn}(\alpha_{(n-1)1})^2.\]

  Po predpostavki je zadnji izraz pozitiven, zato za dovolj velik $M$ velja
  \[
  \alpha_{(n-1)n}
  \left(\sum\limits_{j=1}^n\alpha_{nj}x_j^{-1}\right)^2
  >
  \alpha_{nn}
  \left(\sum\limits_{j=1}^n\alpha_{(n-1)j}x_j^{-1}\right)^2.
  \]
  Od tod sledi
  $(Df(x)e_n)_{n-1}>(Df(x)e_n)_n$, zato $Df(x)e_n\notin K$.

  Ker je $e_n\in K$, velja $Df(x)(K)\nsubseteq K$. Po izreku~\ref{izrLipFreLeb} zato
  preslikava $f$ ne ohranja reda glede na $K$.
\end{zgled}

Podpoglavje zaključimo z nekaj uporabnimi neenakostmi za stožčni
spektralni radij.

\begin{trditev}\label{trd:monotonost_stozcnega_spektralnega_radija}
Naj bosta $K\subseteq C$ zaprta stožca v $V$ ter
$f,g\colon K\longrightarrow K$ zvezni homogeni preslikavi, za kateri
velja $f(x)\leq_Cg(x)$ za vsak $x\in K$. Če vsaj ena izmed preslikav
$f$ in $g$ ohranja red glede na $C$, potem je
$r_K(f)\leq r_K(g)$.
\end{trditev}

\begin{dokaz}
Najprej z indukcijo pokažimo, da je
$f^k(x)\leq_C g^k(x)$ za vsak $x\in K$ in $k\geq1$. Za $k=1$ je to
ravno predpostavka. Predpostavimo, da trditev velja za $k$.

Če $f$ ohranja red glede na $C$, potem
\begin{align*}
f^{k+1}(x)
=
f(f^k(x))
\leq_C
f(g^k(x))
\leq_C
g(g^k(x))
=
g^{k+1}(x).
\end{align*}

Če pa $g$ ohranja red glede na $C$, dobimo
\begin{align*}
f^{k+1}(x)=f(f^k(x))
\leq_Cg(f^k(x))
\leq_Cg(g^k(x))
=g^{k+1}(x).
\end{align*}

Ker je $C$ zaprt stožec v končnorazsežnem vektorskem prostoru, je
normalen. Iz definicije normalnosti sledi, da obstaja $M>0$, tako da iz $u\leq_C v$ sledi
$\|u\|\leq M\|v\|$. 
Ker sta $f^k(x),g^k(x)\in K\subseteq C$, velja
$0\leq_C f^k(x)\leq_C g^k(x)$, zato velja 
$\|f^k(x)\|\leq M\|g^k(x)\|$ za vsak $x\in K$ in $k\geq1$.
Od tod sledi
\begin{align*}
\limsup_{k\to\infty}\|f^k(x)\|^{1/k}
\leq
\limsup_{k\to\infty} M^{1/k}\|g^k(x)\|^{1/k}
= \limsup_{k\to\infty}\|g^k(x)\|^{1/k}.
\end{align*}

Torej je $\mu_K(x,f)\leq\mu_K(x,g)$ za vsak
$x\in K\setminus\{0\}$. Če vzamemo supremum po vseh takih $x$, dobimo
$r_K(f)\leq r_K(g)$.
\end{dokaz}

\begin{trditev}\label{trd:spodnja_meja_stozcnega_spektralnega_radija}
Naj bosta $K\subseteq C$ zaprta stožca v $V$ in naj bo
$f\colon K\longrightarrow K$ zvezna homogena preslikava, ki ohranja
red glede na $C$. Če obstajata $u\in K\setminus\{0\}$ in $\rho>0$,
tako da je $\rho u\leq_C f(u)$, potem velja $r_K(f)\geq\rho$.

Če ima $K$ poleg tega neprazno notranjost, potem za vsak
$x\in\operatorname{int}(K)$ velja
$r_K(f)=\lim\limits_{k\to\infty}\|f^k(x)\|^{1/k}$.
\end{trditev}

\begin{dokaz}
Iz predpostavke $\rho u\leq_C f(u)$ ter homogenosti in ohranjanja reda
preslikave $f$ z indukcijo dobimo
$\rho^k u\leq_C f^k(u)$ za vsak $k\geq1$.
Ker je $C$ normalen, obstaja konstanta $M>0$, tako da
$\rho^k\|u\|\leq M\|f^k(u)\|$. Vzamemo $k$-ti koren in zato sledi
\begin{align*}
\rho
\leq
\limsup_{k\to\infty}\|f^k(u)\|^{1/k}
\leq
r_K(f).
\end{align*}

Naj bo sedaj $x\in\operatorname{int}(K)$. Ker je $x$ notranja točka, obstaja $\varepsilon>0$, tako da je
$B(x, \varepsilon) \subseteq\operatorname{int}(K)$. Če je $ y \in K$ z $\|y\|\leq1$, potem obstaja $\delta>0$, 
tako da je $x-\delta y \in K$, kar pa pomeni $y \leq_K \frac{1}{\delta} x$. Označimo $M' = \frac{1}{\delta}$ in 
potem velja $y\leq_C M'x$ za vsak $y\in K$ z $\|y\|\leq1$. Ker $f$
ohranja red glede na $C$ in je homogena, velja
$f^k(y)\leq_C M'f^k(x)$ za vsak $k\geq1$.

Z normalnostjo stožca $C$ zato dobimo
$\|f^k(y)\|\leq MM'\|f^k(x)\|$ za vsak $y\in K$ z $\|y\|\leq1$.
Če vzamemo supremum po vseh takih $y$, sledi
$\|f^k\|_K\leq MM'\|f^k(x)\|$. Po drugi strani iz definicije
$\|f^k\|_K$ velja
$\|f^k(x)\|\leq\|f^k\|_K\|x\|$, torej dobimo
\begin{align*}
\|f^k\|_K
\leq
M M' \|f^k(x)\|
\leq
M M' \|f^k\|_K\|x\|.
\end{align*}

Ko vse člene potenciramo z $1/k$ in pošljemo $k\to\infty$, konstante
$M$, $M'$ in $\|x\|$ po korenjenju konvergirajo proti $1$. Ker po
izreku~\ref{izr:enakost_stozcnih_spektralnih_radijev} velja
$\lim\limits_{k\to\infty}\|f^k\|_K^{1/k}=r_K(f)$, dobimo
\begin{align*}
\lim\limits_{k\to\infty}\|f^k(x)\|^{1/k}
=
r_K(f).
\end{align*}
\end{dokaz}






\subsection{Lastni vektorji, ki pripadajo stožčnemu spektralnemu radiju}

Tako kot v klasični Perron-Frobeniusovi teoriji je stožčni spektralni
radij zvezne homogene preslikave, ki ohranja red, lastna vrednost s
pripadajočim lastnim vektorjem v stožcu. Dokaz temelji na Brouwerjevem
izreku o negibni točki~\cite{Ciarlet2025}. Ideja, da lahko teorijo negibnih točk uporabimo
za dokaz obstoja lastnih vektorjev, je stara in sega do del
Alexandroffa in Hopfa ter Kreina in Rutmana.

\begin{izrek}[Brouwerjev izrek o negibni točki]\label{Brouwer}
Naj bo $K$ kompaktna in konveksna podmnožica končnorazsežnega normiranega vektorskega prostora 
in naj bo $f:K\to K$ zvezna preslikava. Potem ima $f$ vsaj eno negibno točko.
\end{izrek}

\begin{izrek}\label{izr:obstoj_lastnega_vektorja_perturbacija}
Naj bosta $K\subseteq C$ zaprta stožca v $V$, pri čemer ima $K$
neprazno notranjost, in naj bo
$f\colon\operatorname{int}(K)\longrightarrow K$ zvezna homogena
preslikava, ki ohranja red glede na $C$. Naj bosta
$\vartheta\in\operatorname{int}(C^*)$ in
$u\in\operatorname{int}(K)$. Za $\varepsilon>0$ definirajmo
$f_\varepsilon\colon\operatorname{int}(K)\longrightarrow
\operatorname{int}(K)$ s predpisom
\begin{align*}
f_\varepsilon(x)
=
f(x)+\varepsilon\vartheta(x)u,
\qquad
x\in\operatorname{int}(K).
\end{align*}
Potem veljajo naslednje trditve:
\begin{enumerate}[label=(\roman*)]
\item Preslikava $f_\varepsilon$ je zvezna, homogena in ohranja red
glede na $C$.

\item Obstaja $x_\varepsilon\in\operatorname{int}(K)$, za katerega je
$\vartheta(x_\varepsilon)=1$ in
$f_\varepsilon(x_\varepsilon)
=r_K(f_\varepsilon)x_\varepsilon$.

\item Če je $0<\eta<\varepsilon$, potem je
$r_K(f_\eta)<r_K(f_\varepsilon)$, zato obstaja limita
$r=\lim\limits_{\varepsilon\to0}r_K(f_\varepsilon)$.

\item Obstaja konvergentno zaporedje $(x_{\varepsilon_k})_k$ v množici
$\{x_\varepsilon ~|~0<\varepsilon\leq1\}$, za katero je
$x_{\varepsilon_k}\longrightarrow x$ in
$\varepsilon_k\longrightarrow0$, pri čemer
$f(x_{\varepsilon_k})\longrightarrow y$ in $y=rx$.
\end{enumerate}
\end{izrek}

\begin{dokaz}
(i) Preslikava $f_\varepsilon$ je očitno zvezna, ker sta $f$ in $\vartheta$ zvezni, in homogena, 
ker je $f$ homogena in $\vartheta$ linearna. Ker
$f$ ohranja red glede na $C$ in je $\vartheta\in C^*$, tudi
$f_\varepsilon$ ohranja red glede na $C$. Poleg tega je
$f_\varepsilon(\operatorname{int}(K))
\subseteq\operatorname{int}(K)$.

(ii) Naj bo $\Sigma=\{x\in K ~|~ \vartheta(x)=1\}$. Po
lemi~\ref{lema2.46} je $\Sigma$ kompaktna konveksna množica.
Ker je $u\in\operatorname{int}(K)\subseteq\operatorname{int}(C)$, obstaja $\delta>0$, tako da je 
$B(u,\delta)\subseteq\operatorname{int}(C)$. 
Množica $\Sigma\subseteq K$ je kompaktna in zato omejena, torej obstaja $R>0$, tako da je $\|x\| \leq R$ 
za vsak $x\in\Sigma$. Zdaj izberemo $M_1 > \frac{R}{\delta} > 0$. Za poljuben $x\in\Sigma$ velja 
$\|\frac{x}{M_1}\| \leq \delta$, zato je 
\begin{align}
x\leq_C M_1u,
\qquad
\text{za vsak }x\in\Sigma.
\label{eq:dominacija_sigma_u}
\end{align}
Za $x\in\Sigma\cap\operatorname{int}(K)$ je
$\vartheta(x)=1$, zato iz homogenosti in ohranjanja reda preslikave $f$
sledi
$f_\varepsilon(x)=f(x)+\varepsilon u
\leq_C M_1f(u)+\varepsilon u$.
Ker je $u\in\operatorname{int}(C)$, obstaja tudi $M_2>0$, tako da
\begin{align}
M_1f(u)\leq_C M_2u.
\label{eq:M1fu_M2u}
\end{align}
Zato je
\begin{align}
f_\varepsilon(x)
\leq_C
(M_2+\varepsilon)u,
\qquad
\text{za vsak }
x\in\Sigma\cap\operatorname{int}(K).
\label{eq:feps_zgornja_meja}
\end{align}
Ker je $f(x)\in K\subseteq C$ in $\vartheta\in C^*$, po definiciji dualnega stožca velja
$\vartheta(f(x))\geq0$. Iz
$f_\varepsilon(x)=f(x)+\varepsilon u$ in~\eqref{eq:feps_zgornja_meja} tako dobimo
\begin{align}
\varepsilon\vartheta(u)
\leq
\vartheta(f_\varepsilon(x))
\leq
(M_2+\varepsilon)\vartheta(u)
\qquad
\text{za vsak }
x\in\Sigma\cap\operatorname{int}(K).
\label{eq:vartheta_feps}
\end{align}

Definirajmo sedaj
$g_\varepsilon\colon\Sigma\cap\operatorname{int}(K)
\longrightarrow\Sigma\cap\operatorname{int}(K)$ s predpisom
\begin{align*}
g_\varepsilon(x)
=
\frac{f_\varepsilon(x)}
{\vartheta(f_\varepsilon(x))}.
\end{align*}
Vemo, da za $x\in\Sigma$ velja $f_\varepsilon(x)=f(x)+\varepsilon u$. Ker je $f(x)\in K$, je 
$0\leq_K f(x)$, zato $\varepsilon u\leq_K f(x)+\varepsilon u=f_\varepsilon(x)$. Ker je 
$(M_2+\varepsilon)\vartheta(u)>0$, lahko obe strani delimo in dobimo 
$\frac{\varepsilon u}{(M_2+\varepsilon)\vartheta(u)}\leq_K\frac{f_\varepsilon(x)}{(M_2+\varepsilon)\vartheta(u)}$.
Desno neenakost iz~\eqref{eq:vartheta_feps} invertiramo in množimo z $f_\varepsilon(x)$, da dobimo
$\frac{f_\varepsilon(x)}{(M_2+\varepsilon)\vartheta(u)}\leq_K g_\varepsilon(x)$. S tem dobimo
\begin{align*}
\frac{\varepsilon u}
{(M_2+\varepsilon)\vartheta(u)}
\leq_K
\frac{f_\varepsilon(x)}
{(M_2+\varepsilon)\vartheta(u)}
\leq_K
g_\varepsilon(x).
\end{align*}

Ker je $C$ zaprt stožec v končnorazsežnem vektorskem prostoru, je
normalen. Iz~\eqref{eq:dominacija_sigma_u} in~\eqref{eq:M1fu_M2u} zato obstaja $M_3>0$, tako da je
$\|f(x)\|\leq M_3\|u\|$ za vsak
$x\in\Sigma\cap\operatorname{int}(K)$. Ker je
$f(x)\in K$ in $u\in\operatorname{int}(K)$, obstaja $M>0$, tako da je
$f(x)\leq_K Mu$ za vsak
$x\in\Sigma\cap\operatorname{int}(K)$. Levo neenakost iz~\eqref{eq:vartheta_feps} invertiramo and množimo 
z $f_\varepsilon(x)$, da dobimo $g_\varepsilon(x) \leq \frac{f_\varepsilon(x)}{\varepsilon\vartheta(u)} = 
\frac{f(x)+\varepsilon u}{\varepsilon\vartheta(u)}$. Uporabimo prejšnji razmislek in dobimo
\begin{align*}
g_\varepsilon(x)
\leq_K
\frac{f(x)+\varepsilon u}
{\varepsilon\vartheta(u)}
\leq_K
\frac{(M+\varepsilon)u}
{\varepsilon\vartheta(u)}
\end{align*}
za vsak $x\in\Sigma\cap\operatorname{int}(K)$.

Našli smo zgornjo in spodnjo mejo za $g_\varepsilon (x)$, torej obstajata konstanti $0<A\leq B$, tako da
\begin{align}
g_\varepsilon
\bigl(\Sigma\cap\operatorname{int}(K)\bigr)
\subseteq
\{x\in\Sigma\cap\operatorname{int}(K) ~|~ Au\leq_K x\leq_K Bu\}.
\label{eq:geps_redovni_interval}
\end{align}
 
Definirajmo
$\Omega_\varepsilon
=\{x\in\Sigma\cap\operatorname{int}(K) ~|~ d_H(x,u)\leq\log(B/A)\}$,
kjer je Hilbertova metrika določena s stožcem $K$.
Množica $\Omega_\varepsilon$ je kompaktna, ker je zaprta podmnožica kompaktne množice, in po
lemi~\ref{lema:konveksnost_hilbertovih_krogel} konveksna.
Iz~\eqref{eq:geps_redovni_interval} sledi
$g_\varepsilon(\Omega_\varepsilon)\subseteq\Omega_\varepsilon$.
Po Brouwerjevem izreku o negibni točki~\ref{Brouwer} zato obstaja
$x_\varepsilon\in\Omega_\varepsilon
\subseteq\Sigma\cap\operatorname{int}(K)$, tako da je
$g_\varepsilon(x_\varepsilon)=x_\varepsilon$. Po definiciji preslikave $g_\varepsilon$ dobimo 
$f_\varepsilon(x_\varepsilon)= \vartheta (f_\varepsilon(x_\varepsilon)) x_\varepsilon$.

Postavimo
$\lambda_\varepsilon
=\vartheta(f_\varepsilon(x_\varepsilon))$.
Tedaj je
$f_\varepsilon(x_\varepsilon)
=\lambda_\varepsilon x_\varepsilon$ in
$\vartheta(x_\varepsilon)=1$.
Po trditvi~\ref{trd:spodnja_meja_stozcnega_spektralnega_radija} velja
$r_K(f_\varepsilon)
=\lim\limits_{k\to\infty} \|f_\varepsilon^k(x_\varepsilon)\|^{1/k} 
= \lim\limits_{k\to\infty} \| \lambda_\varepsilon^k x_\varepsilon\|^{1/k}
= \lambda_\varepsilon \lim\limits_{k\to\infty} \| x_\varepsilon\|^{1/k}
=\lambda_\varepsilon$.

(iii) Naj bo sedaj $0<\eta<\varepsilon$. Po definiciji preslikav $f_\eta$ in $f_\varepsilon$ velja 
$f_\eta(x)=f_\eta(x)+\eta \vartheta (x) u$ in $f_\varepsilon(x)=f_\varepsilon(x)+\varepsilon \vartheta (x) u$ za vse $x \in \Sigma \cap \operatorname{int}(K)$.
Ker po točki (ii) velja
$f_\eta(x_\eta)=r_K(f_\eta)x_\eta$ ter $\vartheta (x_\eta)=1$, sledi
$f_\varepsilon(x_\eta)
=r_K(f_\eta)x_\eta+(\varepsilon-\eta)u$.
Ker je $u\in\operatorname{int}(C)$, obstaja $\delta>0$, tako da je
$\delta x_\eta\leq_C(\varepsilon-\eta)u$. Zato
\begin{align}
\bigl(r_K(f_\eta)+\delta\bigr)x_\eta
\leq_C
f_\varepsilon(x_\eta).
\label{eq:feps_eta_spodnja_meja}
\end{align}

Ker je $x_\varepsilon\in\operatorname{int}(K)
\subseteq\operatorname{int}(C)$, obstaja $\alpha>0$, tako da je
$\alpha x_\eta\leq_Cx_\varepsilon$.
Z uporabo \eqref{eq:feps_eta_spodnja_meja}, homogenosti in ohranjanja
reda preslikave $f_\varepsilon$ za vsak $k\geq1$ dobimo
\begin{align*}
\alpha\bigl(r_K(f_\eta)+\delta\bigr)^k x_\eta
\leq_C
\alpha f_\varepsilon^k(x_\eta)
\leq_C
f_\varepsilon^k(x_\varepsilon)
=
r_K(f_\varepsilon)^k x_\varepsilon.
\end{align*}
Obe strani delimo z $(r_K(f_\eta)+\delta)^k$ in dobimo
\begin{align}
\left(
\frac{r_K(f_\varepsilon)}
{r_K(f_\eta)+\delta}
\right)^k
x_\varepsilon
-
\alpha x_\eta
\in C
\qquad
\text{za vsak }k\geq1.
\label{eq:primerjava_feps_feta}
\end{align}

Če bi bilo
$r_K(f_\varepsilon)<r_K(f_\eta)+\delta$, bi pri
$k\to\infty$ iz \eqref{eq:primerjava_feps_feta} in zaprtosti stožca
$C$ sledilo $-\alpha x_\eta\in C$. Ker je tudi
$\alpha x_\eta\in C$ in je $x_\eta\neq0$, je to nemogoče.
Zato je
$r_K(f_\varepsilon)\geq r_K(f_\eta)+\delta$, od koder sledi
$r_K(f_\eta)<r_K(f_\varepsilon)$ za $0<\eta<\varepsilon$.
Posledično obstaja limita
$r=\lim\limits_{\varepsilon\to0}r_K(f_\varepsilon)$.

(iv) Množici
$\{x_\varepsilon ~|~ 0<\varepsilon\leq1\}$ in
$\{f(x_\varepsilon) ~|~ 0<\varepsilon\leq1\}$ sta omejeni podmnožici
končnorazsežnega prostora $V$. Prva zato, ker je podmnožica kompaktne množice $\Sigma$ in posledično tudi kompaktna, 
druga pa zato, ker je $f$ zvezna na $\Sigma$. Ker je prva množica kompaktna, obstaja zaporedje
$(\varepsilon_k)_k$, za katero je
$\varepsilon_k\longrightarrow0$,
$x_{\varepsilon_k}\longrightarrow x$ za neki $x\in\Sigma$ in
$f(x_{\varepsilon_k})\longrightarrow y$ za neki $y\in K$.
Ker je
$f_{\varepsilon_k}(x_{\varepsilon_k})
=r_K(f_{\varepsilon_k})x_{\varepsilon_k}$ in
$\vartheta(x_{\varepsilon_k})=1$, velja
\begin{align*}
f_{\varepsilon_k}(x_{\varepsilon_k})
=
f(x_{\varepsilon_k})+\varepsilon_k u
\longrightarrow y,
\end{align*}
po drugi strani pa po točki (iii) velja
$r_K(f_{\varepsilon_k})x_{\varepsilon_k}\longrightarrow rx$.
Zato je $y=rx$.
\end{dokaz}

Če ima preslikava
$f\colon\operatorname{int}(K)\longrightarrow K$ iz prejšnjega izreka
zvezno homogeno razširitev
$F\colon K\longrightarrow K$, ki ohranja red glede na $C$, potem ima
$F$ lastni vektor v $K$ z lastno vrednostjo $r$. V tem primeru bomo
pokazali, da je $r=r_K(F)$. Zgled~\ref{zgled:nezvezna_razsiritev}
pokaže, da taka zvezna razširitev v splošnem ne obstaja. Če je $K$
poliedrski, pa lahko z argumenti, podobnimi tistim iz
podpoglavja o razširitvah preslikav, ki ohranjajo red, pokažemo, da
$f$ vedno dopušča zvezno homogeno razširitev
$F\colon K\longrightarrow K$, ki ohranja red glede na $C$.

\begin{posledica}\label{posledica:spektralni_radij_lastna_vrednost}
Naj bosta $K\subseteq C$ zaprta stožca v $V$, pri čemer ima $K$
neprazno notranjost. Če je
$f\colon K\longrightarrow K$ zvezna homogena preslikava, ki ohranja
red glede na $C$, potem obstaja $x\in K\setminus\{0\}$, za katerega
velja
$f(x)=r_K(f)x$.
\end{posledica}

\begin{dokaz}
Uporabimo oznake iz
izreka~\ref{izr:obstoj_lastnega_vektorja_perturbacija}.
Po definiciji preslikave $f_\varepsilon$ za vsak $x\in K$ velja $f(x)\leq_C f_\varepsilon(x)$, zato po
trditvi~\ref{trd:monotonost_stozcnega_spektralnega_radija} dobimo
$r_K(f)\leq r_K(f_\varepsilon)$. Ker je
$r=\lim\limits_{\varepsilon\to0}r_K(f_\varepsilon)$, sledi
$r\geq r_K(f)$.

Po točki (iv) izreka~\ref{izr:obstoj_lastnega_vektorja_perturbacija}
obstaja zaporedje $(x_{\varepsilon_k})_k$ v
$\Sigma\cap\operatorname{int}(K)$, tako da je
$x_{\varepsilon_k}\longrightarrow x$,
$\varepsilon_k\longrightarrow0$ in
$f(x_{\varepsilon_k})\longrightarrow y$, pri čemer je $y=rx$.
Ker je $f$ zvezna na $K$, velja
$f(x_{\varepsilon_k})\longrightarrow f(x)$, zato je
$f(x)=y=rx$.

Ker je $x\in\Sigma$, je $x\neq0$. Iz
$f(x)=rx$ zato po definiciji stožčnega spektralnega radija sledi
$r\leq r_K(f)$. Skupaj z že dokazano neenakostjo
$r\geq r_K(f)$ dobimo $r=r_K(f)$ in posledično
$f(x)=r_K(f)x$.
\end{dokaz}

Posledica~\ref{posledica:spektralni_radij_lastna_vrednost} je tesno
povezana s klasičnim Krein-Rutmanovim izrekom, ki
obravnava kompaktne nelinearne preslikave stožcev v Banachovih
prostorih. Neposredne posplošitve na nekompaktne preslikave stožcev
sta podala Nussbaum in Mallet-Paret, pri čemer sta uporabila
zahtevnejše metode teorije negibnih točk, med drugim indeks negibne
točke.






\subsection{Zveznost stožčnega spektralnega radija}

Naj bo $K\subseteq V$ strnjen zaprt stožec in
$f\colon K\longrightarrow K$ zvezna homogena preslikava, ki ohranja red.
Naj bo $(f_k)_k$ zaporedje zveznih homogenih preslikav na $K$, ki ohranjajo red
in za katere velja
\begin{align}
\lim_{k\to\infty}\sup_{x\in\Sigma}\|f(x)-f_k(x)\|=0,
\label{eq:enakomerna_konvergenca_fk}
\end{align}
kjer je $\Sigma=\{x\in K ~|~ \varphi(x)=1\}$ in
$\varphi\in\operatorname{int}(K^*)$.
Naravno se je vprašati, ali tedaj velja
\begin{align}
\lim_{k\to\infty}r_K(f_k)=r_K(f).
\label{eq:zveznost_stozcnega_spektralnega_radija}
\end{align}

\begin{definicija}
  Pravimo, da ima preslikava $f\colon K\longrightarrow K$ zvezen stožčni
  spektralni radij, če~\eqref{eq:zveznost_stozcnega_spektralnega_radija}
  velja za vsako zaporedje $(f_k)_k$ zveznih homogenih preslikav, ki
  ohranjajo red in zadoščajo pogoju~\eqref{eq:enakomerna_konvergenca_fk}.
\end{definicija}


V neskončnorazsežnih Banachovih prostorih obstajajo primeri linearnih
preslikav, ki so kompaktne na zaprtem stožcu, vendar niso kompaktne kot
linearne preslikave na celotnem prostoru in nimajo zveznega stožčnega
spektralnega radija. V končnorazsežnem primeru takšni primeri niso znani.
Tako je problem, ali ima vsaka zvezna homogena preslikava
$f\colon K\longrightarrow K$, ki ohranja red na strnjenem zaprtem stožcu
$K$ v končnorazsežnem vektorskem prostoru $V$, zvezen stožčni spektralni
radij, še vedno odprt.

Stožčni spektralni radij je vedno polzvezen od zgoraj, kar pokaže
naslednja lema.

\begin{lema}\label{lema:zgornja_polzveznost_stozcnega_spektralnega_radija}
Naj bo $f\colon K\longrightarrow K$ zvezna homogena preslikava, ki
ohranja red na strnjenem zaprtem stožcu $K\subseteq V$, in naj bo
$(f_k)_k$ zaporedje zveznih homogenih preslikav na $K$, ki ohranjajo
red in zadoščajo pogoju~\eqref{eq:enakomerna_konvergenca_fk}. Potem velja
\begin{align*}
\limsup_{k\to\infty}r_K(f_k)\leq r_K(f).
\end{align*}
\end{lema}

\begin{dokaz}
Za $k\geq1$ označimo $\lambda_k=r_K(f_k)$. Po posledici~\ref{posledica:spektralni_radij_lastna_vrednost} za vsak
$k$ obstaja $u_k\in K\setminus\{0\}$, za katerega velja
$f_k(u_k)=\lambda_k u_k$. Zaradi homogenosti lahko $u_k$ izberemo tako,
da je $\varphi(u_k)=1$, torej $u_k\in\Sigma$.
Res, ker je $\varphi\in\operatorname{int}(K^*)$ in $u_k\in K\setminus\{0\}$, velja $\varphi(u_k)>0$. 
Zato lahko definiramo nov vektor $\widetilde{u}_k=\frac{u_k}{\varphi(u_k)}$. Potem velja 
$\varphi(\widetilde{u}_k)=\varphi\left(\frac{u_k}{\varphi(u_k)}\right)=\frac{1}{\varphi(u_k)}\varphi(u_k)=1$.

Obstaja podzaporedje $(\lambda_{k_i})_i$, za katero velja
\begin{align*}
\lim_{i\to\infty}\lambda_{k_i}
=
\limsup_{k\to\infty}\lambda_k
=
\lambda.
\end{align*}
Ker je $\Sigma=\{x\in K ~|~ \varphi(x)=1\}$ kompaktna po lemi~\ref{lema2.46}, lahko
po potrebi izberemo še podzaporedje tako, da $(u_{k_i})_i$ konvergira
proti nekemu $v\in\Sigma$. Iz
\eqref{eq:enakomerna_konvergenca_fk}, zveznosti preslikave $f$ in
enakosti $f_{k_i}(u_{k_i})=\lambda_{k_i}u_{k_i}$ sledi
$f(v)=\lambda v$. Ker je $v\neq0$, po definiciji stožčnega spektralnega
radija dobimo $\lambda\leq r_K(f)$. Zato je
$\limsup\limits_{k\to\infty}r_K(f_k)\leq r_K(f)$.
\end{dokaz}

Če je stožec poliedrski, je stožčni spektralni radij vedno zvezen.
To sledi iz naslednjega izreka.

\begin{izrek}\label{izr:zveznost_stozcnega_spektralnega_radija_pogoj_G}
Naj bo $f\colon K\longrightarrow K$ zvezna homogena preslikava, ki
ohranja red na strnjenem zaprtem stožcu $K\subseteq V$. Če obstaja
$u\in K\setminus\{0\}$, tako da $K$ zadošča pogoju G v točki $u$ in
$f(u)=r_K(f)u$, potem ima $f$ zvezen stožčni spektralni radij.
\end{izrek}

\begin{dokaz}
Naj bo $(f_k)_{k\geq1}$ poljubno zaporedje zveznih homogenih preslikav, ki ohranjajo red in 
zadoščajo pogoju~\eqref{eq:enakomerna_konvergenca_fk}.
Za vsak $k\geq1$ postavimo $u_k=f_k(u)$. Iz~\eqref{eq:enakomerna_konvergenca_fk} in predpostavke sledi
$u_k\longrightarrow r_K(f)u$. Ker $K$ zadošča pogoju G v točki $u$,
obstajajo števila $0<\mu_k<1$, za katera je $\mu_k\longrightarrow1$ in
\begin{align*}
\mu_k r_K(f)u\leq_K u_k
\end{align*}
za vsak $k\geq1$. Ker je $u_k=f_k(u)$, dobimo
$\mu_k r_K(f)u\leq_K f_k(u)$. Po trditvi~\ref{trd:spodnja_meja_stozcnega_spektralnega_radija} zato velja
$\mu_k r_K(f)\leq r_K(f_k)$ za vsak $k\geq1$. Če pošljemo
$k\to\infty$, dobimo
\begin{align*}
r_K(f)\leq\liminf\limits_{k\to\infty}r_K(f_k).
\end{align*}
Skupaj z lemo~\ref{lema:zgornja_polzveznost_stozcnega_spektralnega_radija}
sledi
$\lim\limits_{k\to\infty}r_K(f_k)=r_K(f)$.
\end{dokaz}

Spomnimo se, da po lemi~\ref{lema:pogoj_G_poliedrski} poliedrski stožec zadošča pogoju G v vsaki
svoji točki. Iz prejšnjega izreka zato neposredno sledi naslednja
posledica.

\begin{posledica}\label{posledica:zveznost_poliedrski_stozec}
Vsaka zvezna homogena preslikava, ki ohranja red na strnjenem
poliedrskem stožcu, ima zvezen stožčni spektralni radij.
\end{posledica}

Za uporabo izreka~\ref{izr:zveznost_stozcnega_spektralnega_radija_pogoj_G}
 zadošča, da je pogoj G izpolnjen v lastnem vektorju
$u\in K$, za katerega je $f(u)=r_K(f)u$. Posebej zanimiv je primer,
ko je $u\in\operatorname{int}(K)$. V tem primeru je pogoj G v točki
$u$ izpolnjen ne glede na geometrijo stožca $K$, kar smo že pokazali.
Uporabimo izrek~\ref{izr:zveznost_stozcnega_spektralnega_radija_pogoj_G} in dobimo 
naslednjo posledico.

\begin{posledica}\label{posledica:zveznost_notranji_lastni_vektor}
Naj bo $f\colon K\longrightarrow K$ zvezna homogena preslikava, ki
ohranja red na strnjenem zaprtem stožcu $K$. Če obstaja
$u\in\operatorname{int}(K)$, za katerega je
$f(u)=r_K(f)u$, potem ima $f$ zvezen stožčni spektralni radij.
\end{posledica}

Nadaljnji rezultati glede odprtega problema so bili dobljeni z uporabo
indeksa negibne točke. Pokazano je bilo, da ima zvezna homogena
preslikava $f\colon K\longrightarrow K$, ki ohranja red, zvezen
stožčni spektralni radij, če obstaja zaporedje $(\mu_k)_k$ v
$(0,r_K(f)]$, za katero je $\mu_k\longrightarrow r_K(f)$ in noben
$\mu_k$ ni lastna vrednost preslikave $f$. Iz zgleda~\ref{zgled:kontinuum_stozcnega_spektra} vemo, da
takšno zaporedje ne obstaja vedno, saj je v tem zgledu $\sigma_K(f)=[0,1]$ in $r_K(f) = 1$, 
zato je vsako število iz intervala $(0,r_K(f)]$ lastna vrednost preslikave $f$. V tem zgledu pa ima homogena
preslikava, ki ohranja red, lastni vektor v notranjosti stožca, zato
ima po posledici~\ref{posledica:zveznost_notranji_lastni_vektor}
zvezen stožčni spektralni radij.

Podpoglavje zaključimo z naslednjo preprosto lemo, ki jo bomo
potrebovali v nadaljevanju.

\begin{lema}\label{lema:zveznost_stozcnega_spektralnega_radija_dominacija}
Naj bo $f\colon K\longrightarrow K$ zvezna homogena preslikava, ki
ohranja red na strnjenem zaprtem stožcu $K$, in naj bo $(f_k)_k$
zaporedje zveznih homogenih preslikav na $K$, ki ohranjajo red ter
zadoščajo pogoju~\eqref{eq:enakomerna_konvergenca_fk}. Če poleg tega
velja $f(x)\leq_K f_k(x)$ za vsak $x\in K$ in vsak $k\geq1$, potem
\begin{align*}
\lim\limits_{k\to\infty}r_K(f_k)=r_K(f).
\end{align*}
\end{lema}

\begin{dokaz}
Po posledici~\ref{posledica:spektralni_radij_lastna_vrednost} obstaja $u\in K\setminus\{0\}$, za katerega je
$f(u)=r_K(f)u$. Ker po predpostavki velja $f(u)\leq_K f_k(u)$, dobimo
$r_K(f)u\leq_K f_k(u)$. Po trditvi~\ref{trd:spodnja_meja_stozcnega_spektralnega_radija} zato velja
$r_K(f)\leq r_K(f_k)$ za vsak $k\geq1$, od koder sledi
$r_K(f)\leq\liminf\limits_{k\to\infty}r_K(f_k)$. Skupaj z
lemo~\ref{lema:zgornja_polzveznost_stozcnega_spektralnega_radija}
dobimo
$\lim\limits_{k\to\infty}r_K(f_k)=r_K(f)$.
\end{dokaz}







\subsection{Collatz-Wielandtova formula}

Klasični rezultat Collatza in Wielandta pravi, da za vsako nenegativno
matriko $A$ njen spektralni radij $r(A)$ zadošča formuli
\begin{align}
r(A)
=
\inf_{x\in\operatorname{int}(\mathbb{R}_+^n)}
\max_{1\leq i\leq n}
\frac{(Ax)_i}{x_i}.
\label{eq:collatz-wielandt-matrike}
\end{align}
Podobna minimaksna variacijska formula za stožčni spektralni radij
obstaja tudi za splošne zvezne homogene preslikave, ki ohranjajo red.

Naj bo $K\subseteq V$ strnjen zaprt stožec in $u\in\operatorname{int}(K)$.
Spomnimo se oznake
$\Sigma_u^*=\{\varphi\in K^* ~|~ \varphi(u)=1\}$ in naj bo
$\varepsilon_u^*$ množica ekstremnih točk množice $\Sigma_u^*$.
Po formuli~\eqref{oznaka2.6} za $x,y\in K$, kjer je
$y\in\operatorname{int}(K)$, velja
\begin{align*}
M(x/y)
=
\sup_{\varphi\in\varepsilon_u^*}
\frac{\varphi(x)}{\varphi(y)}.
\end{align*}

\begin{definicija}
  Za zvezno homogeno preslikavo $f\colon K\longrightarrow K$, ki ohranja
  red, definiramo \emph{Collatz-Wielandtovo število} s
  \begin{align}
  \operatorname{cw}(f)
  &=
  \inf_{x\in\operatorname{int}(K)} M(f(x)/x) \notag\\
  &=
  \inf_{x\in\operatorname{int}(K)}
  \max_{\varphi\in\operatorname{cl}(\varepsilon_u^*)}
  \frac{\varphi(f(x))}{\varphi(x)}.
  \label{eq:collatz-wielandt-stevilo}
  \end{align}
\end{definicija}


Če je $K=\R_+^n$ in $u=\mathds{1}$, so ekstremne točke množice
$\Sigma_u^*=\{\varphi\in\R_+^n ~|~\sum_{i=1}^n\varphi_i=1\}$
standardni bazni vektorji $e_1,\ldots,e_n$. Ker je
$e_i(x)=x_i$, sledi
$M(f(x)/x)=\max\limits_{1\leq i\leq n}\frac{f_i(x)}{x_i}$.
Zato je
$\operatorname{cw}(f)=
\inf\limits_{x\in\operatorname{int}(\R_+^n)}
\max\limits_{1\leq i\leq n}\frac{f_i(x)}{x_i}$.

\begin{izrek}\label{izr:collatz-wielandt}
Naj bo $f\colon K\longrightarrow K$ zvezna homogena preslikava, ki
ohranja red na strnjenem zaprtem stožcu $K\subseteq V$. Potem velja
\begin{align*}
r_K(f)=\operatorname{cw}(f).
\end{align*}
\end{izrek}

\begin{dokaz}
Naj bo $x\in\operatorname{int}(K)$. Ker je $K$ zaprt, po definiciji
$M(f(x)/x)$ velja
$f(x)\leq_K M(f(x)/x)x$. Po trditvi~\ref{trd:spodnja_meja_stozcnega_spektralnega_radija}
zato dobimo
\begin{align*}
r_K(f)
=
\lim\limits_{k\to\infty}\|f^k(x)\|^{1/k}
\leq
M(f(x)/x).
\end{align*}
Ker to velja za vsak $x\in\operatorname{int}(K)$, sledi
$r_K(f)\leq\inf\limits_{x\in\operatorname{int}(K)}M(f(x)/x)
=\operatorname{cw}(f)$.

Za dokaz obratne neenakosti najprej obravnavajmo primer, ko ima $f$
lastni vektor $x_0\in\operatorname{int}(K)$, za katerega je
$f(x_0)=r_K(f)x_0$. Tedaj je
$M(f(x_0)/x_0)=r_K(f)$, zato
$\operatorname{cw}(f)\leq r_K(f)$. Skupaj s prej dokazano neenakostjo
dobimo $r_K(f)=\operatorname{cw}(f)$.

Predpostavimo sedaj, da $f$ nima lastnega vektorja v
$\operatorname{int}(K)$. Izberimo $u\in\operatorname{int}(K)$ in
$\varphi\in\operatorname{int}(K^*)$. Za vsak $k\geq1$ definirajmo
zvezno homogeno preslikavo $f_k\colon K\longrightarrow K$, ki ohranja
red, s predpisom
\begin{align*}
f_k(x)=f(x)+\frac{\varphi(x)}{k}u.
\end{align*}
Za vsak $x\in K\setminus\{0\}$ je
$f_k(x)\in\operatorname{int}(K)$, saj velja $K+\operatorname{int}(K)\subseteq\operatorname{int}(K)$. 
Vzemimo $\Sigma=\{x\in K ~|~ \varphi(x)=1\}$. Ker je $f(x)-f_k(x) = -\frac{\varphi(x)}{k}u$ in $\varphi(x)=1$, velja
\begin{align*}
\lim\limits_{k\to\infty}
\sup_{x\in\Sigma}\|f(x)-f_k(x)\|= \lim\limits_{k\to\infty} \sup_{x\in\Sigma} \frac{\|u\|}{k} = 0.
\end{align*}
Poleg tega je očitno tudi 
$f(x)\leq_K f_k(x)$ za vsak $x\in K$.

Po lemi~\ref{lema:zveznost_stozcnega_spektralnega_radija_dominacija}
zato velja
$r_K(f)=\lim\limits_{k\to\infty}r_K(f_k)$.
Ker $f_k(x)\in\operatorname{int}(K)$ za vsak $x\in K\setminus\{0\}$,
ima po posledici~\ref{posledica:spektralni_radij_lastna_vrednost} vsaka preslikava $f_k$ lastni vektor v
$x_k\in K\setminus\{0\}$. Ker je $f_k(x_k)\in\operatorname{int}(K)$, mora biti 
$r_K(f_k) x_k \in \operatorname{int}(K)$, zato $x_k\in\operatorname{int}(K)$. 
Po prvem delu dokaza zato velja $\operatorname{cw}(f_k)=r_K(f_k)$.

Ker je $f(x)\leq_K f_k(x)$, velja
$M(f(x)/x)\leq M(f_k(x)/x)$ za vsak $k\geq1$ in vsak
$x\in\operatorname{int}(K)$. Zato je
$\operatorname{cw}(f)\leq\operatorname{cw}(f_k)$ za vsak $k\geq1$.
Če pošljemo $k\to\infty$, dobimo
\begin{align*}
r_K(f)
=
\lim\limits_{k\to\infty}r_K(f_k)
=
\lim\limits_{k\to\infty}\operatorname{cw}(f_k)
\geq
\operatorname{cw}(f).
\end{align*}
Skupaj z že dokazano neenakostjo
$r_K(f)\leq\operatorname{cw}(f)$ sledi
$r_K(f)=\operatorname{cw}(f)$.
\end{dokaz}



\end{document}
